% Options for packages loaded elsewhere
\PassOptionsToPackage{unicode}{hyperref}
\PassOptionsToPackage{hyphens}{url}
%
\documentclass[
]{article}
\usepackage{amsmath,amssymb}
\usepackage{iftex}
\ifPDFTeX
  \usepackage[T1]{fontenc}
  \usepackage[utf8]{inputenc}
  \usepackage{textcomp} % provide euro and other symbols
\else % if luatex or xetex
  \usepackage{unicode-math} % this also loads fontspec
  \defaultfontfeatures{Scale=MatchLowercase}
  \defaultfontfeatures[\rmfamily]{Ligatures=TeX,Scale=1}
\fi
\usepackage{lmodern}
\ifPDFTeX\else
  % xetex/luatex font selection
\fi
% Use upquote if available, for straight quotes in verbatim environments
\IfFileExists{upquote.sty}{\usepackage{upquote}}{}
\IfFileExists{microtype.sty}{% use microtype if available
  \usepackage[]{microtype}
  \UseMicrotypeSet[protrusion]{basicmath} % disable protrusion for tt fonts
}{}
\makeatletter
\@ifundefined{KOMAClassName}{% if non-KOMA class
  \IfFileExists{parskip.sty}{%
    \usepackage{parskip}
  }{% else
    \setlength{\parindent}{0pt}
    \setlength{\parskip}{6pt plus 2pt minus 1pt}}
}{% if KOMA class
  \KOMAoptions{parskip=half}}
\makeatother
\usepackage{xcolor}
\usepackage{color}
\usepackage{fancyvrb}
\newcommand{\VerbBar}{|}
\newcommand{\VERB}{\Verb[commandchars=\\\{\}]}
\DefineVerbatimEnvironment{Highlighting}{Verbatim}{commandchars=\\\{\}}
% Add ',fontsize=\small' for more characters per line
\newenvironment{Shaded}{}{}
\newcommand{\AlertTok}[1]{\textcolor[rgb]{1.00,0.00,0.00}{\textbf{#1}}}
\newcommand{\AnnotationTok}[1]{\textcolor[rgb]{0.38,0.63,0.69}{\textbf{\textit{#1}}}}
\newcommand{\AttributeTok}[1]{\textcolor[rgb]{0.49,0.56,0.16}{#1}}
\newcommand{\BaseNTok}[1]{\textcolor[rgb]{0.25,0.63,0.44}{#1}}
\newcommand{\BuiltInTok}[1]{\textcolor[rgb]{0.00,0.50,0.00}{#1}}
\newcommand{\CharTok}[1]{\textcolor[rgb]{0.25,0.44,0.63}{#1}}
\newcommand{\CommentTok}[1]{\textcolor[rgb]{0.38,0.63,0.69}{\textit{#1}}}
\newcommand{\CommentVarTok}[1]{\textcolor[rgb]{0.38,0.63,0.69}{\textbf{\textit{#1}}}}
\newcommand{\ConstantTok}[1]{\textcolor[rgb]{0.53,0.00,0.00}{#1}}
\newcommand{\ControlFlowTok}[1]{\textcolor[rgb]{0.00,0.44,0.13}{\textbf{#1}}}
\newcommand{\DataTypeTok}[1]{\textcolor[rgb]{0.56,0.13,0.00}{#1}}
\newcommand{\DecValTok}[1]{\textcolor[rgb]{0.25,0.63,0.44}{#1}}
\newcommand{\DocumentationTok}[1]{\textcolor[rgb]{0.73,0.13,0.13}{\textit{#1}}}
\newcommand{\ErrorTok}[1]{\textcolor[rgb]{1.00,0.00,0.00}{\textbf{#1}}}
\newcommand{\ExtensionTok}[1]{#1}
\newcommand{\FloatTok}[1]{\textcolor[rgb]{0.25,0.63,0.44}{#1}}
\newcommand{\FunctionTok}[1]{\textcolor[rgb]{0.02,0.16,0.49}{#1}}
\newcommand{\ImportTok}[1]{\textcolor[rgb]{0.00,0.50,0.00}{\textbf{#1}}}
\newcommand{\InformationTok}[1]{\textcolor[rgb]{0.38,0.63,0.69}{\textbf{\textit{#1}}}}
\newcommand{\KeywordTok}[1]{\textcolor[rgb]{0.00,0.44,0.13}{\textbf{#1}}}
\newcommand{\NormalTok}[1]{#1}
\newcommand{\OperatorTok}[1]{\textcolor[rgb]{0.40,0.40,0.40}{#1}}
\newcommand{\OtherTok}[1]{\textcolor[rgb]{0.00,0.44,0.13}{#1}}
\newcommand{\PreprocessorTok}[1]{\textcolor[rgb]{0.74,0.48,0.00}{#1}}
\newcommand{\RegionMarkerTok}[1]{#1}
\newcommand{\SpecialCharTok}[1]{\textcolor[rgb]{0.25,0.44,0.63}{#1}}
\newcommand{\SpecialStringTok}[1]{\textcolor[rgb]{0.73,0.40,0.53}{#1}}
\newcommand{\StringTok}[1]{\textcolor[rgb]{0.25,0.44,0.63}{#1}}
\newcommand{\VariableTok}[1]{\textcolor[rgb]{0.10,0.09,0.49}{#1}}
\newcommand{\VerbatimStringTok}[1]{\textcolor[rgb]{0.25,0.44,0.63}{#1}}
\newcommand{\WarningTok}[1]{\textcolor[rgb]{0.38,0.63,0.69}{\textbf{\textit{#1}}}}
\usepackage{longtable,booktabs,array}
\usepackage{calc} % for calculating minipage widths
% Correct order of tables after \paragraph or \subparagraph
\usepackage{etoolbox}
\makeatletter
\patchcmd\longtable{\par}{\if@noskipsec\mbox{}\fi\par}{}{}
\makeatother
% Allow footnotes in longtable head/foot
\IfFileExists{footnotehyper.sty}{\usepackage{footnotehyper}}{\usepackage{footnote}}
\makesavenoteenv{longtable}
\setlength{\emergencystretch}{3em} % prevent overfull lines
\providecommand{\tightlist}{%
  \setlength{\itemsep}{0pt}\setlength{\parskip}{0pt}}
\setcounter{secnumdepth}{-\maxdimen} % remove section numbering
\ifLuaTeX
  \usepackage{selnolig}  % disable illegal ligatures
\fi
\usepackage{bookmark}
\IfFileExists{xurl.sty}{\usepackage{xurl}}{} % add URL line breaks if available
\urlstyle{same}
\hypersetup{
  hidelinks,
  pdfcreator={LaTeX via pandoc}}

\author{}
\date{}

\begin{document}

\section{Part 1: The Structures (No
Laws)}\label{part-1-the-structures-no-laws}

\textbf{Author:} Olivier De Wolf, odewolf@gmail.com

\subsection{Table of Contents}\label{table-of-contents}

\begin{itemize}
\tightlist
\item
  \hyperref[introduction]{Introduction}
\item
  \hyperref[a-quick-primer-what-is-a-kind]{A Quick Primer: What is a
  ``Kind''?}
\item
  \hyperref[chapter-1-types-of-kind-type]{Chapter 1: Types of Kind
  \texttt{Type}}

  \begin{itemize}
  \tightlist
  \item
    \hyperref[1-0-inhabitants-uninhabited-type]{1. 0 Inhabitants
    (Uninhabited Type)}

    \begin{itemize}
    \tightlist
    \item
      \hyperref[1-custom-empty-data]{1. Custom Empty Data}
    \item
      \hyperref[2-datavoid-the-built-in-equivalent]{2.
      \texttt{Data.Void} - the built-in equivalent.}

      \begin{itemize}
      \tightlist
      \item
        \hyperref[exercise-1-implementing-the-impossible]{Exercise 1:
        Implementing the Impossible}
      \end{itemize}
    \item
      \hyperref[3-common-idioms]{3. Common Idioms}

      \begin{itemize}
      \tightlist
      \item
        \hyperref[1-type-level-guarantees]{1. Type-Level Guarantees}
      \item
        \hyperref[exercise-2-avoiding-fromright-and-partiality]{Exercise
        2: Avoiding \texttt{fromRight} and Partiality}
      \item
        \hyperref[exercise-3-traversing-without-failure-bonus]{Exercise
        3: Traversing Without Failure (Bonus)}
      \item
        \hyperref[exercise-4-safe-list-processing]{Exercise 4: Safe List
        Processing}
      \item
        \hyperref[2-type-level-phantom-types-for-type-safety]{2.
        Type-Level Phantom Types for Type Safety}
      \item
        \hyperref[exercise-5-phantom-status]{Exercise 5: Phantom Status}
      \item
        \hyperref[exercise-6-a-tree-without-leaves]{Exercise 6: A Tree
        Without Leaves}
      \end{itemize}
    \end{itemize}
  \item
    \hyperref[2-1-inhabitant-unit-type]{2. 1 Inhabitant (Unit Type)}

    \begin{itemize}
    \tightlist
    \item
      \hyperref[1-custom-unit-types]{1. Custom Unit Types}
    \item
      \hyperref[2-the-standard-unit-]{2. The Standard Unit \texttt{()}}
    \item
      \hyperref[3-the-0-tuple-intuition]{3. The ``0-Tuple'' Intuition}
    \item
      \hyperref[4-other-library-unit-types]{4. Other Library Unit Types}
    \item
      \hyperref[5-exercises-the-power-of-one]{5. Exercises: The Power of
      One}

      \begin{itemize}
      \tightlist
      \item
        \hyperref[exercise-7-a-safe-head]{Exercise 7: A Safe
        \texttt{head}}
      \item
        \hyperref[exercise-8-avoiding-fromjust-with-either]{Exercise 8:
        Avoiding \texttt{fromJust} with \texttt{Either}}
      \end{itemize}
    \end{itemize}
  \item
    \hyperref[3-2-inhabitants-boolean-type]{3. 2 Inhabitants (Boolean
    Type)}

    \begin{itemize}
    \tightlist
    \item
      \hyperref[1-the-standard-bool]{1. The Standard \texttt{Bool}}
    \item
      \hyperref[3-custom-enumerations]{3. Custom Enumerations}
    \end{itemize}
  \item
    \hyperref[4-uncountably-infinite-inhabitants]{4. Uncountably
    Infinite Inhabitants}

    \begin{itemize}
    \tightlist
    \item
      \hyperref[1-infinite-streams]{1. Infinite Streams}
    \item
      \hyperref[2-infinite-functions]{2. Infinite Functions}
    \end{itemize}
  \end{itemize}
\item
  \hyperref[annex-]{Annex}

  \begin{itemize}
  \tightlist
  \item
    \hyperref[the-secret-inhabitant-bottom-__]{The Secret Inhabitant:
    Bottom (\texttt{\_\textbar{}\_})}
  \item
    \hyperref[bibliography]{Bibliography}
  \end{itemize}
\end{itemize}

\subsection{Introduction}\label{introduction}

In type theory and functional programming, a profound dichotomy exists
that mirrors mathematics: the distinction between \textbf{Structures}
(Data Types without inherent laws) and \textbf{Algebras} (Typeclasses
with mathematical laws).

\begin{enumerate}
\def\labelenumi{\arabic{enumi}.}
\tightlist
\item
  \textbf{The Structures (Nouns):} A concrete data type (like
  \texttt{Bool}, \texttt{Maybe}, or \texttt{Void}) simply defines a
  ``shape in memory'' by explicitly declaring how many distinct values
  (inhabitants) it can hold. It has no strict mathematical laws
  governing how it must behave; its only rules are structural.
\item
  \textbf{The Algebras (Verbs/Adjectives):} A typeclass (like
  \texttt{Eq}, \texttt{Semigroup}, or \texttt{Functor}) defines an
  interface of behavior mapping across these structures. Because these
  define behavior, they explicitly come with \textbf{Mathematical Laws}
  to ensure that behavior is predictable and compositionally sound.
\end{enumerate}

This document focuses firmly on the first half of that dichotomy:
\textbf{The Structures.} We will explore how classifying types purely by
the number of distinct values they can hold at runtime provides a
phenomenally strong foundation for building robust abstractions.

While minimal types (like those with 0 or 1 inhabitant) are omnipresent
in pure functional languages like Haskell, they can often feel
counter-intuitive or overly abstract to newcomers. Why would we want a
type that holds zero values? What is the point of a type with exactly
one? This document aims to demystify these concepts. To aid in your
learning journey, several hands-on exercises are suggested throughout
this guide.

\textbf{Intended Audience:} From a pedagogical perspective, this guide
is tailored for intermediate Haskell learners and practical software
engineers. If you've ever struggled to understand \emph{why} concepts
like \texttt{Void} or \texttt{Proxy} exist in the standard
library---rather than just \emph{how} to compile them---this resource
provides the foundation. After mastering the \emph{Structures} outlined
here, you will be perfectly prepared to study the \emph{Algebras}
(Typeclasses and their laws) that bring them to life.

\subsection{A Quick Primer: What is a
``Kind''?}\label{a-quick-primer-what-is-a-kind}

Before we dive into counting inhabitants, we need to clarify what we
mean by ``Kind''. Just as \textbf{types} classify \textbf{values} (e.g.,
\texttt{True} is a value of type \texttt{Bool}), \textbf{kinds} classify
\textbf{types}. - \textbf{\texttt{Type}}: This is the kind of a concrete
type that can actually hold values at runtime. For example,
\texttt{Int}, \texttt{Bool}, \texttt{String}, and \texttt{Maybe\ Int}
are all of kind \texttt{Type}. \emph{(Note: Older Haskell code uses
\texttt{*} for this, but modern Haskell uses \texttt{Type} imported from
\texttt{Data.Kind} to avoid conflict with the multiplication operator).}
- \textbf{Also known as}: Concrete Types, Fully Applied Types, Nullary
Type Constructors. - \textbf{Typeclasses}: Typeclasses that operate on
fully realized data (like \texttt{Eq}, \texttt{Ord}, \texttt{Semigroup},
and \texttt{Monoid}) expect variables of kind \texttt{Type}. You can
append two lists of strings, but you cannot append an abstract
\texttt{Maybe}. - \textbf{\texttt{Type\ -\textgreater{}\ Type}}: This is
the kind of a \emph{type constructor} that takes one concrete type and
returns a new concrete type. For example, \texttt{{[}{]}} (List) and
\texttt{Maybe} are of kind \texttt{Type\ -\textgreater{}\ Type} because
they need a type argument (like \texttt{Int}) to become a concrete type
(\texttt{{[}Int{]}} or \texttt{Maybe\ Int}) of kind \texttt{Type}. -
\textbf{Also known as}: Higher-Kinded Types (HKTs), Unary Type
Constructors. - \textbf{Typeclasses}: Typeclasses that operate on
``shapes'' or ``containers'' (like \texttt{Functor},
\texttt{Applicative}, \texttt{Monad}, \texttt{Foldable}, and
\texttt{Traversable}) expect variables of kind
\texttt{Type\ -\textgreater{}\ Type}.

With that in mind, the first chapter will focus on the minimal types of
the simplest kind \texttt{Type} while the second chapter will be devoted
to the minimal types of the second simplest kind
\texttt{Type\ -\textgreater{}\ Type}.

\subsection{\texorpdfstring{Chapter 1: Types of Kind
\texttt{Type}}{Chapter 1: Types of Kind Type}}\label{chapter-1-types-of-kind-type}

These are standard, concrete types that take no parameters. A type of
kind \texttt{Type} can contain any finite number of inhabitants (and
here we will focus specifically on types containing exactly 0, 1, or 2
inhabitants) or an infinite number of inhabitants (we will also later
have a section devoted to minimal countably infinite types).

Furthermore, it is a fundamental property of type theory that any finite
types with exactly \(n\) inhabitants are strictly \textbf{isomorphic} to
each other. Because they have the exact same number of possible distinct
runtime values, they can be perfectly mapped back and forth without
losing any information. For example, the \texttt{Bool} type (which has 2
inhabitants: \texttt{True} and \texttt{False}) is perfectly isomorphic
to the type \texttt{Either\ ()\ ()} (which also has 2 inhabitants:
\texttt{Left\ ()} and \texttt{Right\ ()}). \emph{(Note: As we talked
about in the Annex, as long as we treat Haskell as a total language and
ignore \texttt{\_\textbar{}\_}, this isomorphism holds perfectly true in
the Haskell type-checker!)}

\subsubsection{\texorpdfstring{Section 1.1: \texttt{Void} (0 Inhabitants
/ Initial
Object)}{Section 1.1: Void (0 Inhabitants / Initial Object)}}\label{section-1.1-void-0-inhabitants-initial-object}

\textbf{Categorical Analog:} \emph{The Initial Object} (denoted as
\(0\)). In Category Theory, an initial object has exactly one unique
morphism \emph{to} every other object in the category, but none coming
in from populated objects. In Haskell, this corresponds to the
\texttt{absurd} function (which can produce any type from
\texttt{Void}).

An uninhabited type (or empty type) is a type that has absolutely no
data constructors. Because there are no constructors, it is impossible
to create a value of this type at runtime.

\paragraph{1. Custom Empty Data}\label{custom-empty-data}

You can construct your own empty type with no constructors simply by
providing no right-hand side.

\begin{Shaded}
\begin{Highlighting}[]
\KeywordTok{data} \DataTypeTok{Never}

\CommentTok{{-}{-} This is impossible! There is no constructor to provide.}
\CommentTok{{-}{-} myNever :: Never}
\CommentTok{{-}{-} myNever = ???}
\end{Highlighting}
\end{Shaded}

\emph{(Note for beginners: The name \texttt{Never} is completely
arbitrary! You could name this \texttt{Void}, \texttt{Empty},
\texttt{MyImpossibleType}, or anything else; the compiler only cares
that it lacks data constructors.)}

\paragraph{\texorpdfstring{2. \texttt{Data.Void} - the built-in
equivalent.}{2. Data.Void - the built-in equivalent.}}\label{data.void---the-built-in-equivalent.}

Haskell provides a standard empty type called \texttt{Void} in the
\texttt{Data.Void} module (see the
\href{https://hackage.haskell.org/package/base/docs/Data-Void.html}{Hackage
\texttt{Data.Void} documentation}).

\begin{Shaded}
\begin{Highlighting}[]
\KeywordTok{import} \DataTypeTok{Data.Void}\NormalTok{ (}\DataTypeTok{Void}\NormalTok{)}
\CommentTok{{-}{-} Void has 0 constructors.}

\CommentTok{{-}{-} Like Never, it cannot be instantiated.}
\CommentTok{{-}{-} myVoid :: Void}
\CommentTok{{-}{-} myVoid = ???}
\end{Highlighting}
\end{Shaded}

Because \texttt{Data.Void} has exactly 0 inhabitants just like our
custom \texttt{Never} type, they are strictly \textbf{isomorphic}.
However, \texttt{Data.Void} is overwhelmingly preferred in practice
because it is the standard empty type and comes bundled with powerful
tooling for working with impossible values: -
\texttt{absurd\ ::\ Void\ -\textgreater{}\ a}: Since it is impossible to
ever actually have a value of type \texttt{Void}, a function taking it
can return \emph{any} type \texttt{a}. This is used to exhaustively
handle impossible code branches (see the Annex). -
\texttt{vacuous\ ::\ Functor\ f\ =\textgreater{}\ f\ Void\ -\textgreater{}\ f\ a}:
If you map over a functor that ``contains'' \texttt{Void} (which means
it's structurally empty, like \texttt{Right} on an
\texttt{Either\ Void\ a}), you can safely cast it to contain any type
\texttt{a} instead.

\begin{quote}
{[}!NOTE{]} \#\#\#\# Functions Returning an Uninhabited Type

What happens if we try to write a function that \emph{returns} our empty
type \texttt{Never}? For a standard function like
\texttt{createNever\ ::\ a\ -\textgreater{}\ Never}, it is
\textbf{impossible} to provide a valid total implementation (for a
discussion on computations that crash or loop indefinitely, see the
section on ``Bottom'' in the Annex). Because we cannot instantiate a
\texttt{Never}, we can never return one.

But \textbf{importantly}, it \emph{is} possible to implement functions
returning \texttt{Never} \textbf{if one of the inputs is also impossible
to provide!} For example, consider the signature
\texttt{f\ ::\ a\ -\textgreater{}\ Never\ -\textgreater{}\ b\ -\textgreater{}\ Never}.
Because the caller can never properly supply the second argument
(\texttt{Never}), it is mathematically sound (based on the
\textbf{Principle of Explosion}, or \emph{ex falso quodlibet}, which
states that from a false premise any conclusion can be drawn) to return
a \texttt{Never}. This principle is a direct application of the
Curry-Howard correspondence {[}1{]}. We do this by proving to the
compiler that the code branch is unreachable.

Here is how you can implement such a function using an empty case
expression:

\begin{Shaded}
\begin{Highlighting}[]
\CommentTok{{-}{-} Requires \{{-}\# LANGUAGE EmptyCase \#{-}\}}
\OtherTok{f ::}\NormalTok{ a }\OtherTok{{-}\textgreater{}} \DataTypeTok{Never} \OtherTok{{-}\textgreater{}}\NormalTok{ b }\OtherTok{{-}\textgreater{}} \DataTypeTok{Never}
\NormalTok{f \_ impossibleValue \_ }\OtherTok{=} \KeywordTok{case}\NormalTok{ impossibleValue }\KeywordTok{of}\NormalTok{ \{\}}
\end{Highlighting}
\end{Shaded}

There are also a few other ways to implement this without explicitly
writing \texttt{case\ ...\ of\ \{\}}: 1. \textbf{Using
\texttt{LambdaCase}}: With \texttt{\{-\#\ LANGUAGE\ LambdaCase\ \#-\}},
you can drop the variable name and write
\texttt{\textbackslash{}case\ \{\}}. 2. \textbf{Using standard library
tools (\texttt{absurd})}: If you use the standard library's
\texttt{Void} type instead of a custom \texttt{Never}, the library
provides a function called
\texttt{absurd\ ::\ Void\ -\textgreater{}\ a}. Because \texttt{absurd}
can return \emph{any} type \texttt{a}, you can simply use it to return
\texttt{Void} right back! \texttt{f\ \_\ imposs\ \_\ =\ absurd\ imposs}.

\emph{(For the deep technical details on how the compiler handles
matching on uninhabited types, refer to the
\href{https://ghc.gitlab.haskell.org/ghc/doc/users_guide/exts/empty_case.html}{GHC
User Guide on EmptyCase}).}

However, we rarely need to write our own custom empty types because
Haskell's standard library provides a built-in one! \#\#\#\#\# Exercise
1: Implementing the Impossible It is actually a great exercise to
understand how to implement the standard empty type tooling yourself!
\end{quote}

\begin{enumerate}
\def\labelenumi{\arabic{enumi}.}
\tightlist
\item
  Given your own custom empty type \texttt{data\ Never}, how would you
  implement your own \texttt{absurd\ ::\ Never\ -\textgreater{}\ a}?
\item
  Now, using your defined \texttt{absurd} function, how would you
  implement
  \texttt{vacuous\ ::\ Functor\ f\ =\textgreater{}\ f\ Never\ -\textgreater{}\ f\ a}?
\end{enumerate}

View Solutions

\textbf{1.} By enabling the \texttt{EmptyCase} language extension, we
can pattern match on the impossible value. Since the compiler sees there
are 0 constructors for \texttt{Never}, we don't even have to provide a
right-hand side for the case expression!

\begin{Shaded}
\begin{Highlighting}[]
\OtherTok{\{{-}\# LANGUAGE EmptyCase \#{-}\}}

\KeywordTok{data} \DataTypeTok{Never}

\OtherTok{absurd ::} \DataTypeTok{Never} \OtherTok{{-}\textgreater{}}\NormalTok{ a}
\NormalTok{absurd v }\OtherTok{=} \KeywordTok{case}\NormalTok{ v }\KeywordTok{of}\NormalTok{ \{\}}
\end{Highlighting}
\end{Shaded}

\textbf{2.} Because \texttt{absurd} can turn a \texttt{Never} into any
type \texttt{a}, all we need to do is map it over the functor!

\begin{Shaded}
\begin{Highlighting}[]
\OtherTok{vacuous ::} \DataTypeTok{Functor}\NormalTok{ f }\OtherTok{=\textgreater{}}\NormalTok{ f }\DataTypeTok{Never} \OtherTok{{-}\textgreater{}}\NormalTok{ f a}
\NormalTok{vacuous }\OtherTok{=} \FunctionTok{fmap}\NormalTok{ absurd}
\end{Highlighting}
\end{Shaded}

\emph{(Note on safety: Why does \texttt{vacuous} not crash if calling
\texttt{absurd} crashes? Because \texttt{absurd} only crashes if you
actually give it a \texttt{Never} value! If an \texttt{f\ Never} exists
at runtime, the structure \texttt{f} must logically be ``empty'' of
values---such as \texttt{Nothing}, an empty list \texttt{{[}{]}}, or a
\texttt{Right}. Consequently, \texttt{fmap} traverses the container but
never actually finds a \texttt{Never} value to apply \texttt{absurd} to,
meaning the code safely evaluates without crashing!)}

\paragraph{3. Common Idioms}\label{common-idioms}

Uninhabited (Empty) types might seem useless at first glance since you
can never construct them. However, they are incredibly powerful tools
for the compiler. Here are some very useful common idioms using
uninhabited types:

\subparagraph{1. Type-Level Guarantees}\label{type-level-guarantees}

One of the most frequent patterns when dealing with impossible states is
safely extracting a value from a sum type where one branch can never
happen. This section will demonstrate how to elegantly establish and
resolve these type-level guarantees by leveraging the
\texttt{Either\ Void\ a} structure alongside the
\texttt{either\ absurd\ id} idiom. \emph{(Note: you can seamlessly apply
the exact same logic to the right side using \texttt{Either\ a\ Void}
and \texttt{either\ id\ absurd}!)}

By encoding the impossibility of failure directly into the type
signature (e.g.~\texttt{Either\ Void\ a}), you mathematically prove a
computation is \emph{guaranteed} to succeed! This approach is a far
safer alternative to relying on notorious partial functions---like
\texttt{head}, \texttt{fromJust}, \texttt{read}, or the list index
operator \texttt{!!}---that will crash your entire program at runtime if
handed unexpected input.

\begin{quote}
{[}!NOTE{]} If we have an \texttt{Either\ Void\ a}, we know the
\texttt{Left} branch is impossible. The compiler is actually smart
enough to know that since \texttt{Void} cannot exist, the \emph{entire
\texttt{Left} constructor} also cannot exist because it demands a
\texttt{Void} at runtime! If you enable the \texttt{EmptyCase} extension
(or use newer compiler versions with it built-in), you can pattern match
exclusively on the \texttt{Right} side. You can completely omit the
\texttt{Left} case, and the compiler will cleanly accept it
\emph{without} issuing any ``incomplete pattern match'' warnings!
\end{quote}

\begin{Shaded}
\begin{Highlighting}[]
\OtherTok{\{{-}\# LANGUAGE EmptyCase \#{-}\}}
\OtherTok{collapseLeft ::} \DataTypeTok{Either} \DataTypeTok{Void}\NormalTok{ a }\OtherTok{{-}\textgreater{}}\NormalTok{ a}
\NormalTok{collapseLeft (}\DataTypeTok{Right}\NormalTok{ x) }\OtherTok{=}\NormalTok{ x}
\CommentTok{{-}{-} The compiler mathematically proves Left is impossible and doesn\textquotesingle{}t require us to write it!}
\end{Highlighting}
\end{Shaded}

Instead of writing custom pattern-matching functions like this, it is
highly idiomatic in general to just use the standard library's
\texttt{either} function combined with \texttt{id} and \texttt{absurd}
to collapse the ``impossible'' branch.

In fact, a dedicated \texttt{collapseLeft} function is almost never
explicitly defined in real Haskell codebases because developers simply
use \texttt{either\ absurd\ id} directly inline! This provides a
one-line mathematical proof to the compiler, allowing safe extraction
without ever using dangerous partial functions like \texttt{fromRight}:

\begin{Shaded}
\begin{Highlighting}[]
\CommentTok{{-}{-} In practice, developers just use \textasciigrave{}either absurd id\textasciigrave{} inline!}
\OtherTok{collapseLeft ::} \DataTypeTok{Either} \DataTypeTok{Void}\NormalTok{ a }\OtherTok{{-}\textgreater{}}\NormalTok{ a}
\NormalTok{collapseLeft }\OtherTok{=} \FunctionTok{either}\NormalTok{ absurd }\FunctionTok{id}
\end{Highlighting}
\end{Shaded}

\subparagraph{\texorpdfstring{Exercise 2: Avoiding \texttt{fromRight}
and
Partiality}{Exercise 2: Avoiding fromRight and Partiality}}\label{exercise-2-avoiding-fromright-and-partiality}

Imagine you inherit a codebase where a previous developer used
\texttt{Data.Either.fromRight} to extract a value from a guaranteed
computation by passing a runtime \texttt{error} as the default value:

\begin{Shaded}
\begin{Highlighting}[]
\KeywordTok{import} \DataTypeTok{Data.Either}\NormalTok{ (fromRight)}

\OtherTok{unsafeExtract ::} \DataTypeTok{Either} \DataTypeTok{Void}\NormalTok{ a }\OtherTok{{-}\textgreater{}}\NormalTok{ a}
\NormalTok{unsafeExtract comp }\OtherTok{=}\NormalTok{ fromRight (}\FunctionTok{error} \StringTok{"This should be impossible!"}\NormalTok{) comp}
\end{Highlighting}
\end{Shaded}

Why is this approach dangerous even if the \texttt{Left} branch is
\texttt{Void}, and how can you write a new \texttt{safeExtract} function
in a single line using standard library combinators, completely removing
the \texttt{error}?

View Solution

By relying on \texttt{error} as a default value, the developer bypassed
the type-checker and introduced a potential program crash (a bottom
value, \texttt{\_\textbar{}\_}) if the code was ever refactored to
something other than \texttt{Void}!

You can replace the entirely unsafe helper with the idiomatic
mathematical proof we learned! By passing \texttt{absurd} to the
\texttt{Left} handler, we safely prove to the compiler that the left
branch is unreachable, which would instantly fail to compile if anyone
ever changed \texttt{Void} to a real type later on.

\begin{Shaded}
\begin{Highlighting}[]
\OtherTok{safeExtract ::} \DataTypeTok{Either} \DataTypeTok{Void}\NormalTok{ a }\OtherTok{{-}\textgreater{}}\NormalTok{ a}
\NormalTok{safeExtract }\OtherTok{=} \FunctionTok{either}\NormalTok{ absurd }\FunctionTok{id}
\end{Highlighting}
\end{Shaded}

\emph{(Tip: If the type is ever flipped so the impossible \texttt{Void}
is on the right side, like \texttt{Either\ a\ Void}, you can simply flip
the handlers to extract your valid value: \texttt{either\ id\ absurd}!)}

\subparagraph{Exercise 3: Traversing Without Failure
(Bonus)}\label{exercise-3-traversing-without-failure-bonus}

The \texttt{traverse} function is commonly used to map a fallible
function over a sequence of elements:
\texttt{traverse\ ::\ Applicative\ f\ =\textgreater{}\ (a\ -\textgreater{}\ f\ b)\ -\textgreater{}\ {[}a{]}\ -\textgreater{}\ f\ {[}b{]}}
When specialized to \texttt{Either\ e}, its signature effectively
becomes:
\texttt{traverse\ ::\ (a\ -\textgreater{}\ Either\ e\ b)\ -\textgreater{}\ {[}a{]}\ -\textgreater{}\ Either\ e\ {[}b{]}}
Imagine you have a list of valid inputs \texttt{{[}a{]}}, and a specific
function \texttt{process\ ::\ a\ -\textgreater{}\ Either\ Void\ b} that
is \emph{mathematically guaranteed} to succeed (perhaps reusing a parser
or computation that theoretically \emph{could} fail on some data, but
not on this specific data). How can you use \texttt{either\ absurd\ id}
to write a function
\texttt{processAll\ ::\ {[}a{]}\ -\textgreater{}\ {[}b{]}} that
completely sheds the \texttt{Either} wrapper from the resulting list?

View Solution

Because \texttt{process} returns \texttt{Either\ Void\ b}, mapping it
via \texttt{traverse\ process} will return an
\texttt{Either\ Void\ {[}b{]}}. Since the type system proves the
sequence of computations cannot possibly fail, we can safely extract our
final list of results using the exact same idiom!

\begin{Shaded}
\begin{Highlighting}[]
\OtherTok{processAll ::}\NormalTok{ [a] }\OtherTok{{-}\textgreater{}}\NormalTok{ [b]}
\NormalTok{processAll xs }\OtherTok{=} \FunctionTok{either}\NormalTok{ absurd }\FunctionTok{id}\NormalTok{ (}\FunctionTok{traverse}\NormalTok{ process xs)}
\end{Highlighting}
\end{Shaded}

\subparagraph{Exercise 4: Safe List
Processing}\label{exercise-4-safe-list-processing}

Suppose a trusted API parser guarantees that a deeply-nested JSON
payload sequence will never return an error branch, returning
\texttt{{[}Either\ Void\ User{]}}. Using the idiom we just learned, how
can you elegantly extract a clean \texttt{{[}User{]}} list?

View Solution

Since \texttt{either\ absurd\ id} is a basic, total function
\texttt{Either\ Void\ a\ -\textgreater{}\ a}, we just pass it to
\texttt{map} (or \texttt{fmap}) to unwrap every single element! Because
it leverages absolute mathematical proofs, it's significantly safer than
\texttt{Data.Either.rights}, which filters elements without statically
proving none were lost.

\begin{Shaded}
\begin{Highlighting}[]
\OtherTok{extractUsers ::}\NormalTok{ [}\DataTypeTok{Either} \DataTypeTok{Void} \DataTypeTok{User}\NormalTok{] }\OtherTok{{-}\textgreater{}}\NormalTok{ [}\DataTypeTok{User}\NormalTok{]}
\NormalTok{extractUsers }\OtherTok{=} \FunctionTok{map}\NormalTok{ (}\FunctionTok{either}\NormalTok{ absurd }\FunctionTok{id}\NormalTok{)}
\end{Highlighting}
\end{Shaded}

\subparagraph{2. Type-Level Phantom Types for Type
Safety}\label{type-level-phantom-types-for-type-safety}

In many mainstream languages like Java, C++, or Python, developers often
rely on \texttt{const} modifiers, \texttt{final} keywords, or empty
``marker interfaces'' to tag data and enforce invariants at
compile-time. Haskell achieves a much more powerful and flexible version
of this exact same concept using \textbf{Phantom Types}.

Empty type declarations are commonly used as tags for these phantom
types. A phantom type parameter is one that appears on the left side of
a type definition but not on the right.

\begin{Shaded}
\begin{Highlighting}[]
\KeywordTok{data} \DataTypeTok{USD} \CommentTok{{-}{-} 0 inhabitants}
\KeywordTok{data} \DataTypeTok{EUR} \CommentTok{{-}{-} 0 inhabitants}

\KeywordTok{newtype} \DataTypeTok{Money}\NormalTok{ currency }\OtherTok{=} \DataTypeTok{Money} \DataTypeTok{Double}

\OtherTok{fiveDollars ::} \DataTypeTok{Money} \DataTypeTok{USD}
\NormalTok{fiveDollars }\OtherTok{=} \DataTypeTok{Money} \FloatTok{5.0}

\OtherTok{fiveEuros ::} \DataTypeTok{Money} \DataTypeTok{EUR}
\NormalTok{fiveEuros }\OtherTok{=} \DataTypeTok{Money} \FloatTok{5.0}

\CommentTok{{-}{-} This will result in a compile{-}time error:}
\CommentTok{{-}{-} error: Couldn\textquotesingle{}t match type ‘EUR’ with ‘USD’}
\CommentTok{{-}{-} illegalSum = fiveDollars \textasciigrave{}addMoney\textasciigrave{} fiveEuros}
\end{Highlighting}
\end{Shaded}

Because \texttt{USD} and \texttt{EUR} have no constructors, we never
intended to instantiate them. We only use them as ``labels'' at
compile-time to prevent mixing up currencies. The compiler will now
throw an error if we accidentally try to add dollars and euros together,
completely eliminating a whole class of bugs at zero runtime cost!

\subparagraph{Exercise 5: Phantom
Status}\label{exercise-5-phantom-status}

Imagine a \texttt{Document\ status} type where \texttt{status} can be
\texttt{Draft} or \texttt{Published} (both uninhabited types). Write a
function signature
\texttt{publish\ ::\ Document\ Draft\ -\textgreater{}\ Document\ Published}
and explain why you cannot accidentally pass a \texttt{Published}
document to \texttt{publish}.

View Solution

Because \texttt{publish} explicitly requires a \texttt{Document\ Draft},
providing a \texttt{Document\ Published} will result in a compile-time
type mismatch error. This guarantees at compile time that we only
publish drafts, and prevents re-publishing already published documents!

\subparagraph{Exercise 6: A Tree Without
Leaves}\label{exercise-6-a-tree-without-leaves}

Consider a simple parameterised binary tree:

\begin{Shaded}
\begin{Highlighting}[]
\KeywordTok{data} \DataTypeTok{Tree}\NormalTok{ a }\OtherTok{=} \DataTypeTok{Leaf}\NormalTok{ a }\OperatorTok{|} \DataTypeTok{Node}\NormalTok{ (}\DataTypeTok{Tree}\NormalTok{ a) (}\DataTypeTok{Tree}\NormalTok{ a)}
\end{Highlighting}
\end{Shaded}

If we use the \texttt{Void} type as the type parameter \texttt{a} to
declare the type \texttt{Tree\ Void}, what happens when we try to
construct a value of this type? What does this imply about the shape of
the tree?

View Solution

Because it is impossible to instantiate a \texttt{Void}, we can never
use the \texttt{Leaf} constructor (which requires a \texttt{Void}
value). This means any valid value of type \texttt{Tree\ Void} can
\emph{only} be constructed using \texttt{Node}s. Therefore, a
\texttt{Tree\ Void} must be an \textbf{infinitely deep tree} containing
no leaves!

\emph{(Bonus thought: Any attempt to manually construct a finite
\texttt{Tree\ Void} in Haskell would require ``cheating'' the type
system by explicitly placing a bottom (\texttt{\_\textbar{}\_}) value,
such as \texttt{undefined} or \texttt{error}, in the \texttt{Leaf}
position!)}

\subsubsection{\texorpdfstring{Section 1.2: \texttt{()} (1 Inhabitant /
Terminal
Object)}{Section 1.2: () (1 Inhabitant / Terminal Object)}}\label{section-1.2-1-inhabitant-terminal-object}

\textbf{Categorical Analog:} \emph{The Terminal Object} (denoted as
\(1\)). In Category Theory, a terminal object has exactly one unique
morphism coming in \emph{from} every other object. In Haskell, this
corresponds to the fact that you can map any arbitrary Type into the
Unit type simply by throwing away the input value
(\texttt{\textbackslash{}x\ -\textgreater{}\ ()}).

A 1-inhabitant type has exactly one possible value. Inspecting the value
tells you nothing new---it simply conveys ``this computation finished''
or acts as a structural placeholder.

\paragraph{1. Custom Unit Types}\label{custom-unit-types}

You can easily define your own single-inhabitant types if you want them
to carry a specific semantic meaning in your application rather than
using generic units.

\begin{Shaded}
\begin{Highlighting}[]
\KeywordTok{data} \DataTypeTok{Acknowledged} \OtherTok{=} \DataTypeTok{Acknowledged}

\CommentTok{{-}{-} We can instantiate it using its single constructor:}
\OtherTok{myAck ::} \DataTypeTok{Acknowledged}
\NormalTok{myAck }\OtherTok{=} \DataTypeTok{Acknowledged}
\end{Highlighting}
\end{Shaded}

\paragraph{\texorpdfstring{2. The Standard Unit
\texttt{()}}{2. The Standard Unit ()}}\label{the-standard-unit}

The most common 1-inhabitant type in Haskell is \texttt{()} (pronounced
``unit''). Its only value is also written as \texttt{()}.

\begin{Shaded}
\begin{Highlighting}[]
\CommentTok{{-}{-} The type is (), the value is ()}
\OtherTok{myUnit ::}\NormalTok{ ()}
\NormalTok{myUnit }\OtherTok{=}\NormalTok{ ()}
\end{Highlighting}
\end{Shaded}

\paragraph{3. The ``0-Tuple'' Intuition}\label{the-0-tuple-intuition}

It is highly insightful to understand \emph{why} the unit type is
written with empty parentheses \texttt{()}. In type algebra, tuples
represent multiplication (Product types). For example: - A pair
\texttt{(a,\ b)} has an inhabitant cardinality of \(n \times m\).\\
- A triple \texttt{(a,\ b,\ c)} has an inhabitant cardinality of
\(n \times m \times p\).

If we extend this pattern downwards to an ``empty tuple'' \texttt{()},
it formally represents the \textbf{empty product}. In mathematics, the
empty product evaluates exactly to \textbf{1}. This beautifully explains
why \texttt{()} has exactly 1 inhabitant and fits perfectly into our
mathematical theme!

Another brilliant consequence of this relates to function arity.
Mathematically, an \(n\)-ary function requires \(n\) elements to provide
one result. - A \textbf{2-ary} function takes 2 inputs to provide 1
output, written as \texttt{(a,\ a)\ -\textgreater{}\ a} or curried as
\texttt{a\ -\textgreater{}\ a\ -\textgreater{}\ a}. (Examples:
\texttt{(+)} or \texttt{(*)}). - A \textbf{1-ary} function takes 1 input
to provide 1 output, written as \texttt{(a)\ -\textgreater{}\ a} or
simply \texttt{a\ -\textgreater{}\ a}. (Examples: \texttt{not} or
\texttt{negate}).

So what is a \textbf{0-ary} function? It must provide an output out of
\emph{no} input. It is equivalent to choosing an element from a set with
no prior information. Mathematically, this is modeled via a function
from a \textbf{singleton set} (a set with exactly one element, like
\texttt{\{*\}}) to a target type \texttt{a}. The function
\texttt{f\ :\ \{*\}\ -\textgreater{}\ a} simply picks exactly one value
of type \texttt{a} when you evaluate \texttt{f(*)}.

Because all singleton sets are strictly \textbf{isomorphic} to one
another (i.e., \texttt{()} \(\cong\) \texttt{\{myElement\}} \(\cong\)
\texttt{\{*\}}), the actual value of their single element is completely
irrelevant. In Haskell, we use the unit type \texttt{()} as the
standard, generic representation of this singleton concept. Therefore, a
0-ary computation has the exact signature
\texttt{()\ -\textgreater{}\ a}. This beautifully explains why we use
\texttt{()} to represent an action or computation that requires no
meaningful input to provide a result!

\paragraph{4. Other Library Unit Types}\label{other-library-unit-types}

While \texttt{()} is standard, Haskell libraries often use specialized
1-inhabitant types for specific contexts: -
\textbf{\texttt{Data.Functor.Identity}}: The \texttt{Identity\ ()} type
has only one inhabitant (\texttt{Identity\ ()}). It is used as a base
functor that doesn't add any effects. - \textbf{Type Equality
\texttt{(:\textasciitilde{}:)}}: From \texttt{Data.Type.Equality}, a
value of type \texttt{a\ :\textasciitilde{}:\ a} has exactly one
inhabitant, \texttt{Refl}, representing a proof that two types are
equal.

Because \texttt{Acknowledged}, \texttt{()}, \texttt{Identity\ ()}, and
\texttt{a\ :\textasciitilde{}:\ a} all have an identical cardinality of
1 (a single constructor), they are all formally \textbf{isomorphic} to
one another.

\paragraph{5. Exercises: The Power of
One}\label{exercises-the-power-of-one}

\subparagraph{\texorpdfstring{Exercise 7: A Safe
\texttt{head}}{Exercise 7: A Safe head}}\label{exercise-7-a-safe-head}

The standard library's \texttt{head\ ::\ {[}a{]}\ -\textgreater{}\ a}
function is notorious for crashing if given an empty list because it
lacks a value to return. How could you write a total, non-crashing
\texttt{safeHead} function using \texttt{Either}? What minimal type is
the most appropriate for the \texttt{Left} error branch if you don't
actually need to provide an error message?

View Solution

When we simply want to signal that a computation failed without
attaching any explanation, the Unit type \texttt{()} (which has exactly
1 inhabitant) is the perfect minimal error type!

\begin{Shaded}
\begin{Highlighting}[]
\OtherTok{safeHead ::}\NormalTok{ [a] }\OtherTok{{-}\textgreater{}} \DataTypeTok{Either}\NormalTok{ () a}
\NormalTok{safeHead []    }\OtherTok{=} \DataTypeTok{Left}\NormalTok{ ()}
\NormalTok{safeHead (x}\OperatorTok{:}\NormalTok{\_) }\OtherTok{=} \DataTypeTok{Right}\NormalTok{ x}
\end{Highlighting}
\end{Shaded}

\emph{(Note: Haskell's \texttt{Maybe\ a} is effectively isomorphic to
\texttt{Either\ ()\ a} and is historically preferred for this exact
scenario!)}

\subparagraph{\texorpdfstring{Exercise 8: Avoiding \texttt{fromJust}
with
\texttt{Either}}{Exercise 8: Avoiding fromJust with Either}}\label{exercise-8-avoiding-fromjust-with-either}

Similarly, \texttt{fromJust\ ::\ Maybe\ a\ -\textgreater{}\ a} forces an
extraction and crashes if given \texttt{Nothing}. If you need to
integrate a function returning a \texttt{Maybe\ a} into a pipeline that
strictly uses \texttt{Either} to handle failures safely, how would you
convert it without risking a crash?

View Solution

We can eliminate the risk of a crash by pattern matching to convert the
\texttt{Nothing} case into our 1-inhabitant minimal error type
\texttt{()}, and wrapping the \texttt{Just} value into a \texttt{Right}!

\begin{Shaded}
\begin{Highlighting}[]
\OtherTok{safeFromJust ::} \DataTypeTok{Maybe}\NormalTok{ a }\OtherTok{{-}\textgreater{}} \DataTypeTok{Either}\NormalTok{ () a}
\NormalTok{safeFromJust }\DataTypeTok{Nothing}  \OtherTok{=} \DataTypeTok{Left}\NormalTok{ ()}
\NormalTok{safeFromJust (}\DataTypeTok{Just}\NormalTok{ x) }\OtherTok{=} \DataTypeTok{Right}\NormalTok{ x}
\end{Highlighting}
\end{Shaded}

\subsubsection{\texorpdfstring{Section 1.3: \texttt{Bool} (2 Inhabitants
/ Coproduct of Terminal
Objects)}{Section 1.3: Bool (2 Inhabitants / Coproduct of Terminal Objects)}}\label{section-1.3-bool-2-inhabitants-coproduct-of-terminal-objects}

\textbf{Categorical Analog:} \emph{The Coproduct (Sum) of two Terminal
Objects} (denoted as \(1 + 1 = 2\)). In Category theory, adding two
single-element sets together creates a two-element set. This
fundamentally models branching logic (e.g., ``Left or Right'', ``True or
False'').

A 2-inhabitants type represents a binary choice. It contains exactly two
distinct values.

\paragraph{\texorpdfstring{1. The Standard
\texttt{Bool}}{1. The Standard Bool}}\label{the-standard-bool}

The most ubiquitous 2-inhabitant type is \texttt{Bool}, which represents
logical truth.

\begin{Shaded}
\begin{Highlighting}[]
\KeywordTok{data} \DataTypeTok{Bool} \OtherTok{=} \DataTypeTok{False} \OperatorTok{|} \DataTypeTok{True}
\end{Highlighting}
\end{Shaded}

\paragraph{\texorpdfstring{2. Using
\texttt{Either\ ()\ ()}}{2. Using Either () ()}}\label{using-either}

Just as tuples \texttt{(a,\ b)} represent \textbf{multiplication}
(Product types) in type algebra, \texttt{Either\ a\ b} represents
\textbf{addition} (Sum types). A sum type simply adds the number of
inhabitants of its branches.

Since \texttt{()} has exactly 1 inhabitant, \texttt{Either\ ()\ ()} has
\(1 + 1 = 2\) inhabitants: \texttt{Left\ ()} and \texttt{Right\ ()}. It
is perfectly structurally isomorphic to \texttt{Bool}.

\begin{Shaded}
\begin{Highlighting}[]
\KeywordTok{type} \DataTypeTok{IsTrue} \OtherTok{=} \DataTypeTok{Either}\NormalTok{ () ()}
\end{Highlighting}
\end{Shaded}

\paragraph{3. Custom Enumerations}\label{custom-enumerations}

Instead of using \texttt{Bool} everywhere (which can lead to ``boolean
blindness''), it's highly idiomatic in Haskell to define custom
2-inhabitant types.

\begin{Shaded}
\begin{Highlighting}[]
\KeywordTok{data} \DataTypeTok{SwitchState} \OtherTok{=} \DataTypeTok{On} \OperatorTok{|} \DataTypeTok{Off}
\KeywordTok{data} \DataTypeTok{AccessLevel} \OtherTok{=} \DataTypeTok{Admin} \OperatorTok{|} \DataTypeTok{User}
\end{Highlighting}
\end{Shaded}

\subsubsection{Section 1.4: Infinite Inhabitants (Countable and
Uncountable)}\label{section-1.4-infinite-inhabitants-countable-and-uncountable}

When we step beyond finite algebraic structures, we enter the realm of
infinity. However, in mathematics and type theory, not all infinities
are equal in size!

\paragraph{\texorpdfstring{1. Countably Infinite
(\(\aleph_0\))}{1. Countably Infinite (\textbackslash aleph\_0)}}\label{countably-infinite-aleph_0}

Countably infinite types have exactly as many inhabitants as the Natural
numbers (\(0, 1, 2, 3\dots\)). You can theoretically line them up and
count them one by one.

\textbf{The Base Standard:} Haskell's \texttt{Integer} is the canonical
countably infinite type, as it models unbounded whole numbers.

\textbf{The Structural Example (\texttt{{[}(){]}}):} Structurally, we
can construct a countably infinite type using a recursive Sum type.
Consider a standard list containing \emph{only} the Unit type:
\texttt{{[}(){]}}.

Because \texttt{()} carries zero information, the only piece of
information an entire \texttt{{[}(){]}} list can possibly hold is its
length. It is perfectly isomorphic to the Natural numbers: *
\texttt{{[}{]}} represents 0 * \texttt{{[}(){]}} represents 1 *
\texttt{{[}(),\ (){]}} represents 2 * \ldots and so on.

\textbf{The Unordered Pair (\texttt{UnorderedPair\ Integer}):} Another
simple example is the unordered pair we discussed previously. Consider
the type:

\begin{Shaded}
\begin{Highlighting}[]
\KeywordTok{data} \DataTypeTok{UnorderedPair}\NormalTok{ a }\OtherTok{=} \DataTypeTok{UPair}\NormalTok{ a a}
\end{Highlighting}
\end{Shaded}

If we instantiate this with a countably infinite type like
\texttt{Integer}, the resulting type \texttt{UnorderedPair\ Integer} is
also countably infinite. Even though it pairs two numbers (like
\texttt{UPair\ 1\ 2}), because order doesn't mathematically matter in an
\emph{unordered} set (meaning \texttt{UPair\ 1\ 2} is equivalent to
\texttt{UPair\ 2\ 1}), the total number of distinct pairs remains
countably infinite!

\textbf{Other Utterly Useful Envelopes:} If you need a countably
infinite type but must conform to specific structural API shapes, you
will frequently see basic infinite types wrapped in ``containers'' that
preserve their exact cardinality (\emph{Note: these containers are of
kind \texttt{Type\ -\textgreater{}\ Type} which we will explore in
Chapter 2, but when instantiated with an \texttt{Integer} they return to
kind \texttt{Type}!}): * \textbf{The Ordered Pair
(\texttt{(Integer,\ Integer)}):} Unlike our unordered pair, the standard
tuple distinguishes order, meaning \texttt{(1,\ 2)} and \texttt{(2,\ 1)}
are distinct values. Because mathematically
\(\aleph_0 \times \aleph_0 = \aleph_0\), the standard tuple remains
perfectly countably infinite. It's the ubiquitous foundation for
returning multiple values, representing 2D Euclidean coordinates, and
building Rational numbers (fractions). * \textbf{The Constant Payload
(\texttt{Const\ Integer\ String}):} \texttt{Const} is a highly
specialized, absolutely pervasive wrapper from
\texttt{Data.Functor.Const} that physically holds a value of its first
type parameter (e.g.~\texttt{Integer}) while completely ignoring and
pretending to hold its second type parameter (\texttt{String}). It
actually has exactly the same number of inhabitants as \texttt{Integer}!
It is utterly indispensable in advanced Haskell (like building Lenses or
writing Traversals), as it lets you ``hijack'' a mapping operation that
expects to modify a \texttt{String}, and instead use it to accumulate a
constant \texttt{Integer}. * \textbf{The Transparent Box
(\texttt{Identity\ Integer}):} Often, complex abstractions like Monad
Transformers require a value to be wrapped in a structural shape. The
\texttt{Identity} type perfectly fills this requirement without adding
any branching logic or effects, acting as a transparent zero-cost box
around an \texttt{Integer}.

\paragraph{\texorpdfstring{2. Uncountably Infinite
(\(2^{\aleph_0}\))}{2. Uncountably Infinite (2\^{}\{\textbackslash aleph\_0\})}}\label{uncountably-infinite-2aleph_0}

Haskell's laziness and function semantics allow us to define types that
contain an \textbf{uncountably infinite} number of possible inhabitants
(specifically \(2^{\aleph_0}\), possessing the same cardinality as the
real numbers).

To achieve an uncountable number of distinct values, a type must
represent an \emph{infinite sequence of choices}, where every single
step offers at least two branching options.

\subparagraph{A. Infinite Streams}\label{a.-infinite-streams}

\begin{Shaded}
\begin{Highlighting}[]
\CommentTok{{-}{-} Structurally infinite sequence of choices}
\KeywordTok{data} \DataTypeTok{Stream}\NormalTok{ a }\OtherTok{=} \DataTypeTok{Cons}\NormalTok{ a (}\DataTypeTok{Stream}\NormalTok{ a)}

\CommentTok{{-}{-} Uncountably infinite inhabitants}
\KeywordTok{type} \DataTypeTok{CoinFlips} \OtherTok{=} \DataTypeTok{Stream} \DataTypeTok{Bool}
\end{Highlighting}
\end{Shaded}

Every \texttt{CoinFlips} stream is an infinite sequence of
\texttt{True}/\texttt{False} decisions
(e.g.~\texttt{True,\ False,\ False,\ True...}). Because this represents
an infinite branching path of binary choices, it is mathematically
isomorphic to a real number between \(0.0\) and \(1.0\) in binary. There
is no mapping that can linearly ``count'' every possible infinite
stream; it is strictly uncountably infinite.

\subparagraph{B. Infinite Functions}\label{b.-infinite-functions}

A mathematically pure function mapping from a countably infinite domain
(like \texttt{Integer}) to a type with at least two inhabitants (like
\texttt{Bool}) is also uncountably infinite.

\begin{Shaded}
\begin{Highlighting}[]
\CommentTok{{-}{-} Uncountably infinite inhabitants}
\KeywordTok{type} \DataTypeTok{IntegerPredicate} \OtherTok{=} \DataTypeTok{Integer} \OtherTok{{-}\textgreater{}} \DataTypeTok{Bool}
\end{Highlighting}
\end{Shaded}

This is because defining a single complete function of type
\texttt{Integer\ -\textgreater{}\ Bool} requires making an infinite
number of binary choices (the outputs corresponding to the inputs
\texttt{0,\ 1,\ 2...}). Thus, the total number of unique
\texttt{Integer\ -\textgreater{}\ Bool} functions is exactly
\(2^{\aleph_0}\).

\begin{center}\rule{0.5\linewidth}{0.5pt}\end{center}

\subsection{Annex}\label{annex}

\subsubsection{\texorpdfstring{The Secret Inhabitant: Bottom
(\texttt{\_\textbar{}\_})}{The Secret Inhabitant: Bottom (\_\textbar\_)}}\label{the-secret-inhabitant-bottom-__}

In mathematical logic, it is impossible to return a value from a
function without inputs if there are no values to return. However,
Haskell is not a pure mathematical language; it is a real programming
language. Because of this, every single type in Haskell has a secret,
hidden inhabitant known as \texttt{\_\textbar{}\_} (pronounced
``bottom'').

``Bottom'' represents a computation that fails to complete successfully.
This can happen in two main ways: 1. The program crashes (e.g., by
throwing an error, or by evaluating \texttt{undefined}). 2. The program
gets stuck in an infinite loop and never returns.

If you wrote \texttt{createNever\ ::\ a\ -\textgreater{}\ Never}, you
could technically ``implement'' it by writing
\texttt{createNever\ x\ =\ undefined} or
\texttt{createNever\ x\ =\ createNever\ x}. It would compile and satisfy
the type checker, but it's fundamentally cheating because it avoids
returning altogether by crashing or looping forever! Because
\texttt{\_\textbar{}\_} inhabits every type, you can use it to satisfy
any signature, even impossible ones.

However, Haskellers typically reason about their code by assuming it
terminates and doesn't crash, treating it as if it were a total
language. This approach is formally justified in the well-known paper
\emph{``Fast and Loose Reasoning is Morally Correct''} {[}2{]}.

\paragraph{\texorpdfstring{Interacting with Bottom safely using
\texttt{IO}}{Interacting with Bottom safely using IO}}\label{interacting-with-bottom-safely-using-io}

Because pure functions in Haskell must evaluate consistently, you
\textbf{cannot} catch a crash (a bottom value like \texttt{error}) in
pure code. Doing so would make the pure function's behavior depend on
exactly \emph{when} the crash was evaluated, breaking referential
transparency.

If you absolutely must interact with legacy code that might crash, you
must acknowledge that observing a crash is a side effect. Therefore,
catching it must happen in the \texttt{IO} monad! You can use
\texttt{Control.Exception.try} combined with
\texttt{Control.Exception.evaluate} (which forces the pure computation
to run so its exceptions can be caught) to wrap a dangerous pure
function into a safe \texttt{Either}:

\begin{Shaded}
\begin{Highlighting}[]
\KeywordTok{import} \DataTypeTok{Control.Exception}\NormalTok{ (try, }\DataTypeTok{SomeException}\NormalTok{, evaluate)}

\OtherTok{legacyCrashingCall ::} \DataTypeTok{String} \OtherTok{{-}\textgreater{}} \DataTypeTok{Book}
\NormalTok{legacyCrashingCall }\OtherTok{=} \FunctionTok{undefined} \CommentTok{{-}{-} Imagine this throws an error}

\CommentTok{{-}{-} We must wrap it in IO to observe the crash safely}
\OtherTok{safeWrapper ::} \DataTypeTok{String} \OtherTok{{-}\textgreater{}} \DataTypeTok{IO}\NormalTok{ (}\DataTypeTok{Either} \DataTypeTok{SomeException} \DataTypeTok{Book}\NormalTok{)}
\NormalTok{safeWrapper title }\OtherTok{=}\NormalTok{ try (evaluate (legacyCrashingCall title))}
\end{Highlighting}
\end{Shaded}

By doing this, a complete crash (\texttt{\_\textbar{}\_}) is safely
intercepted and converted into a \texttt{Left\ SomeException} within the
\texttt{IO} boundary.

\begin{center}\rule{0.5\linewidth}{0.5pt}\end{center}

\begin{quote}
For references, papers, and further reading on these algebraic
structures, refer to \href{09_bibliography.md}{Part 9: Bibliography}.
\end{quote}

\section{Part 2: The Algebras (Laws) for Concrete
Types}\label{part-2-the-algebras-laws-for-concrete-types}

Welcome to the second part of Universe 1. In Part 1, we defined our core
\textbf{Structures}---the bare mathematical geometry of how many values
a type can hold. We looked at the Initial Object (\texttt{Void}), the
Terminal Object (\texttt{()}), the Coproduct of Terminal Objects
(\texttt{Bool}), and both Countable (e.g., \texttt{{[}(){]}},
\texttt{Integer}) and Uncountable (e.g., \texttt{Stream\ Bool},
\texttt{Integer\ -\textgreater{}\ Bool}) infinite inhabitants.

But structures alone are sterile. To actually perform computation, we
need \textbf{Algebras}. An algebra assigns specific \emph{behaviors} to
our structures. In Haskell, we implement these algebras using
Typeclasses. But unlike simple interfaces in other programming
languages, a true algebra must come with \textbf{Mathematical Laws} to
ensure the behavior is predictably sound.

In this document, we will build out the fundamental algebras that
operate directly on concrete types of kind \texttt{Type}.

\subsection{Chapter 1: Equivalence and
Ordering}\label{chapter-1-equivalence-and-ordering}

Before we can combine values or map over structures, the most
fundamental operation a computer can perform is determining if two
things are the same.

\subsubsection{\texorpdfstring{Section 1.1: \texttt{Eq} (The Laws of
Mathematical
Equivalence)}{Section 1.1: Eq (The Laws of Mathematical Equivalence)}}\label{section-1.1-eq-the-laws-of-mathematical-equivalence}

Before looking at the operations of \texttt{Eq}, we must ask: what
\emph{kinds} of types can have an \texttt{Eq} instance? Because the
operators compare fully instantiated runtime values, any type
implementing \texttt{Eq} \textbf{must be a simple kind \texttt{Type}}.
We cannot test if two uninstantiated type constructors (like
\texttt{Maybe}) are equal; we can only test if two concrete values (like
\texttt{Maybe\ Int}) are equal.

The \texttt{Eq} typeclass provides the \texttt{(==)} and \texttt{(/=)}
operators.

\begin{Shaded}
\begin{Highlighting}[]
\KeywordTok{class} \DataTypeTok{Eq}\NormalTok{ a }\KeywordTok{where}
\OtherTok{    (==) ::}\NormalTok{ a }\OtherTok{{-}\textgreater{}}\NormalTok{ a }\OtherTok{{-}\textgreater{}} \DataTypeTok{Bool}
\OtherTok{    (/=) ::}\NormalTok{ a }\OtherTok{{-}\textgreater{}}\NormalTok{ a }\OtherTok{{-}\textgreater{}} \DataTypeTok{Bool}
    \OtherTok{\{{-}\# MINIMAL (==) | (/=) \#{-}\}}
\end{Highlighting}
\end{Shaded}

Notice the \texttt{\{-\#\ MINIMAL\ (==)\ \textbar{}\ (/=)\ \#-\}}
pragma. This means we only \emph{need} to implement one of the two
operators. If we write \texttt{(==)}, Haskell provides a default
implementation \texttt{x\ /=\ y\ =\ not\ (x\ ==\ y)} (and vice versa).

To be a valid instance, it must rigorously satisfy the three
mathematical laws of an \textbf{equivalence relation}:

\begin{enumerate}
\def\labelenumi{\arabic{enumi}.}
\tightlist
\item
  \textbf{Reflexivity}: Everything is equal to itself. \texttt{x\ ==\ x}
  must be \texttt{True}.
\item
  \textbf{Symmetry}: Order of comparison doesn't matter.
  \texttt{x\ ==\ y} implies \texttt{y\ ==\ x}.
\item
  \textbf{Transitivity}: Equality chains perfectly. If \texttt{x\ ==\ y}
  and \texttt{y\ ==\ z}, then \texttt{x\ ==\ z}.
\end{enumerate}

\textbf{Testing the Laws (\texttt{tasty-quickcheck} \&
\texttt{testBatch})}

As we discussed in the Introduction, these laws are mathematically
absolute. But we don't need to manually verify them! We can write
property tests using \texttt{tasty-quickcheck} to systematically
generate random values and mechanically assert all three laws hold. For
many standard typeclasses, libraries even provide pre-built test
batches.

\begin{Shaded}
\begin{Highlighting}[]
\CommentTok{{-}{-} Automatically tests Reflexivity, Symmetry, and Transitivity!}
\NormalTok{testBatch (eq (}\FunctionTok{undefined}\OtherTok{ ::} \DataTypeTok{MyData}\NormalTok{))}
\end{Highlighting}
\end{Shaded}

\textbf{The Minimal Implementations:} - \textbf{0 Inhabitants
(\texttt{Void})}: As explored in Part 1 (Functions Returning an
Uninhabited Type), any function computing a value from \texttt{Void} is
mathematically valid via the Principle of Explosion. Since the signature
becomes
\texttt{(==)\ ::\ Void\ -\textgreater{}\ Void\ -\textgreater{}\ Bool},
substituting \texttt{Void} requires at least one \texttt{Void} as input.
Using \texttt{absurd} (or the empty case), this gives us the \emph{only}
implementation possible:
\texttt{haskell\ \ \ instance\ Eq\ Void\ where\ \ \ \ \ \ \ v1\ ==\ \_\ =\ absurd\ v1}
And because we can never instantiate the values at runtime to break
them, the property laws of \texttt{Eq} are trivially (vacuously)
satisfied: * \textbf{Reflexivity}: \(\forall v\), \(v == v\)? We can
never provide any \(v\), so the statement is vacuously true. *
\textbf{Symmetry}: \(\forall (v_1, v_2)\),
\(v_1 == v_2 \Rightarrow v_2 == v_1\)? We can never provide \(v_1\) or
\(v_2\), so yes. * \textbf{Transitivity}: \(\forall (v_1, v_2, v_3)\),
\(v_1 == v_2 \land v_2 == v_3 \Rightarrow v_1 == v_3\)? We can never
provide \(v_1\), \(v_2\), or \(v_3\), so yes.

\textbf{Exercise 1: Point-Free \texttt{(==)}} How would you manually
implement
\texttt{(==)\ ::\ Void\ -\textgreater{}\ Void\ -\textgreater{}\ Bool} in
a fully point-free style using only \texttt{absurd} and \texttt{const}?

View Solution

\texttt{absurd} has the type \texttt{Void\ -\textgreater{}\ a}, which
can be specialized to \texttt{Void\ -\textgreater{}\ Bool}.
\texttt{const} takes a value and ignores its second argument, with type
\texttt{x\ -\textgreater{}\ y\ -\textgreater{}\ x}. By passing
\texttt{absurd} to \texttt{const}, we get a function
\texttt{const\ absurd} with the signature
\texttt{y\ -\textgreater{}\ (Void\ -\textgreater{}\ Bool)}. When used in
our instance, \texttt{y} aligns with \texttt{Void}, giving
\texttt{Void\ -\textgreater{}\ Void\ -\textgreater{}\ Bool}. Thus:
\texttt{haskell\ \ \ instance\ Eq\ Void\ where\ \ \ \ \ \ \ (==)\ =\ const\ absurd}

\textbf{Exercise 2: The Derived \texttt{(/=)}} Since \texttt{Eq} has a
minimal pragma prescribing either \texttt{(==)} or \texttt{(/=)},
Haskell will automatically derive \texttt{(/=)} from our \texttt{(==)}
implementation. What is the effective full implementation of
\texttt{Eq\ Void} that the compiler generates?

View Solution

Haskell uses the default implementation
\texttt{x\ /=\ y\ =\ not\ (x\ ==\ y)}. Combined with our implementation
of \texttt{(==)}, the full effective code becomes:
\texttt{haskell\ \ \ instance\ Eq\ Void\ where\ \ \ \ \ \ \ v1\ ==\ \_\ =\ absurd\ v1\ \ \ \ \ \ \ v1\ /=\ v2\ =\ not\ (absurd\ v1)}

\begin{itemize}
\item
  \textbf{1 Inhabitant (\texttt{()})}: There is only one possible value,
  so \texttt{()\ ==\ ()} is always \texttt{True}. This trivially
  respects all the laws.

  \textbf{Exercise 3: The Trivial Inequality} Without relying on
  \texttt{(==)}, how would you implement the simplest possible
  \texttt{(/=)\ ::\ ()\ -\textgreater{}\ ()\ -\textgreater{}\ Bool}
  directly?

  View Solution

  Since both inputs must be \texttt{()}, they are always exactly the
  same value. Thus, they can never be not equal. The implementation is
  universally \texttt{False}:

\begin{Shaded}
\begin{Highlighting}[]
\KeywordTok{instance} \DataTypeTok{Eq}\NormalTok{ () }\KeywordTok{where}
\NormalTok{    () }\OperatorTok{/=}\NormalTok{ () }\OtherTok{=} \DataTypeTok{False}
\CommentTok{{-}{-} or simply:}
\CommentTok{{-}{-}  \_ /= \_ = False}
\end{Highlighting}
\end{Shaded}

  \textbf{Exercise 4: Breaking the Unit Laws} Is it possible to write a
  mathematically invalid \texttt{Eq} instance for \texttt{()}? If so,
  what is it and which law does it break?

  View Solution

  Yes! Although there's only one way to define equality that obeys the
  mathematical laws, Haskell still lets you write whatever code you
  want:

\begin{Shaded}
\begin{Highlighting}[]
\KeywordTok{instance} \DataTypeTok{Eq}\NormalTok{ () }\KeywordTok{where}
\NormalTok{    () }\OperatorTok{==}\NormalTok{ () }\OtherTok{=} \DataTypeTok{False}
\end{Highlighting}
\end{Shaded}

  This immediately breaks \textbf{Reflexivity}, which strictly mandates
  that for all values \(x\), \(x == x\) must evaluate to \texttt{True}.
  Because our instance returns \texttt{False}, it is an unlawful,
  mathematically invalid \texttt{Eq}!
\item
  \textbf{2 Inhabitants (\texttt{Bool})}: \textbf{Reflexivity} strictly
  forces our hand to define \texttt{True\ ==\ True} and
  \texttt{False\ ==\ False}. This leaves us with exactly two lawful
  possibilities for how we handle the cross-comparisons
  (\texttt{True\ ==\ False}). Both of these configurations naturally
  fulfill the remaining laws (Symmetry and Transitivity), giving us two
  valid cases:

  \begin{enumerate}
  \def\labelenumi{\arabic{enumi}.}
  \tightlist
  \item
    \textbf{The Useful Case (\texttt{\_\ ==\ \_\ =\ False} for differing
    values)}: By defining the cross-comparisons to evaluate to
    \texttt{False}, we preserve \texttt{True} and \texttt{False} as two
    completely distinct semantic concepts.
  \item
    \textbf{The Non-Useful Case (\texttt{\_\ ==\ \_\ =\ True} for
    differing values)}: Mathematically, if we evaluate this to
    \texttt{True}, we construct a perfectly lawful equality where
    \texttt{True} and \texttt{False} belong to the exact same
    \emph{equivalence class}. While \texttt{Bool} would still
    structurally have two distinct memory tags internally, this
    algebraically groups them together into a single \emph{Quotient
    Type}. It mathematically collapses the concept of \texttt{Bool} into
    behaving like a 1-inhabitant type whenever evaluated through
    \texttt{(==)}. This is a perfectly correct and mathematically lawful
    instance, but it is hardly useful in any programming context since
    we would entirely lose the ability to differentiate branches!
  \end{enumerate}

  \textbf{Exercise 5: The Logic Gate} Focusing on the ``useful''
  implementation of \texttt{Eq} for \texttt{Bool}, the \texttt{(==)} and
  \texttt{(/=)} operators are both functions of type
  \texttt{Bool\ -\textgreater{}\ Bool\ -\textgreater{}\ Bool}. If you
  were building a physical circuit board, what standard logic gates do
  these two operators correspond to?

  View Solution

  \begin{itemize}
  \tightlist
  \item
    \texttt{(==)} outputs \texttt{True} only if both inputs are the same
    (both \texttt{True} or both \texttt{False}). This corresponds
    perfectly to an \textbf{XNOR (Exclusive-NOR)} gate (also called the
    logical biconditional).
  \item
    \texttt{(/=)} outputs \texttt{True} only if the inputs differ
    (\texttt{True/False} or \texttt{False/True}). This corresponds
    perfectly to an \textbf{XOR (Exclusive-OR)} gate!
  \end{itemize}
\item
  \textbf{\(\infty\) Inhabitants (\texttt{Fraction})}: When a type has
  many inhabitants, the default compiler-derived \texttt{Eq} (which
  checks identical memory structure) may not reflect true mathematical
  parity. We often have to manually implement structural equality.

  \textbf{Exercise 6: A Meaningful Custom \texttt{Eq}} Suppose you are
  working with fractions defined as a numerator and denominator:

\begin{Shaded}
\begin{Highlighting}[]
\KeywordTok{data} \DataTypeTok{Fraction} \OtherTok{=} \DataTypeTok{Fraction} \DataTypeTok{Integer} \DataTypeTok{Integer}
\end{Highlighting}
\end{Shaded}

  If we let the compiler automatically derive an \texttt{Eq} instance
  for us, it would only check if the exact fields match. Under that
  default instance, \texttt{Fraction\ 1\ 2\ ==\ Fraction\ 2\ 4} would
  evaluate to \texttt{False}. How would you write a custom \texttt{Eq}
  instance that mathematically reflects truly equivalent fractions?

  View Solution

  Two fractions \(a/b\) and \(c/d\) are equal if their
  cross-multiplication matches (\(a \times d = b \times c\)).

\begin{Shaded}
\begin{Highlighting}[]
\KeywordTok{instance} \DataTypeTok{Eq} \DataTypeTok{Fraction} \KeywordTok{where}
    \DataTypeTok{Fraction}\NormalTok{ a b }\OperatorTok{==} \DataTypeTok{Fraction}\NormalTok{ c d }\OtherTok{=}\NormalTok{ (a }\OperatorTok{*}\NormalTok{ d) }\OperatorTok{==}\NormalTok{ (b }\OperatorTok{*}\NormalTok{ c)}
\end{Highlighting}
\end{Shaded}

  By defining equality based on the mathematical properties of the
  values rather than their literal data layout, we ensure a robust
  equivalence relation that fully obeys Reflexivity, Symmetry, and
  Transitivity!
\end{itemize}

\subsubsection{\texorpdfstring{Section 1.2: \texttt{Ord} (The Laws of
Total
Ordering)}{Section 1.2: Ord (The Laws of Total Ordering)}}\label{section-1.2-ord-the-laws-of-total-ordering}

If \texttt{Eq} tells us if things are the same, \texttt{Ord} tells us
how to line them up in a sequence. \texttt{Ord} provides operations like
\texttt{compare}, \texttt{\textless{}=}, and \texttt{\textgreater{}}.

\begin{Shaded}
\begin{Highlighting}[]
\KeywordTok{data} \DataTypeTok{Ordering} \OtherTok{=} \DataTypeTok{LT} \OperatorTok{|} \DataTypeTok{EQ} \OperatorTok{|} \DataTypeTok{GT}

\KeywordTok{class} \DataTypeTok{Eq}\NormalTok{ a }\OtherTok{=\textgreater{}} \DataTypeTok{Ord}\NormalTok{ a }\KeywordTok{where}
\OtherTok{    compare ::}\NormalTok{ a }\OtherTok{{-}\textgreater{}}\NormalTok{ a }\OtherTok{{-}\textgreater{}} \DataTypeTok{Ordering}
\OtherTok{    (\textless{})     ::}\NormalTok{ a }\OtherTok{{-}\textgreater{}}\NormalTok{ a }\OtherTok{{-}\textgreater{}} \DataTypeTok{Bool}
\OtherTok{    (\textless{}=)    ::}\NormalTok{ a }\OtherTok{{-}\textgreater{}}\NormalTok{ a }\OtherTok{{-}\textgreater{}} \DataTypeTok{Bool}
\OtherTok{    (\textgreater{})     ::}\NormalTok{ a }\OtherTok{{-}\textgreater{}}\NormalTok{ a }\OtherTok{{-}\textgreater{}} \DataTypeTok{Bool}
\OtherTok{    (\textgreater{}=)    ::}\NormalTok{ a }\OtherTok{{-}\textgreater{}}\NormalTok{ a }\OtherTok{{-}\textgreater{}} \DataTypeTok{Bool}
\OtherTok{    max     ::}\NormalTok{ a }\OtherTok{{-}\textgreater{}}\NormalTok{ a }\OtherTok{{-}\textgreater{}}\NormalTok{ a}
\OtherTok{    min     ::}\NormalTok{ a }\OtherTok{{-}\textgreater{}}\NormalTok{ a }\OtherTok{{-}\textgreater{}}\NormalTok{ a}
    \OtherTok{\{{-}\# MINIMAL compare | (\textless{}=) \#{-}\}}
\end{Highlighting}
\end{Shaded}

Notice the \texttt{Ordering} type at the top of the block. This is a
standard built-in Haskell enumeration with exactly three inhabitants:
\texttt{LT} (Less Than), \texttt{EQ} (Equal To), and \texttt{GT}
(Greater Than). It exists purely to represent the concrete result of
evaluating which of two values comes first.

Just like with \texttt{Eq}, the
\texttt{\{-\#\ MINIMAL\ compare\ \textbar{}\ (\textless{}=)\ \#-\}}
pragma dictates what we need to provide. To satisfy the compiler and get
a full \texttt{Ord} instance with all seven functions, we only need to
implement \emph{either} the \texttt{compare} function \emph{or} the
\texttt{(\textless{}=)} operator. If we define \texttt{compare}, all
other functions like \texttt{\textless{}}, \texttt{\textgreater{}}, and
\texttt{max} are automatically derived from it.

It is a fundamental rule that any type with an \texttt{Ord} instance
\emph{must} also have an \texttt{Eq} instance.

Mathematically, \texttt{Ord} defines a \textbf{Total Order}. It inherits
the rules of \texttt{Eq} and adds:

\begin{enumerate}
\def\labelenumi{\arabic{enumi}.}
\tightlist
\item
  \textbf{Antisymmetry}: If \texttt{x\ \textless{}=\ y} and
  \texttt{y\ \textless{}=\ x}, then they must actually be the same value
  (\texttt{x\ ==\ y}).
\item
  \textbf{Transitivity}: If \texttt{x\ \textless{}=\ y} and
  \texttt{y\ \textless{}=\ z}, then \texttt{x\ \textless{}=\ z}.
\item
  \textbf{Strong Connexity}: For any two values, one must be smaller
  than or equal to the other (\texttt{x\ \textless{}=\ y} or
  \texttt{y\ \textless{}=\ x}). In other words, every single value in
  the type can be compared to every other value without exception.
\end{enumerate}

\textbf{The Minimal Implementations:} - \textbf{0 Inhabitants
(\texttt{Void})}: Vacuously true. - \textbf{1 Inhabitant (\texttt{()})}:
\texttt{()} is always equal to (and therefore \texttt{\textless{}=} to)
\texttt{()}. - \textbf{2 Inhabitants (\texttt{Bool})}: \texttt{False} is
canonically ordered before \texttt{True}
(\texttt{False\ \textless{}=\ True}).

\begin{itemize}
\item
  \textbf{3 Inhabitants (\texttt{RPS})}: Three inhabitants is the
  minimum number required to demonstrate a cyclic relationship, meaning
  we can mathematically break a Total Order!

  \textbf{Exercise 7: Breaking the Total Order} Consider a hypothetical
  game of Rock-Paper-Scissors. Can we mathematically construct a valid
  sequence of all choices? Let's try writing an \texttt{Ord} instance:

\begin{Shaded}
\begin{Highlighting}[]
\KeywordTok{data} \DataTypeTok{RPS} \OtherTok{=} \DataTypeTok{Rock} \OperatorTok{|} \DataTypeTok{Paper} \OperatorTok{|} \DataTypeTok{Scissors} \KeywordTok{deriving}\NormalTok{ (}\DataTypeTok{Eq}\NormalTok{)}

\KeywordTok{instance} \DataTypeTok{Ord} \DataTypeTok{RPS} \KeywordTok{where}
    \FunctionTok{compare} \DataTypeTok{Rock} \DataTypeTok{Paper} \OtherTok{=} \DataTypeTok{LT}     \CommentTok{{-}{-} Rock loses to Paper  (Rock \textless{} Paper)}
    \FunctionTok{compare} \DataTypeTok{Paper} \DataTypeTok{Scissors} \OtherTok{=} \DataTypeTok{LT} \CommentTok{{-}{-} Paper loses to Scissors (Paper \textless{} Scissors)}
    \FunctionTok{compare} \DataTypeTok{Scissors} \DataTypeTok{Rock} \OtherTok{=} \DataTypeTok{LT}  \CommentTok{{-}{-} Scissors loses to Rock (Scissors \textless{} Rock)}
    \FunctionTok{compare}\NormalTok{ x y }\OperatorTok{|}\NormalTok{ x }\OperatorTok{==}\NormalTok{ y    }\OtherTok{=} \DataTypeTok{EQ}
                \OperatorTok{|} \FunctionTok{otherwise} \OtherTok{=} \DataTypeTok{GT}
\end{Highlighting}
\end{Shaded}

  While this compiles and type-checks, which of the mathematical laws of
  a Total Order (\texttt{Ord}) does it violate?

  View Solution

  It violates \textbf{Transitivity}! Our instance defines
  \texttt{Rock\ \textless{}=\ Paper} and
  \texttt{Paper\ \textless{}=\ Scissors}. By the mathematical rule of
  Transitivity, it would follow that
  \texttt{Rock\ \textless{}=\ Scissors} must be true. However, our code
  specifically defines \texttt{compare\ Scissors\ Rock\ =\ LT} (meaning
  \texttt{Scissors\ \textless{}\ Rock}, and thus
  \texttt{Rock\ \textgreater{}\ Scissors}), making
  \texttt{Rock\ \textless{}=\ Scissors} evaluate to \texttt{False}!
  Therefore, Rock-Paper-Scissors is a mathematical \emph{cycle},
  rendering it physically impossible to fulfill a sequence of Total
  Order!

  \textbf{Exercise 8: Deriving the Rest from \texttt{compare}} Assume
  you have provided a valid
  \texttt{compare\ ::\ a\ -\textgreater{}\ a\ -\textgreater{}\ Ordering}
  for your type. How would you mathematically define the other operators
  (\texttt{\textless{}}, \texttt{\textless{}=}, \texttt{\textgreater{}},
  \texttt{\textgreater{}=}, \texttt{max}, \texttt{min}) solely in terms
  of \texttt{compare}?

  View Solution

\begin{Shaded}
\begin{Highlighting}[]
\NormalTok{x }\OperatorTok{\textless{}}\NormalTok{  y }\OtherTok{=} \FunctionTok{compare}\NormalTok{ x y }\OperatorTok{==} \DataTypeTok{LT}
\NormalTok{x }\OperatorTok{\textless{}=}\NormalTok{ y }\OtherTok{=} \FunctionTok{compare}\NormalTok{ x y }\OperatorTok{/=} \DataTypeTok{GT}
\NormalTok{x }\OperatorTok{\textgreater{}}\NormalTok{  y }\OtherTok{=} \FunctionTok{compare}\NormalTok{ x y }\OperatorTok{==} \DataTypeTok{GT}
\NormalTok{x }\OperatorTok{\textgreater{}=}\NormalTok{ y }\OtherTok{=} \FunctionTok{compare}\NormalTok{ x y }\OperatorTok{/=} \DataTypeTok{LT}

\FunctionTok{max}\NormalTok{ x y }\OtherTok{=} \KeywordTok{case} \FunctionTok{compare}\NormalTok{ x y }\KeywordTok{of}
              \DataTypeTok{LT} \OtherTok{{-}\textgreater{}}\NormalTok{ y}
\NormalTok{              \_  }\OtherTok{{-}\textgreater{}}\NormalTok{ x}

\FunctionTok{min}\NormalTok{ x y }\OtherTok{=} \KeywordTok{case} \FunctionTok{compare}\NormalTok{ x y }\KeywordTok{of}
              \DataTypeTok{GT} \OtherTok{{-}\textgreater{}}\NormalTok{ y}
\NormalTok{              \_  }\OtherTok{{-}\textgreater{}}\NormalTok{ x}
\end{Highlighting}
\end{Shaded}

  This beautifully shows how the entirety of total ordering logic neatly
  cascades out of a single comparison query!
\end{itemize}

Because we have firmly established how to compare and order concrete
values, we can finally move on to \emph{combining} them.

\begin{center}\rule{0.5\linewidth}{0.5pt}\end{center}

\subsection{\texorpdfstring{Chapter 2: Associative Binary Operations
(\(+\) and
\(\times\))}{Chapter 2: Associative Binary Operations (+ and \textbackslash times)}}\label{chapter-2-associative-binary-operations-and-times}

\subsubsection{\texorpdfstring{Section 2.1: \texttt{Semigroup} and
\texttt{Monoid}}{Section 2.1: Semigroup and Monoid}}\label{section-2.1-semigroup-and-monoid}

While \texttt{Eq} and \texttt{Ord} define relationships between static
values, \texttt{Semigroup} and \texttt{Monoid} give us a fundamental way
to dynamically \emph{aggregate} concrete values together. To be a valid
\texttt{Monoid} in Haskell, a type must satisfy two main conditions:

\paragraph{\texorpdfstring{1. A Well-Kinded Type
(\texttt{Type})}{1. A Well-Kinded Type (Type)}}\label{a-well-kinded-type-type}

Unlike Functors which must be type constructors of kind
\texttt{Type\ -\textgreater{}\ Type} (like \texttt{{[}{]}} or
\texttt{Maybe}), a Monoid must have kind \texttt{Type}. It operates on
fully saturated, concrete value types like \texttt{{[}Int{]}},
\texttt{String}, or \texttt{Sum\ Double}. You cannot have a
\texttt{Monoid} instance for a bare constructor like \texttt{Maybe},
only for a specific type like \texttt{Maybe\ Int}.

\begin{quote}
{[}!WARNING{]} \textbf{What about function types?} It is a very common
trap to look at a function type like
\texttt{Integer\ -\textgreater{}\ Bool} and intuitively guess its kind
is \texttt{Type\ -\textgreater{}\ Type} because of the single arrow. But
this is an illusion of syntax!

The arrow operator \texttt{(-\textgreater{})} is actually an infix type
constructor. It takes \emph{two} concrete types to build a final
concrete type. Thus, its base kind is
\texttt{Type\ -\textgreater{}\ Type\ -\textgreater{}\ Type}. *
\texttt{(-\textgreater{})} alone has kind
\texttt{Type\ -\textgreater{}\ Type\ -\textgreater{}\ Type} *
\texttt{(-\textgreater{})\ Integer} (partially applied) has kind
\texttt{Type\ -\textgreater{}\ Type} *
\texttt{(-\textgreater{})\ Integer\ Bool} (fully applied, normally
written as \texttt{Integer\ -\textgreater{}\ Bool}) has kind
\texttt{Type}.

Because \texttt{Integer\ -\textgreater{}\ Bool} is fully saturated and
has kind \texttt{Type}, it operates as a concrete value type and
perfectly qualifies to be a Monoid! In fact, Haskell automatically
provides a Monoid instance for any function
\texttt{a\ -\textgreater{}\ b} provided that the return type \texttt{b}
is a Monoid.
\end{quote}

\begin{quote}
{[}!NOTE{]} However, the \emph{concept} of a monoid absolutely exists
for \texttt{Type\ -\textgreater{}\ Type}! In Haskell, the monoid for
types of kind \texttt{Type\ -\textgreater{}\ Type} is captured by the
\texttt{Alternative} typeclass (which we will talk about later). If you
map the signatures conceptually, it's an exact match: \texttt{mempty}
becomes \texttt{empty\ ::\ f\ a} and \texttt{\textless{}\textgreater{}}
becomes
\texttt{\textless{}\textbar{}\textgreater{}\ ::\ f\ a\ -\textgreater{}\ f\ a\ -\textgreater{}\ f\ a}.
\end{quote}

\paragraph{2. Two Core Operations}\label{two-core-operations}

A type is a monoid if it has two operations: a 0-ary identity and a
2-ary combination. 1. \texttt{mempty\ ::\ a}: An identity ``empty''
value. 2.
\texttt{mappend\ ::\ a\ -\textgreater{}\ a\ -\textgreater{}\ a} (or
\texttt{\textless{}\textgreater{}}): A binary associative operation to
combine two values.

\textbf{The Monoid Laws and Testing} Just like Functors and
Applicatives, instances of \texttt{Monoid} must rigidly obey
mathematical laws: 1. \textbf{Left Identity}:
\texttt{mempty\ \textless{}\textgreater{}\ x\ ==\ x} 2. \textbf{Right
Identity}: \texttt{x\ \textless{}\textgreater{}\ mempty\ ==\ x} 3.
\textbf{Associativity}:
\texttt{(x\ \textless{}\textgreater{}\ y)\ \textless{}\textgreater{}\ z\ ==\ x\ \textless{}\textgreater{}\ (y\ \textless{}\textgreater{}\ z)}

\begin{quote}
{[}!NOTE{]} \textbf{Wait, what about Commutativity?} Notice that
commutativity
(\texttt{x\ \textless{}\textgreater{}\ y\ ==\ y\ \textless{}\textgreater{}\ x})
is strictly \textbf{NOT} one of the Monoid laws! If a Monoid
\emph{happens} to also be commutative (like numeric \texttt{Sum} or
\texttt{Product}), it is given a special name: an \textbf{Abelian
Monoid}. However, the vast majority of useful structural monoids in
programming are strictly non-commutative. For example, \texttt{List}
(\texttt{"A"\ \textless{}\textgreater{}\ "B"\ /=\ "B"\ \textless{}\textgreater{}\ "A"}),
\texttt{First} (keeps the first \texttt{Just} value), and \texttt{Endo}
(function composition \(f \circ g \neq g \circ f\)) rigidly obey
associativity but deliberately break commutativity!
\end{quote}

\textbf{Developer Responsibility}: The Haskell compiler will perfectly
compile a \texttt{Monoid} instance even if it violently breaks these
laws! It is solely the developer's responsibility to ensure algebraic
correctness.

However, notice that these laws rely on strict \texttt{==} equality
(unlike the natural Left/Right identities of Bifunctors from Chapter 1,
which strictly relied on structural isomorphism \(\cong\)). Because they
rely on simple equality, we can trivially automate their validation
using \texttt{tasty-quickcheck} and \texttt{testBatch}:

\begin{Shaded}
\begin{Highlighting}[]
\CommentTok{{-}{-} Automatically tests Associativity and Left/Right Identity!}
\NormalTok{testBatch (monoid (}\FunctionTok{undefined}\OtherTok{ ::} \DataTypeTok{All}\NormalTok{))}
\end{Highlighting}
\end{Shaded}

What are the top minimal implementations of a Monoid? Of course, because
we mathematically require an identity value, a monoid cannot be
\texttt{Void} (a type with 0 inhabitants).

\paragraph{1. The Absolute Minimum (1
Inhabitant)}\label{the-absolute-minimum-1-inhabitant}

\textbf{The Unit \texttt{()}}: There is only one value, so
\texttt{mempty\ =\ ()} and
\texttt{()\ \textless{}\textgreater{}\ ()\ =\ ()}. This is the absolute
minimum implementation of a Monoid, and the identity and associativity
laws are trivially fulfilled because \texttt{()} is the only possible
value.

\paragraph{\texorpdfstring{2. Types with 2 Inhabitants
(\texttt{Bool})}{2. Types with 2 Inhabitants (Bool)}}\label{types-with-2-inhabitants-bool}

A type with exactly 2 values (like \texttt{Bool} with \texttt{True} and
\texttt{False}) has \(2 \times 2 = 4\) possible input combinations for a
binary function. For each input, it must choose one of 2 outputs,
yielding \(2^4 = 16\) mathematically possible binary operations.

Here is the exhaustive list of all 16 possible logical operations for a
Boolean type: 1. \textbf{Contradiction} (⊥): Always returns
\texttt{False} (ignores both inputs). 2. \textbf{NOR} (↓): Returns
\texttt{True} only if both are \texttt{False}. 3. \textbf{Converse
Nonimplication} (↚): Returns \texttt{True} only if \(B\) is True and
\(A\) is False. 4. \textbf{Negation A} (¬A): Always returns
\texttt{Not\ A} (ignores the second argument). 5. \textbf{Material
Nonimplication} (↛): Returns \texttt{True} only if \(A\) is True and
\(B\) is False. 6. \textbf{Negation B} (¬B): Always returns
\texttt{Not\ B} (ignores the first argument). 7. \textbf{XOR} (⊕):
Returns \texttt{True} if inputs are different. 8. \textbf{NAND} (↑):
Returns \texttt{False} only if both are \texttt{True}. 9. \textbf{AND}
(∧): Returns \texttt{True} only if both are \texttt{True}. 10.
\textbf{Equivalence} (↔): Returns \texttt{True} if inputs are the same.
11. \textbf{Projection B} (B): Always returns \(B\) (ignores the first
argument). 12. \textbf{Material Implication} (→): Returns \texttt{False}
only if \(A\) is True and \(B\) is False. 13. \textbf{Projection A} (A):
Always returns \(A\) (ignores the second argument). 14. \textbf{Converse
Implication} (←): Returns \texttt{False} only if \(B\) is True and \(A\)
is False. 15. \textbf{OR} (∨): Returns \texttt{True} if at least one is
\texttt{True}. 16. \textbf{Tautology} (⊤): Always returns \texttt{True}
(ignores both inputs).

In fact, any 2-inhabitant operation that possesses a valid two-sided
identity is mathematically \emph{guaranteed} to be associative! (See the
mathematical proof of this anomaly in \textbf{Annex A}).

To find our Monoids, we can mathematically filter these down by
rigorously testing the identity laws!

\textbf{1. Which ones fail the Left Identity requirement?
(\(e \diamond x = x\))} An operation must have some constant \(e\)
(\texttt{True} or \texttt{False}) that leaves the right side \(x\)
unchanged. Exactly \textbf{9 operations utterly fail} to have a left
identity. These include the ones that ignore the right argument
(Contradiction, Tautology, Projection A, Negation A), as well as NOR,
NAND, Negation B, Material Nonimplication (↛), and Converse Implication
(←). Discarding those 9 leaves us with exactly 7 operations possessing a
valid left identity.

\textbf{2. Which ones fail the Right Identity requirement?
(\(x \diamond e = x\))} Of the 7 surviving operations, 3 of them fail to
have a corresponding right identity element: * \textbf{Projection B}
(B): Has a left identity but evaluation always yields \(e \neq x\) on
the right. * \textbf{Material Implication} (→): \texttt{T\ →\ x\ =\ x}
(Left Identity is \texttt{T}), but \texttt{x\ →\ T\ =\ True} (Fails
Right Identity). * \textbf{Converse Nonimplication} (↚):
\texttt{F\ ↚\ x\ =\ x} (Left Identity is \texttt{F}), but
\texttt{x\ ↚\ F\ =\ False} (Fails Right Identity).

Discarding those 3 leaves us with exactly 4 operations that possess a
complete, \textbf{two-sided} identity element. At parameter size 2,
proving that these surviving 4 operations also satisfy the final Monoid
Law (Associativity) becomes trivial.

These remaining 4 operations perfectly form our 4 Boolean Monoids: *
\textbf{Boolean \texttt{All} (AND)} (\texttt{\&\&}) * \textbf{Boolean
\texttt{Any} (OR)} (\texttt{\textbar{}\textbar{}}) * \textbf{Boolean
Equivalence (XNOR)} (\texttt{==}) * \textbf{Boolean Exclusive OR (XOR)}
(\texttt{/=})

\paragraph{\texorpdfstring{3. Types with 3 Inhabitants (e.g.,
\texttt{Ordering})}{3. Types with 3 Inhabitants (e.g., Ordering)}}\label{types-with-3-inhabitants-e.g.-ordering}

What happens when we jump to a type with exactly 3 values (like
\texttt{LT}, \texttt{EQ}, \texttt{GT})? We witness a massive
combinatorial explosion, but it is still small enough to mathematically
map out!

\begin{enumerate}
\def\labelenumi{\arabic{enumi}.}
\tightlist
\item
  \textbf{Total Binary Operations}: A binary function takes two
  arguments, so there are \(3 \times 3 = 9\) possible input combinations
  \texttt{(x,\ y)}. For each of those 9 inputs, the function must choose
  one of 3 outputs. This yields \(3^9 = \mathbf{19,683}\) mathematically
  possible binary operations!
\item
  \textbf{Operations with an Identity}: To be a Monoid, we must possess
  an identity element. We have 3 choices for our identity (let's pick
  \texttt{EQ}). By setting \texttt{EQ} as the identity, we instantly
  lock in the required answers for 5 of our 9 input pairs (e.g.,
  \texttt{(LT,\ EQ)\ -\textgreater{}\ LT}). This leaves only 4 remaining
  input pairs where we still have the freedom to choose any of the 3
  outputs. Therefore, there are exactly \(3^4 = 81\) operations where
  \texttt{EQ} is the identity. Since any of the 3 elements could have
  been chosen, there are exactly \(3 \times 81 = \mathbf{243}\) total
  operations that possess a valid Identity Element (these are known
  mathematically as \emph{Unital Magmas}).
\item
  \textbf{Operations that are Associative}: Out of those 243 Unital
  Magmas, we must filter out any that break the Law of Associativity
  \texttt{(x\ \textless{}\textgreater{}\ y)\ \textless{}\textgreater{}\ z\ ==\ x\ \textless{}\textgreater{}\ (y\ \textless{}\textgreater{}\ z)}.
  If we explicitly calculate this for all 27 possible combinations of
  \texttt{(x,\ y,\ z)}, the math reveals that exactly \textbf{33} of
  them survive.
\end{enumerate}

\begin{quote}
{[}!TIP{]} \textbf{The Loss of the Mathematical Freebie} Notice a
fascinating anomaly here! In the 2-inhabitant (\texttt{Bool}) case,
every single one of the 4 operations that possessed an identity
\emph{automatically} passed associativity. You get associativity
completely ``for free''.

However, the moment you jump up to 3 inhabitants, this mathematical
freebie violently vanishes. Out of the 243 operations that possessed a
perfect identity element, a massive \textbf{210 operations} (243 - 33)
had to be discarded \emph{specifically} because they broke the Law of
Association!
\end{quote}

Therefore, for a type with 3 values, out of 19,683 possible operations,
exactly \textbf{33 form perfectly valid Monoids}!

\paragraph{\texorpdfstring{4. Types with Countably Infinite Inhabitants
(e.g.,
\texttt{Integer})}{4. Types with Countably Infinite Inhabitants (e.g., Integer)}}\label{types-with-countably-infinite-inhabitants-e.g.-integer}

What if the type has an infinite number of values? In this case, there
are an \textbf{infinite} number of valid Monoids. For example, for
standard numeric types (\texttt{Integer}), you trivially have
\texttt{Sum} (\(0, \mathbf{+}\)) and \texttt{Product}
(\(1, \mathbf{\times}\)), but also \texttt{Max} (\(-\infty, \max\)) and
\texttt{Min} (\(\infty, \min\)), along with infinite logical bitwise
operations like \texttt{And} and \texttt{Xor}.

\paragraph{\texorpdfstring{5. The Free Monoid
(\texttt{{[}a{]}})}{5. The Free Monoid ({[}a{]})}}\label{the-free-monoid-a}

An incredibly special case of a countably infinite type is the List
(\texttt{{[}a{]}}). Lists form what mathematicians term the \textbf{Free
Monoid} over a set \texttt{a}. A ``Free'' object in algebra is one that
satisfies the minimal laws required, and absolutely nothing else.

By concatenating elements end-to-end (\texttt{++}), lists perfectly obey
the Monoid laws: * \texttt{{[}{]}\ ++\ xs\ =\ xs} (Left Identity) *
\texttt{xs\ ++\ {[}{]}\ =\ xs} (Right Identity) *
\texttt{(as\ ++\ bs)\ ++\ cs\ =\ as\ ++\ (bs\ ++\ cs)} (Associativity)

The List monoid does no ``computation'' or ``squashing'' like
\texttt{Sum} or \texttt{Product} do; it purely memorizes the order and
elements over time. Because it is the ``purest'' monoid with no
additional baggage, you can map any List into \emph{any other valid
Monoid} using \texttt{foldMap}:

\begin{Shaded}
\begin{Highlighting}[]
\CommentTok{{-}{-} The Free Monoid perfectly translates into any other Monoid:}
\NormalTok{sumList xs }\OtherTok{=}\NormalTok{ getSum (}\FunctionTok{foldMap} \DataTypeTok{Sum}\NormalTok{ xs)}
\NormalTok{mulList xs }\OtherTok{=}\NormalTok{ getProduct (}\FunctionTok{foldMap} \DataTypeTok{Product}\NormalTok{ xs)}
\end{Highlighting}
\end{Shaded}

In fact, \texttt{Foldable} is entirely defined by a data structure's
ability to be collapsed down into this Free Monoid!

\paragraph{\texorpdfstring{6. Why do \texttt{Sum}, \texttt{Product},
\texttt{Max}, and \texttt{Min} stand
out?}{6. Why do Sum, Product, Max, and Min stand out?}}\label{why-do-sum-product-max-and-min-stand-out}

You might notice that while there are infinitely many ways to combine
integers, we almost always reach for these four. What makes them
``atomic''?

Just as Functors can be built from ``atoms'' (Identity, Constant,
Either, Pair) using composition, these Monoids are the \textbf{natural
algebraic projections} of underlying structures:

\begin{enumerate}
\def\labelenumi{\arabic{enumi}.}
\tightlist
\item
  \textbf{Additive/Multiplicative Monoids}: These are derived from the
  fact that \texttt{Integer} is a \textbf{Semiring}. A Semiring is a
  type with two monoidal operations that interact via the Distributive
  Law (\(a \times (b + c) = a \times b + a \times c\)).
\item
  \textbf{Max/Min Monoids}: These are derived from the fact that
  \texttt{Integer} is a \textbf{Bounded Lattice}. Any type with a total
  ordering (\texttt{Ord}) can form a Monoid using the ``least upper
  bound'' (\texttt{max}) or ``greatest lower bound'' (\texttt{min}).
\end{enumerate}

In this sense, these monoids aren't arbitrary; they are the
\textbf{unique} ways to satisfy the Monoid laws while preserving the
deeper algebraic relationships (like distribution or ordering) already
present in the type.

\textbf{Is Parametricity Helping Here?} Unlike Functors
(\texttt{Type\ -\textgreater{}\ Type}), which are parameterized over
\emph{any} type, Monoids operate on concrete types (\texttt{Type}). This
means parametricity \emph{does not} force a single, unique
implementation. For example, the type \texttt{Double} could form a
monoid under addition (\texttt{0} and \texttt{+}) or under
multiplication (\texttt{1} and \texttt{*}). Haskell uses
\texttt{newtype} wrappers like \texttt{Sum} and \texttt{Product} to
explicitly choose the monoidal behavior.

\paragraph{Category Theory Origin: The Single-Object
Category}\label{category-theory-origin-the-single-object-category}

In Category Theory, a Monoid is rigorously defined as a \textbf{Category
with exactly one object}.

If a category only has a single object, what do the morphisms (the
arrows) represent? Because there are no other objects to point to, every
single morphism must be an endomorphism (an arrow pointing from the
object back to itself).

This yields a profound translation: 1. \textbf{The Identity Value
(\texttt{mempty})}: Every category must have a mandated identity
morphism for each object. In our single-object category, this single
identity morphism geometrically represents \texttt{mempty}! 2.
\textbf{The Binary Operation (\texttt{\textless{}\textgreater{}})}:
Every category must allow the composition of morphisms (\(\circ\)).
Because all our morphisms start and end at the exact same single object,
\emph{any} two morphisms can be freely composed with each other. This
categorical composition perfectly corresponds to \texttt{mappend}! 3.
\textbf{The Laws}: The fundamental categorical laws of Morphism Identity
and Morphism Associativity directly translate into our Monoid laws!

Therefore, every time you combine two concrete values using
\texttt{\textless{}\textgreater{}}, you are musically and mathematically
composing two structural morphisms on a single, invisible categorical
object.

\section{Part 3: The Parameterized Structures (No
Laws)}\label{part-3-the-parameterized-structures-no-laws}

\textbf{Author:} Olivier De Wolf, odewolf@gmail.com

\subsection{\texorpdfstring{Chapter 2: Parameterized Types of Kind
\texttt{Type\ -\textgreater{}\ Type}}{Chapter 2: Parameterized Types of Kind Type -\textgreater{} Type}}\label{chapter-2-parameterized-types-of-kind-type---type}

These are type constructors that require one type argument \texttt{a}
before they become concrete types. Because they take another type as an
argument, they are categorically referred to as \textbf{Higher-Kinded
Types (HKTs)}. The number of inhabitants discussed here applies
\emph{regardless} of what \texttt{a} is instantiated to (i.e.~the type
parameter \texttt{a} is completely ignored at the value level).

\subsubsection{\texorpdfstring{Section 2.1: \texttt{VoidFoldable} (0
Inhabitants)}{Section 2.1: VoidFoldable (0 Inhabitants)}}\label{section-2.1-voidfoldable-0-inhabitants}

These parameterized types cannot be constructed, no matter what
\texttt{a} is.

\paragraph{1. Standard Parameterized Empty
Data}\label{standard-parameterized-empty-data}

\begin{Shaded}
\begin{Highlighting}[]
\KeywordTok{data} \DataTypeTok{VoidFoldable}\NormalTok{ a}
\end{Highlighting}
\end{Shaded}

\paragraph{2. Using GADT Syntax}\label{using-gadt-syntax}

\begin{Shaded}
\begin{Highlighting}[]
\OtherTok{\{{-}\# LANGUAGE GADTs \#{-}\}}
\KeywordTok{data} \DataTypeTok{VoidFoldable}\NormalTok{ a }\KeywordTok{where}\NormalTok{ \{\}}
\end{Highlighting}
\end{Shaded}

\paragraph{\texorpdfstring{3. Phantom Wrapping
\texttt{Data.Void}}{3. Phantom Wrapping Data.Void}}\label{phantom-wrapping-data.void}

\begin{Shaded}
\begin{Highlighting}[]
\KeywordTok{import} \DataTypeTok{Data.Void}\NormalTok{ (}\DataTypeTok{Void}\NormalTok{)}
\KeywordTok{newtype} \DataTypeTok{VoidFoldable}\NormalTok{ a }\OtherTok{=} \DataTypeTok{VoidFoldable} \DataTypeTok{Void}
\end{Highlighting}
\end{Shaded}

\paragraph{4. Reusing Standard Library
Structures}\label{reusing-standard-library-structures}

GHC provides existing parameterized empty types for generic programming,
like \texttt{V1}, or we can combine \texttt{Const} and \texttt{Void}.

\begin{Shaded}
\begin{Highlighting}[]
\KeywordTok{import} \DataTypeTok{GHC.Generics}\NormalTok{ (}\DataTypeTok{V1}\NormalTok{)}
\KeywordTok{import} \DataTypeTok{Data.Functor.Const}\NormalTok{ (}\DataTypeTok{Const}\NormalTok{)}
\KeywordTok{import} \DataTypeTok{Data.Void}\NormalTok{ (}\DataTypeTok{Void}\NormalTok{)}

\CommentTok{{-}{-} V1 a }
\CommentTok{{-}{-} Const Void a}
\end{Highlighting}
\end{Shaded}

\subsubsection{\texorpdfstring{Section 2.2: \texttt{Proxy} (1
Inhabitant)}{Section 2.2: Proxy (1 Inhabitant)}}\label{section-2.2-proxy-1-inhabitant}

These parameterized types have exactly one value, irrespective of
\texttt{a}.

\paragraph{\texorpdfstring{1.
\texttt{Data.Proxy}}{1. Data.Proxy}}\label{data.proxy}

\texttt{Proxy} is used to pass \emph{type-level} information around at
runtime without needing an actual value of that type.

\begin{Shaded}
\begin{Highlighting}[]
\KeywordTok{import} \DataTypeTok{Data.Proxy}\NormalTok{ (}\DataTypeTok{Proxy}\NormalTok{(..))}
\CommentTok{{-}{-} The type is \textasciigrave{}Proxy a\textasciigrave{}, the only value is \textasciigrave{}Proxy\textasciigrave{}}
\OtherTok{myProxy ::} \DataTypeTok{Proxy} \DataTypeTok{Int}
\NormalTok{myProxy }\OtherTok{=} \DataTypeTok{Proxy}
\end{Highlighting}
\end{Shaded}

\paragraph{\texorpdfstring{2. \texttt{Constants} and
\texttt{Generics}}{2. Constants and Generics}}\label{constants-and-generics}

GHC generic programming uses \texttt{U1} to represent constructors with
no fields. Alteratively, \texttt{Const\ ()\ a} yields exactly 1
inhabitant.

\begin{Shaded}
\begin{Highlighting}[]
\KeywordTok{import} \DataTypeTok{GHC.Generics}\NormalTok{ (}\DataTypeTok{U1}\NormalTok{(..))}
\KeywordTok{import} \DataTypeTok{Data.Functor.Const}\NormalTok{ (}\DataTypeTok{Const}\NormalTok{(..))}

\CommentTok{{-}{-} U1 a (value is U1)}
\CommentTok{{-}{-} Const () a (value is Const ())}
\end{Highlighting}
\end{Shaded}

\subsubsection{\texorpdfstring{Section 2.3: \texttt{Const\ Bool\ a} (2
Inhabitants)}{Section 2.3: Const Bool a (2 Inhabitants)}}\label{section-2.3-const-bool-a-2-inhabitants}

These parameterized types have precisely two values, regardless of
\texttt{a}.

\paragraph{1. Custom Parameterized
Tags}\label{custom-parameterized-tags}

\begin{Shaded}
\begin{Highlighting}[]
\KeywordTok{data} \DataTypeTok{TwoOptions}\NormalTok{ a }\OtherTok{=} \DataTypeTok{Option1} \OperatorTok{|} \DataTypeTok{Option2}
\end{Highlighting}
\end{Shaded}

\paragraph{\texorpdfstring{2.
\texttt{Const\ Bool\ a}}{2. Const Bool a}}\label{const-bool-a}

The \texttt{Const} functor holding a \texttt{Bool} gives exactly two
possible states.

\begin{Shaded}
\begin{Highlighting}[]
\KeywordTok{import} \DataTypeTok{Data.Functor.Const}\NormalTok{ (}\DataTypeTok{Const}\NormalTok{(..))}

\CommentTok{{-}{-} Const False :: Const Bool a}
\CommentTok{{-}{-} Const True  :: Const Bool a}
\end{Highlighting}
\end{Shaded}

\begin{center}\rule{0.5\linewidth}{0.5pt}\end{center}

\section{Part 4: The Algebras of
Shape}\label{part-4-the-algebras-of-shape}

\subsection{1. Introduction}\label{introduction-1}

Walking through the exercise of constructing ``minimal'' instances is
one of the best ways to deeply understand Functors, Applicatives, and
Monads in Haskell. By stripping away domain-specific noise (like state
management, I/O, or failure), we demystify a lot of features that
initially look like magic. It reveals the underlying mechanics at play.

In mathematics, there is a beautiful, recurring pattern: we like to
start with the absolute simplest ``atoms'' (axiomatic primitives) and
establish a clear set of combinators to build more complicated
structures. We then look for the \emph{closure}---the minimal set that
contains all those starting axioms and remains perfectly valid under
every possible combination of those operations.

In this document, we will apply exactly that mathematical lens. We will
start by defining the absolute minimal ``atomic'' Functors and
Bifunctors. Next, we will introduce the combinators of our algebra (Sums
and Products). Finally, by exploring the closure of these operations, we
will demonstrate how you can transparently build incredibly complex,
robust algebraic data types (Molecules) without ever breaking the
foundational laws of the atoms.

\subsubsection{\texorpdfstring{The Updated \(\to\) Ranking (Simplicity
to
Power)}{The Updated \textbackslash to Ranking (Simplicity to Power)}}\label{the-updated-to-ranking-simplicity-to-power}

\begin{longtable}[]{@{}
  >{\raggedright\arraybackslash}p{(\columnwidth - 6\tabcolsep) * \real{0.2500}}
  >{\raggedright\arraybackslash}p{(\columnwidth - 6\tabcolsep) * \real{0.2500}}
  >{\raggedright\arraybackslash}p{(\columnwidth - 6\tabcolsep) * \real{0.2500}}
  >{\raggedright\arraybackslash}p{(\columnwidth - 6\tabcolsep) * \real{0.2500}}@{}}
\toprule\noalign{}
\begin{minipage}[b]{\linewidth}\raggedright
Rank
\end{minipage} & \begin{minipage}[b]{\linewidth}\raggedright
Type Class
\end{minipage} & \begin{minipage}[b]{\linewidth}\raggedright
Core Idea
\end{minipage} & \begin{minipage}[b]{\linewidth}\raggedright
Why it's here
\end{minipage} \\
\midrule\noalign{}
\endhead
\bottomrule\noalign{}
\endlastfoot
1 & Functor & Mapping & Only 1 law/method. Purely transformative. \\
2 & Foldable & Reducing & High utility, intuitive ``summary'' logic. \\
3 & Applicative & Multi-context & Applying functions across multiple
containers. \\
4 & Alternative & Selection & Adds ``OR'' logic and ``Failure'' to
Applicative. \\
5 & Monad & Sequencing & Introduces ``Flattening'' and step-by-step
dependency. \\
6 & Traversable & Commuting & The most abstract; requires understanding
all above. \\
\end{longtable}

While the core concepts structured here are foundational to modern
Haskell, this specific teaching narrative---starting with absolute
minimalism to actually ``prove'' the forced hand of parametricity---is
something usually only found scattered across different resources. We
will synthesize foundational ideas found in Philip Wadler's
\emph{``Theorems for free!''} and Sandy Maguire's \emph{``Thinking with
Types''}.

\textbf{Intended Audience:} This journey is designed for mathematicians,
computer scientists, or intermediate Haskell programmers who already
grasp the basic syntax and perhaps have a surface-level intuition of
Category Theory or Abstract Algebra. If you have ever used a Functor or
a Monad but felt a lingering desire to derive them from the absolute
mathematical ``scratch''---to build an unshakeable, axiomatic
understanding of \emph{why} they must exist and behave exactly as they
do---this exploration is for you!

In this exploration, our scope is specific: we are focusing entirely on
Endofunctors operating within the category of Haskell types (from
\texttt{Hask} to \texttt{Hask}).

\begin{quote}
\textbf{Note on \texttt{Hask}}: Technically, \texttt{Hask} is not a
strict mathematical category due to non-terminating programs
(represented by \texttt{\_\textbar{}\_} or ``bottom''). For the purposes
of reasoning about types, we generally ignore this, a practice validated
by the famous paper
\href{https://www.cse.chalmers.se/~nad/publications/danielsson-et-al-popl2006.pdf}{``Fast
and Loose Reasoning is Morally Correct''}.
\end{quote}

The true protagonist of this journey is \textbf{Parametricity}. Due to
parametric polymorphism (the inability to inspect types at runtime), the
implementation of most functor, applicative, and monad instances for
simple structures is mathematically forced to be unique. This provides
immense ``intellectual economy'' for Haskell developers: operations like
\texttt{bind}, \texttt{pure}, \texttt{fmap}, \texttt{apply}
(\texttt{\textless{}*\textgreater{}}), and Kleisli composition
(\texttt{\textgreater{}=\textgreater{}}) generally have exactly one
possible correct implementation for simple structural types. The
compiler practically writes the code for you.

\emph{(Note: There are rare counterexamples where multiple valid
implementations might exist---for example, traversing a complex tree
structure in different orders, or the \texttt{List} monad which actually
has exactly two valid implementations for \texttt{bind}---but these are
atypical for the minimal types we are exploring).}

\begin{center}\rule{0.5\linewidth}{0.5pt}\end{center}

\subsection{Chapter 1: Functor \& Bifunctor (Shape
Preservation)}\label{chapter-1-functor-bifunctor-shape-preservation}

\subsubsection{Section 1.1: What is a
Functor?}\label{section-1.1-what-is-a-functor}

If you ask a mathematician, they will point you to Saunders Mac Lane,
one of the founders of Category Theory. In Category Theory, a functor is
a structure-preserving mapping between two categories. It is a
ubiquitous concept in mathematics; for instance, you have
\emph{Forgetful functors} (which strip algebraic structure) and
\emph{Free functors} (which automatically build algebraic structure).

In Haskell, the \texttt{Functor} typeclass is a specific implementation
of a categorical functor. To be a valid \texttt{Functor} in Haskell, you
must satisfy three distinct conditions:

\paragraph{\texorpdfstring{1. A Well-Kinded Type Constructor
(\texttt{Type\ -\textgreater{}\ Type})}{1. A Well-Kinded Type Constructor (Type -\textgreater{} Type)}}\label{a-well-kinded-type-constructor-type---type}

You must be an \emph{Endofunctor} on the category \texttt{Hask}. This
means you map from \texttt{Hask} back to \texttt{Hask}.

\begin{itemize}
\tightlist
\item
  \emph{Invalid Kind}: \texttt{Int} (kind \texttt{Type}) or \texttt{(,)}
  (kind \texttt{Type\ -\textgreater{}\ Type\ -\textgreater{}\ Type}) are
  not functors on their own. They don't have the right ``shape'' to be a
  container/wrapper. A functor must be a ``context'' that can hold any
  type \texttt{a}.
\end{itemize}

\paragraph{\texorpdfstring{2. Unconstrained Morphism Mapping
(\texttt{fmap})}{2. Unconstrained Morphism Mapping (fmap)}}\label{unconstrained-morphism-mapping-fmap}

You must provide a function
\texttt{fmap\ ::\ (a\ -\textgreater{}\ b)\ -\textgreater{}\ f\ a\ -\textgreater{}\ f\ b}.
This is the implementation of how ``arrows'' are mapped between
categories. Crucially, in Haskell, this mapping must be
\textbf{unconstrained}: it must work for \emph{any} type \texttt{a} and
\texttt{b}. You cannot require headers or properties (like \texttt{Eq}
or \texttt{Ord}).

This restriction is forced by \textbf{Parametricity}. When we write a
polymorphic function in Haskell, the function must be completely
ignorant of the types going into it.

If we have a generic type \texttt{a} and need to produce a generic type
\texttt{b}, we cannot inspect the value, switch on its type, or conjure
a \texttt{b} out of thin air. This drastic restriction essentially
forces our implementations to \textbf{preserve structure}. This concept
is famously codified in Philip Wadler's paper
\href{https://people.mpi-sws.org/~dreyer/tor/papers/wadler.pdf}{``Theorems
for free!''}, which proves that simply reading the type signature of a
polymorphic function tells you almost everything about what the function
physically \emph{must} do.

\paragraph{3. Mathematical Laws}\label{mathematical-laws}

You must satisfy the Identity and Composition laws: 1. \textbf{Identity
Law}: \texttt{fmap\ id\ ==\ id} 2. \textbf{Composition Law}:
\texttt{fmap\ (f\ .\ g)\ ==\ (fmap\ f)\ .\ (fmap\ g)}

\begin{quote}
{[}!IMPORTANT{]} \textbf{The Parametricity Shortcut}: A remarkable
result from Category Theory and Haskell's type system is that \textbf{if
a parametric function satisfies the Identity Law, it automatically
satisfies the Composition Law.}

This stems from the fact that \texttt{fmap} is a parametrically
polymorphic function. Its behavior is so constrained by its type
signature that it cannot ``sneak in'' extra logic that would
specifically target composed functions differently than identity. This
is a core result of Philip Wadler's famous paper:
\href{https://people.mpi-sws.org/~dreyer/tor/papers/wadler.pdf}{\textbf{``Theorems
for free!''}}. A formal proof of this is provided in the
\hyperref[proof-of-identity-implies-composition]{Annex}.
\end{quote}

While parametricity gives us elegant mathematical proofs, we can also
automate verification empirically. Using property testing libraries like
\texttt{quickcheck-classes}, verifying \texttt{Maybe} is reduced to one
simple line:

\begin{Shaded}
\begin{Highlighting}[]
\KeywordTok{import} \DataTypeTok{Test.QuickCheck.Classes}

\CommentTok{{-}{-} Automatically tests all Functor laws!}
\NormalTok{testProperties }\StringTok{"Maybe Functor"} \OperatorTok{$}\NormalTok{ functor (}\DataTypeTok{Proxy}\OtherTok{ ::} \DataTypeTok{Proxy} \DataTypeTok{Maybe}\NormalTok{)}
\end{Highlighting}
\end{Shaded}

\begin{quote}
{[}!WARNING{]} \#\#\#\# Caveat: Testing vs.~Proof (Property-Based
Testing) While \texttt{testBatch} provides extreme confidence by
checking thousands of random inputs, it is not a formal mathematical
proof. Because it relies on \textbf{Property-Based Testing}
(QuickCheck), it is probabilistic.

In Haskell, there is a fundamental difference between: 1.
\textbf{Verification (Testing)}: Checking that the laws hold for
\emph{many} random cases. 2. \textbf{Proof (Types/Parametricity)}: Using
the compiler and Category Theory (the ``Shortcut'') to guarantee the
laws hold for \emph{all} cases.
\end{quote}

\paragraph{4. Almost Functors}\label{almost-functors}

Many structures look like Functors but fail one of the strict Haskell
criteria or the mathematical laws.

\subparagraph{4.1. The Constrained
Functors}\label{the-constrained-functors}

Many structures in Haskell \emph{are} valid functors in Category Theory
but fail the Haskell unconstrained mapping condition.

\begin{enumerate}
\def\labelenumi{\arabic{enumi}.}
\item
  \textbf{The Forgetful Functor (\texttt{Monoid} -\textgreater{}
  \texttt{Hask})}:
  \texttt{haskell\ \ \ \ \ -\/-\ In\ Category\ Theory,\ this\ is\ a\ functor\ between\ different\ categories.\ \ \ \ \ forget\ ::\ Monoid\ a\ =\textgreater{}\ a\ -\textgreater{}\ a\ \ \ \ \ forget\ =\ id}

  \begin{itemize}
  \tightlist
  \item
    \emph{Haskell Status}: Not a \texttt{Functor} instance because it
    requires the \texttt{Monoid\ a} constraint.
  \end{itemize}
\item
  \textbf{The Type Inspector (\texttt{isInt})}:

\begin{Shaded}
\begin{Highlighting}[]
\KeywordTok{import} \DataTypeTok{Data.Typeable}\NormalTok{ (}\DataTypeTok{Typeable}\NormalTok{, cast)}
\KeywordTok{import} \DataTypeTok{Data.Maybe}\NormalTok{ (isJust)}

\CommentTok{{-}{-} This is a functor that inspects the object in Category Theory.}
\OtherTok{isInt ::} \DataTypeTok{Typeable}\NormalTok{ a }\OtherTok{=\textgreater{}}\NormalTok{ a }\OtherTok{{-}\textgreater{}} \DataTypeTok{Bool}
\NormalTok{isInt x }\OtherTok{=}\NormalTok{ isJust (cast}\OtherTok{ x ::} \DataTypeTok{Maybe} \DataTypeTok{Int}\NormalTok{)}
\end{Highlighting}
\end{Shaded}

  \begin{itemize}
  \tightlist
  \item
    \emph{Haskell Status}: Not a \texttt{Functor} instance because it
    requires \texttt{Typeable\ a}. It won't work for \emph{any}
    \texttt{a}, only those the compiler can reify. This is why we call
    it a ``backdoor'': it bypasses the intentional ``blindness'' of
    parametric polymorphism.
  \end{itemize}
\item
  \textbf{The Balanced Tree (\texttt{Data.Set})}:
  \texttt{haskell\ \ \ \ \ -\/-\ Rebuilds\ a\ BST\ based\ on\ new\ values.\ \ \ \ \ mapSet\ ::\ Ord\ b\ =\textgreater{}\ (a\ -\textgreater{}\ b)\ -\textgreater{}\ Set\ a\ -\textgreater{}\ Set\ b}

  \begin{itemize}
  \tightlist
  \item
    \emph{Logic}: A \texttt{Set} is implemented as a balanced
    \textbf{Binary Search Tree (BST)}. To maintain the invariant
    (ordered and unique), every map operation must rebuild the tree
    using comparisons of the \emph{new} values \texttt{b}. Since this
    requires \texttt{Ord\ b}, it is a \textbf{Restricted Functor}
    mapping to the subcategory of ordered types.
  \end{itemize}
\end{enumerate}

\textbf{The Great Synthesis: Everything is a Restricted Functor} In all
three cases above, we can ``fix'' the problem by adding a constraint
like \texttt{Monoid\ a\ =\textgreater{}},
\texttt{Typeable\ a\ =\textgreater{}}, or
\texttt{Ord\ b\ =\textgreater{}}.

\textbf{In Haskell, almost every ``non-functor'' is actually just a
functor on a subcategory.} By adding a constraint, you are explicitly
telling the compiler: ``I am no longer operating on the category of all
types (\texttt{Hask}); I am now operating only on a subcategory.'' The
standard \texttt{Functor} typeclass is simply the special case where
that subcategory is the entire category \texttt{Hask}.

If we look at valid candidates in Haskell: * \texttt{Maybe} is a valid
functor candidate (Kind \texttt{Type\ -\textgreater{}\ Type}). *
\texttt{Identity} is a valid functor candidate (Kind
\texttt{Type\ -\textgreater{}\ Type}).

\subparagraph{4.2. The Malicious Functor (Hidden
Law-Breaker)}\label{the-malicious-functor-hidden-law-breaker}

This example illustrates why testing alone isn't proof. It has the
correct signature and is parametric, but it ``hides'' its law-breaking
behavior behind a conditional:

\begin{Shaded}
\begin{Highlighting}[]
\KeywordTok{data} \DataTypeTok{MyBox}\NormalTok{ a }\OtherTok{=} \DataTypeTok{MyBox} \DataTypeTok{Int}\NormalTok{ a}

\KeywordTok{instance} \DataTypeTok{Functor} \DataTypeTok{MyBox} \KeywordTok{where}
    \FunctionTok{fmap}\NormalTok{ f (}\DataTypeTok{MyBox}\NormalTok{ x val) }
      \OperatorTok{|}\NormalTok{ x }\OperatorTok{==} \DecValTok{12345} \OtherTok{=} \DataTypeTok{MyBox}\NormalTok{ (x }\OperatorTok{+} \DecValTok{1}\NormalTok{) (f val) }\CommentTok{{-}{-} Breaking Identity}
      \OperatorTok{|} \FunctionTok{otherwise}  \OtherTok{=} \DataTypeTok{MyBox}\NormalTok{ x (f val)      }\CommentTok{{-}{-} Looking Lawful}
\end{Highlighting}
\end{Shaded}

If \texttt{testBatch} never randomly generates the integer
\texttt{12345}, this structure will \textbf{pass all your tests} while
remaining mathematically invalid!

\begin{center}\rule{0.5\linewidth}{0.5pt}\end{center}

\subsubsection{Section 1.2: Minimal
Functors}\label{section-1.2-minimal-functors}

Now that we have explored several examples of types that are \emph{not}
valid functors, let's reverse the approach. We will define the absolute
simplest, most minimal structural types we can physically imagine
building in Haskell. We will conduct this exercise for both standard
\textbf{Functors} (types with a single parameter,
\texttt{Type\ -\textgreater{}\ Type}) and \textbf{Bifunctors} (types
with two parameters,
\texttt{Type\ -\textgreater{}\ Type\ -\textgreater{}\ Type}).

The beautiful consequence of choosing structures this simple is that it
perfectly demonstrates the ``forced hand'' of \textbf{parametricity}.
Because these minimal types contain almost no data, there is
mathematically only a single possible way to map over them without
violating the type signature. Once we define the type, the compiler
practically writes the unique \texttt{Functor} and \texttt{Bifunctor}
instances for us!

These minimal structures act as the ``atoms'' from which the rest of the
algebraic universe is built.

\paragraph{\texorpdfstring{1. The Absolute Bottom:
\texttt{Zero}}{1. The Absolute Bottom: Zero}}\label{the-absolute-bottom-zero}

\emph{(Zero constructors, Zero computational data, Zero contextual data.
Mathematically, it uniquely forms the \textbf{Initial Object} of the
\texttt{Hask} category, with the usual caveat of bottom/undefined values
(\texttt{\_\textbar{}\_}) slightly muddying strict categorical purity).}

The mathematically absolute smallest possible Functor has no
constructors at all. It represents an uninhabited type---it's
mathematically impossible to construct a value of this type. It
represents total ``nothingness''.

\begin{Shaded}
\begin{Highlighting}[]
\OtherTok{\{{-}\# LANGUAGE EmptyCase \#{-}\}}
\OtherTok{\{{-}\# LANGUAGE InstanceSigs \#{-}\}}

\KeywordTok{data} \DataTypeTok{Zero}\NormalTok{ a }\CommentTok{{-}{-} No constructors! (Note: While not built{-}in by this name, an identical structure exists in base as \textasciigrave{}V1\textasciigrave{} from \textasciigrave{}GHC.Generics\textasciigrave{})}

\KeywordTok{instance} \DataTypeTok{Functor} \DataTypeTok{Zero} \KeywordTok{where}
\OtherTok{    fmap ::}\NormalTok{ (a }\OtherTok{{-}\textgreater{}}\NormalTok{ b) }\OtherTok{{-}\textgreater{}} \DataTypeTok{Zero}\NormalTok{ a }\OtherTok{{-}\textgreater{}} \DataTypeTok{Zero}\NormalTok{ b}
    \FunctionTok{fmap}\NormalTok{ \_ z }\OtherTok{=} \KeywordTok{case}\NormalTok{ z }\KeywordTok{of}\NormalTok{ \{\} }
\end{Highlighting}
\end{Shaded}

\textbf{The ``Why''}: Because \texttt{Zero\ a} has no constructors, we
can never actually instantiate it at runtime. However, the type
signature
\texttt{(a\ -\textgreater{}\ b)\ -\textgreater{}\ Zero\ a\ -\textgreater{}\ Zero\ b}
is perfectly valid. Actually, any data type with zero constructors (an
uninhabited type), regardless of how many type parameters it takes (like
\texttt{data\ Zero\ a\ b\ c}), is ALWAYS guaranteed to be a perfectly
lawful Functor, Bifunctor, Profunctor, etc. If we were somehow handed a
value \texttt{z} of type \texttt{Zero\ a}, we prove to the compiler we
can produce a \texttt{Zero\ b} by pattern matching on its non-existent
constructors, leading to an empty case. Parametricity holds because the
transformation is forced by the absolute absence of data.

\textbf{Law Verification}: * \emph{Identity}: \texttt{fmap\ id\ z} where
\texttt{z\ ::\ Zero\ a}. Pattern matching on \texttt{z} (empty case)
immediately satisfies the law as no value exists to violate it. *
\emph{Composition}: Guaranteed automatically by parametricity
(``Theorems for free!'') since the Identity law is satisfied.

\textbf{Category Theory Equivalent}: This represents the constant
functor \(\Delta_0\). It maps every object in the category space to the
Initial Object \(0\) (the empty set \(\emptyset\)) and every morphism to
the empty function \(id_0\).

\paragraph{\texorpdfstring{2. The Empty Box:
\texttt{Proxy}}{2. The Empty Box: Proxy}}\label{the-empty-box-proxy}

\emph{(One constructor, Zero computational data, Zero contextual data).}

The smallest possible Functor holds absolutely the minimum amount of
data: \textbf{none}.

\begin{Shaded}
\begin{Highlighting}[]
\KeywordTok{data} \DataTypeTok{Proxy}\NormalTok{ a }\OtherTok{=} \DataTypeTok{Proxy}
\end{Highlighting}
\end{Shaded}

The key here is that there is only one possible way to construct a
\texttt{Proxy\ a}: using the empty constructor that produces an empty
\texttt{Proxy}. It maps any phantom type \texttt{a} to a constructor
that contains zero term-level data. The type \texttt{a} exists only at
compile time; at runtime, the box is completely empty.

\textbf{Functor Implementation}:

\begin{Shaded}
\begin{Highlighting}[]
\KeywordTok{instance} \DataTypeTok{Functor} \DataTypeTok{Proxy} \KeywordTok{where}
\OtherTok{    fmap ::}\NormalTok{ (a }\OtherTok{{-}\textgreater{}}\NormalTok{ b) }\OtherTok{{-}\textgreater{}} \DataTypeTok{Proxy}\NormalTok{ a }\OtherTok{{-}\textgreater{}} \DataTypeTok{Proxy}\NormalTok{ b}
    \FunctionTok{fmap}\NormalTok{ \_ }\DataTypeTok{Proxy} \OtherTok{=} \DataTypeTok{Proxy}
\end{Highlighting}
\end{Shaded}

\textbf{The ``Why''}: Due to parametricity, there is exactly one
possible implementation that compiles. We are given a function
\texttt{(a\ -\textgreater{}\ b)}. We have a \texttt{Proxy\ a} (value
\texttt{Proxy}). We must return a \texttt{Proxy\ b} (value
\texttt{Proxy}). We have no \texttt{a} to feed into the function.
Therefore, the function \emph{must} be ignored.

\textbf{Law Verification}: * \emph{Identity}:
\texttt{fmap\ id\ Proxy\ ==\ Proxy\ ==\ id\ Proxy} * \emph{Composition}:
Guaranteed automatically by parametricity (``Theorems for free!'') since
the Identity law is satisfied.

\emph{(Note: As proven by Wadler's ``Theorems for free!'', satisfying
the Identity law automatically guarantees the Composition law for any
parametrically polymorphic functor. We explicitly verify both here and
throughout this section purely for the sake of a complete, explicit
proof).}

\textbf{Category Theory Equivalent}: This represents the constant
functor \(\Delta_1\). It maps every object in the category space to the
Terminal Object \(1\) (the singleton set \(\{*\}\)) and every morphism
to \(id_1\).

\paragraph{\texorpdfstring{3. The Constant Context:
\texttt{Const\ r}}{3. The Constant Context: Const r}}\label{the-constant-context-const-r}

\emph{(Zero computational data, Some contextual data \texttt{r}).}

If \texttt{Proxy} holds no data, \texttt{Const} holds zero
\emph{computational} data \texttt{a}, but stores an orthogonal
contextual value \texttt{r}. (We will see later in Chapter 2 that this
structure acts as an ``Accumulator'' once it is upgraded to an
\texttt{Applicative}).

\begin{Shaded}
\begin{Highlighting}[]
\KeywordTok{newtype} \DataTypeTok{Const}\NormalTok{ r a }\OtherTok{=} \DataTypeTok{Const}\NormalTok{ r}
\end{Highlighting}
\end{Shaded}

\textbf{Functor Implementation}:

\begin{Shaded}
\begin{Highlighting}[]
\KeywordTok{instance} \DataTypeTok{Functor}\NormalTok{ (}\DataTypeTok{Const}\NormalTok{ r) }\KeywordTok{where}
\OtherTok{    fmap ::}\NormalTok{ (a }\OtherTok{{-}\textgreater{}}\NormalTok{ b) }\OtherTok{{-}\textgreater{}} \DataTypeTok{Const}\NormalTok{ r a }\OtherTok{{-}\textgreater{}} \DataTypeTok{Const}\NormalTok{ r b}
    \FunctionTok{fmap}\NormalTok{ \_ (}\DataTypeTok{Const}\NormalTok{ x) }\OtherTok{=} \DataTypeTok{Const}\NormalTok{ x}
\end{Highlighting}
\end{Shaded}

\textbf{The ``Why''}: We need to create an instance of
\texttt{Const\ r\ b}. To do this, we need an instance of \texttt{r}. The
mapping function \texttt{f} cannot help us because we don't have any
\texttt{a} to feed it! So the only way is to extract the \texttt{r} from
the passed instance of \texttt{Const\ r\ a} (via \texttt{getConst} or,
as done here, simple pattern matching). There is mathematically no other
choice. Note that at the Functor level, \texttt{r} requires no special
structure (it doesn't need to be a \texttt{Monoid}).

\textbf{Law Verification}: * \emph{Identity}:
\texttt{fmap\ id\ (Const\ r)\ ==\ Const\ r\ ==\ id\ (Const\ r)} *
\emph{Composition}: Guaranteed automatically by parametricity
(``Theorems for free!'') since the Identity law is satisfied.

\textbf{Notes on Specializing \texttt{Const}:} *
\textbf{\texttt{Const\ Void\ =\ Zero}}: If we specialize \texttt{r} to
\texttt{Void} (a type with zero inhabitants, logically defined as
\texttt{data\ Void}), \texttt{Const\ Void} becomes impossible to
instantiate at runtime. Thus, \texttt{Const\ Void} is mathematically
isomorphic to our completely empty \texttt{Zero} functor. It is actually
very common in real-world Haskell to write \texttt{Const\ Void} instead
of defining a custom \texttt{Zero}! *
\textbf{\texttt{Const\ ()\ =\ Proxy}}: If we specialize \texttt{r} to
the unit type \texttt{()} (a type with exactly one inhabitant, logically
defined as \texttt{data\ ()\ =\ ()}), we get a functor that safely
exists but carries zero bits of information. Thus, \texttt{Const\ ()} is
mathematically isomorphic to our empty box \texttt{Proxy}! You can
translate back and forth between \texttt{Proxy} and \texttt{Const\ ()}
without losing any data (i.e., you can write functions
\texttt{f\ (Const\ ())\ =\ Proxy} and \texttt{g\ Proxy\ =\ Const\ ()}
where applying both functions always returns the exact original value).
* \textbf{\texttt{Const\ Bool}}: If we specialize \texttt{r} to
\texttt{Bool} (a type with exactly two inhabitants, logically defined as
\texttt{data\ Bool\ =\ False\ \textbar{}\ True}), we get a functor that
safely exists and carries exactly one bit of information (True or
False). Thus, \texttt{Const\ Bool} is mathematically isomorphic to
\texttt{Either\ (Proxy\ a)\ (Proxy\ a)}, where \texttt{Left\ Proxy} acts
as \texttt{False} and \texttt{Right\ Proxy} acts as \texttt{True}!

\emph{(We will see in Section 1.3 how these three specific
specializations intimately link to the numbers \(0\), \(1\), and \(2\)
in algebraic arithmetic!)}

\textbf{Category Theory Equivalent}: This represents the general
constant functor \(\Delta_r\). It collapses the entire category, mapping
every object to the specific fixed object \(r\), and every morphism
mathematically to the identity morphism \(id_r\).

\paragraph{\texorpdfstring{4. The Wrapper:
\texttt{Identity}}{4. The Wrapper: Identity}}\label{the-wrapper-identity}

\emph{(One computational data, Zero contextual data).}

Next is the minimal structure with exactly \emph{one} value: a
transparent wrapper.

\begin{Shaded}
\begin{Highlighting}[]
\KeywordTok{newtype} \DataTypeTok{Identity}\NormalTok{ a }\OtherTok{=} \DataTypeTok{Identity}\NormalTok{ a}
\end{Highlighting}
\end{Shaded}

\textbf{Functor Implementation}:

\begin{Shaded}
\begin{Highlighting}[]
\KeywordTok{instance} \DataTypeTok{Functor} \DataTypeTok{Identity} \KeywordTok{where}
    \FunctionTok{fmap}\NormalTok{ f (}\DataTypeTok{Identity}\NormalTok{ x) }\OtherTok{=} \DataTypeTok{Identity}\NormalTok{ (f x)}
\end{Highlighting}
\end{Shaded}

\textbf{The ``Why''}: The type signature demands we produce an
\texttt{Identity\ b}. We possess an \texttt{x\ ::\ a} and a function
\texttt{f\ ::\ a\ -\textgreater{}\ b}. The \emph{only} mathematical way
to obtain a \texttt{b} is to apply \texttt{f} to \texttt{x}.

\textbf{Law Verification}: * \emph{Identity}:
\texttt{fmap\ id\ (Identity\ x)\ ==\ Identity\ (id\ x)\ ==\ Identity\ x\ ==\ id\ (Identity\ x)}
* \emph{Composition}:
\texttt{fmap\ (f\ .\ g)\ (Identity\ x)\ ==\ Identity\ ((f\ .\ g)\ x)\ ==\ Identity\ (f\ (g\ x))\ ==\ fmap\ f\ (Identity\ (g\ x))\ ==\ fmap\ f\ (fmap\ g\ (Identity\ x))}

\textbf{Category Theory Equivalent}: This represents the Identity
Functor \(Id_{\mathbf{C}}\). It strictly maps every object to itself
(\(X \mapsto X\)) and every morphism to itself (\(f \mapsto f\)). It is
the perfectly transparent container.

\paragraph{\texorpdfstring{5. The Exponential:
\texttt{(-\textgreater{})\ r} (The
Reader)}{5. The Exponential: (-\textgreater) r (The Reader)}}\label{the-exponential---r-the-reader}

\emph{(Infinite computational data, delayed by domain \texttt{r}).}

While \texttt{Either} and \texttt{(,)} represent algebra's polynomial
addition (\(+\)) and multiplication (\(\times\)), functions represent
exponents (\(a^r\)). This forms the ``Reader'' functor: an environment
\texttt{r} waiting to produce our \texttt{a}.

\begin{Shaded}
\begin{Highlighting}[]
\CommentTok{{-}{-} The type constructor is \textasciigrave{}({-}\textgreater{}) r\textasciigrave{}. The parameter is \textasciigrave{}a\textasciigrave{}.}
\KeywordTok{instance} \DataTypeTok{Functor}\NormalTok{ ((}\OtherTok{{-}\textgreater{}}\NormalTok{) r) }\KeywordTok{where}
    \FunctionTok{fmap}\NormalTok{ f g }\OtherTok{=}\NormalTok{ f }\OperatorTok{.}\NormalTok{ g }
    \CommentTok{{-}{-} Equivalently: fmap = (.)}
\end{Highlighting}
\end{Shaded}

\textbf{The ``Why''}: We need to produce a function of type
\texttt{(r\ -\textgreater{}\ b)}. We possess a function
\texttt{g\ ::\ r\ -\textgreater{}\ a} and a mapping function
\texttt{f\ ::\ a\ -\textgreater{}\ b}. The only mathematical way to
obtain a \texttt{b} from an \texttt{r} without cheating is to pipe the
argument \texttt{r} through \texttt{g} to get an \texttt{a}, and then
pipe that \texttt{a} into \texttt{f}. This is exactly function
composition \texttt{(.)}.

\textbf{Law Verification}: * \emph{Identity}:
\texttt{fmap\ id\ g\ ==\ id\ .\ g\ ==\ g\ ==\ id\ g} *
\emph{Composition}: Guaranteed automatically by parametricity
(``Theorems for free!'') since the Identity law is satisfied.

\textbf{Category Theory Equivalent}: This represents the Covariant
\(Hom\)-functor \(Hom(r, -)\). In any category, \(Hom(A, B)\) represents
the set of all morphisms passing from object \(A\) to object \(B\). In
Haskell, fixing the input type \(r\) forms the functor mapping
\(a \mapsto Hom(r, a)\).

\subparagraph{Exponential Blends and Higher-Order
Exponentials}\label{exponential-blends-and-higher-order-exponentials}

To truly illustrate the power of parametricity, consider what happens
when we combine our building blocks (Sums, Products, and Exponentials).
Even for these complex concepts, parametricity completely forces the
only mathematically valid implementation!

\begin{itemize}
\item
  \textbf{1. The Blended Exponent/Product (\texttt{State\ s})}:
  Mathematically \((A \times S)^S\). It computes an \(a\) while
  modifying an environment \(s\).

\begin{Shaded}
\begin{Highlighting}[]
\KeywordTok{newtype} \DataTypeTok{State}\NormalTok{ s a }\OtherTok{=} \DataTypeTok{State}\NormalTok{ (s }\OtherTok{{-}\textgreater{}}\NormalTok{ (a, s))}

\KeywordTok{instance} \DataTypeTok{Functor}\NormalTok{ (}\DataTypeTok{State}\NormalTok{ s) }\KeywordTok{where}
    \FunctionTok{fmap}\NormalTok{ f (}\DataTypeTok{State}\NormalTok{ g) }\OtherTok{=} \DataTypeTok{State} \OperatorTok{$}\NormalTok{ \textbackslash{}s }\OtherTok{{-}\textgreater{}} 
        \KeywordTok{let}\NormalTok{ (a, new\_s) }\OtherTok{=}\NormalTok{ g s }
        \KeywordTok{in}\NormalTok{ (f a, new\_s)}
\end{Highlighting}
\end{Shaded}

  \emph{The ``Why''}: We must produce a function returning
  \texttt{(b,\ s)}. We possess an initial state \texttt{s} and a
  function \texttt{g} returning \texttt{(a,\ s)}. The only legal move is
  to apply \texttt{s} to \texttt{g}, extract the resulting \texttt{a},
  hit that \texttt{a} with our \texttt{f}, and return it bundled tightly
  with the new \texttt{s}! The state piping is practically written for
  us by the type system.
\item
  \textbf{2. The Higher-Order Exponential (\texttt{Cont\ r})}:
  Mathematically \(R^{(R^A)}\). It is a function that takes a callback
  \texttt{(a\ -\textgreater{}\ r)} and eventually produces an
  \texttt{r}.

\begin{Shaded}
\begin{Highlighting}[]
\KeywordTok{newtype} \DataTypeTok{Cont}\NormalTok{ r a }\OtherTok{=} \DataTypeTok{Cont}\NormalTok{ ((a }\OtherTok{{-}\textgreater{}}\NormalTok{ r) }\OtherTok{{-}\textgreater{}}\NormalTok{ r)}

\KeywordTok{instance} \DataTypeTok{Functor}\NormalTok{ (}\DataTypeTok{Cont}\NormalTok{ r) }\KeywordTok{where}
    \FunctionTok{fmap}\NormalTok{ f (}\DataTypeTok{Cont}\NormalTok{ g) }\OtherTok{=} \DataTypeTok{Cont} \OperatorTok{$}\NormalTok{ \textbackslash{}callback\_b }\OtherTok{{-}\textgreater{}} 
\NormalTok{        g (\textbackslash{}a }\OtherTok{{-}\textgreater{}}\NormalTok{ callback\_b (f a))}
\end{Highlighting}
\end{Shaded}

  \emph{The ``Why''}: This is a brain-bender, but parametricity saves
  us. We must return an \texttt{r}. We possess
  \texttt{callback\_b\ ::\ (b\ -\textgreater{}\ r)} and
  \texttt{g\ ::\ ((a\ -\textgreater{}\ r)\ -\textgreater{}\ r)}. We are
  forced to pass \emph{something} to \texttt{g} that looks like
  \texttt{(a\ -\textgreater{}\ r)}. Since we possess a
  \texttt{b\ -\textgreater{}\ r}, and an \texttt{a\ -\textgreater{}\ b},
  the only legal move is to compose them: \texttt{callback\_b\ .\ f} is
  of type \texttt{a\ -\textgreater{}\ r}. We feed that exact composition
  to \texttt{g}. The types dictate the entire callback logic!
\end{itemize}

\subsubsection{Section 1.3: Minimal
Bifunctors}\label{section-1.3-minimal-bifunctors}

Just as we started Chapter 1 by looking at the simplest possible
Functors (\texttt{Proxy}, \texttt{Const}, \texttt{Identity}), we can
apply the exact same ``shrinking'' exercise to Bifunctors
(\texttt{Type\ -\textgreater{}\ Type\ -\textgreater{}\ Type}). While
\texttt{Either} (Sum) and \texttt{(,)} (Product) are the fundamental
operations of our algebra, they both contain term-level data. We can go
simpler in three distinct ways:

\subparagraph{1. The Absolute Simplest: The ``Bi-Proxy'' (Zero
Data)}\label{the-absolute-simplest-the-bi-proxy-zero-data}

Just like \texttt{Proxy} ignoring its \texttt{a}, the simplest Bifunctor
ignores \emph{both} \texttt{a} and \texttt{b}. It is essentially an
empty box with two phantom types.

\begin{Shaded}
\begin{Highlighting}[]
\KeywordTok{data} \DataTypeTok{BiProxy}\NormalTok{ a b }\OtherTok{=} \DataTypeTok{BiProxy}
\end{Highlighting}
\end{Shaded}

\textbf{Bifunctor Implementation}:

\begin{Shaded}
\begin{Highlighting}[]
\KeywordTok{instance} \DataTypeTok{Bifunctor} \DataTypeTok{BiProxy} \KeywordTok{where}
\NormalTok{    bimap \_ \_ }\DataTypeTok{BiProxy} \OtherTok{=} \DataTypeTok{BiProxy}
\end{Highlighting}
\end{Shaded}

\textbf{The ``Why''}: The signature demands we produce a
\texttt{BiProxy\ c\ d} (value \texttt{BiProxy}). We are given two
functions \texttt{(a\ -\textgreater{}\ c)} and
\texttt{(b\ -\textgreater{}\ d)}. Because we possess neither an
\texttt{a} nor a \texttt{b} to apply the functions to, parametricity
forces us to ignore both functions entirely.

\textbf{Law Verification}: * \emph{Identity}:
\texttt{bimap\ id\ id\ BiProxy\ ==\ BiProxy\ ==\ id\ BiProxy} *
\emph{Composition}:
\texttt{bimap\ (f\ .\ g)\ (h\ .\ i)\ BiProxy\ ==\ BiProxy\ ==\ bimap\ f\ h\ BiProxy\ ==\ bimap\ f\ h\ (bimap\ g\ i\ BiProxy)}

\subparagraph{2. The Unrelated Constant (Context Data
Only)}\label{the-unrelated-constant-context-data-only}

Just like \texttt{Const\ r\ a} holds an \texttt{r} but ignores
\texttt{a}, we can have a Bifunctor that holds an \texttt{r} but ignores
both \texttt{a} and \texttt{b}. \emph{(Notice the exact same parallel
here: if we specialize \texttt{r} to the unit type \texttt{()}, we get
\texttt{ConstContext\ ()}, which is mathematically isomorphic to
\texttt{BiProxy}!)}

\begin{Shaded}
\begin{Highlighting}[]
\KeywordTok{newtype} \DataTypeTok{ConstContext}\NormalTok{ r a b }\OtherTok{=} \DataTypeTok{ConstContext}\NormalTok{ r}
\end{Highlighting}
\end{Shaded}

\textbf{Bifunctor Implementation}:

\begin{Shaded}
\begin{Highlighting}[]
\KeywordTok{instance} \DataTypeTok{Bifunctor}\NormalTok{ (}\DataTypeTok{ConstContext}\NormalTok{ r) }\KeywordTok{where}
\NormalTok{    bimap \_ \_ (}\DataTypeTok{ConstContext}\NormalTok{ r) }\OtherTok{=} \DataTypeTok{ConstContext}\NormalTok{ r}
\end{Highlighting}
\end{Shaded}

\textbf{The ``Why''}: We must produce a \texttt{ConstContext\ r\ c\ d}.
We possess an orthogonal context value \texttt{r}. Because we have no
\texttt{a} or \texttt{b} to transform, we are forced to discard the
mapping functions and return the unadulterated context.

\textbf{Law Verification}: * \emph{Identity}:
\texttt{bimap\ id\ id\ (ConstContext\ r)\ ==\ ConstContext\ r\ ==\ id\ (ConstContext\ r)}
* \emph{Composition}:
\texttt{bimap\ (f\ .\ g)\ (h\ .\ i)\ (ConstContext\ r)\ ==\ ConstContext\ r\ ==\ bimap\ f\ h\ (ConstContext\ r)\ ==\ bimap\ f\ h\ (bimap\ g\ i\ (ConstContext\ r))}

\subparagraph{3. The One-Sided Constants (Left and
Right)}\label{the-one-sided-constants-left-and-right}

A Bifunctor takes two arguments. We can define Bifunctors that act like
\texttt{Identity} on one side, and \texttt{Proxy} on the other.

\textbf{The Left identity (ignoring the right):}

\begin{Shaded}
\begin{Highlighting}[]
\KeywordTok{newtype} \DataTypeTok{ConstLeft}\NormalTok{ a b }\OtherTok{=} \DataTypeTok{ConstLeft}\NormalTok{ a}
\end{Highlighting}
\end{Shaded}

\textbf{Bifunctor Implementation}:

\begin{Shaded}
\begin{Highlighting}[]
\KeywordTok{instance} \DataTypeTok{Bifunctor} \DataTypeTok{ConstLeft} \KeywordTok{where}
\NormalTok{    bimap f \_ (}\DataTypeTok{ConstLeft}\NormalTok{ a) }\OtherTok{=} \DataTypeTok{ConstLeft}\NormalTok{ (f a)}
\end{Highlighting}
\end{Shaded}

\textbf{The ``Why''}: We need a \texttt{ConstLeft\ c\ d}. We possess an
\texttt{a} and a function \texttt{(a\ -\textgreater{}\ c)}. We are
mathematically forced to apply \texttt{f} to \texttt{a} to produce the
required \texttt{c}. Since we possess no \texttt{b}, the second function
is ignored.

\textbf{Law Verification}: * \emph{Identity}:
\texttt{bimap\ id\ id\ (ConstLeft\ a)\ ==\ ConstLeft\ (id\ a)\ ==\ ConstLeft\ a\ ==\ id\ (ConstLeft\ a)}
* \emph{Composition}:
\texttt{bimap\ (f\ .\ g)\ (h\ .\ i)\ (ConstLeft\ a)\ ==\ ConstLeft\ ((f\ .\ g)\ a)\ ==\ ConstLeft\ (f\ (g\ a))\ ==\ bimap\ f\ h\ (ConstLeft\ (g\ a))\ ==\ bimap\ f\ h\ (bimap\ g\ i\ (ConstLeft\ a))}

\textbf{The Right identity (ignoring the left):}

\begin{Shaded}
\begin{Highlighting}[]
\KeywordTok{newtype} \DataTypeTok{ConstRight}\NormalTok{ a b }\OtherTok{=} \DataTypeTok{ConstRight}\NormalTok{ b}
\end{Highlighting}
\end{Shaded}

\textbf{Bifunctor Implementation}:

\begin{Shaded}
\begin{Highlighting}[]
\KeywordTok{instance} \DataTypeTok{Bifunctor} \DataTypeTok{ConstRight} \KeywordTok{where}
\NormalTok{    bimap \_ g (}\DataTypeTok{ConstRight}\NormalTok{ b) }\OtherTok{=} \DataTypeTok{ConstRight}\NormalTok{ (g b)}
\end{Highlighting}
\end{Shaded}

\textbf{The ``Why''}: We need a \texttt{ConstRight\ c\ d}. We possess a
\texttt{b} and a function \texttt{(b\ -\textgreater{}\ d)}.
Parametricity dictates we must apply \texttt{g} to \texttt{b} to produce
the required \texttt{d}. The first function is ignored.

\textbf{Law Verification}: * \emph{Identity}:
\texttt{bimap\ id\ id\ (ConstRight\ b)\ ==\ ConstRight\ (id\ b)\ ==\ ConstRight\ b\ ==\ id\ (ConstRight\ b)}
* \emph{Composition}:
\texttt{bimap\ (f\ .\ g)\ (h\ .\ i)\ (ConstRight\ b)\ ==\ ConstRight\ ((h\ .\ i)\ b)\ ==\ ConstRight\ (h\ (i\ b))\ ==\ bimap\ f\ h\ (ConstRight\ (i\ b))\ ==\ bimap\ f\ h\ (bimap\ g\ i\ (ConstRight\ b))}

\subparagraph{\texorpdfstring{4. The Sum Molecule:
\texttt{Either}}{4. The Sum Molecule: Either}}\label{the-sum-molecule-either}

The fundamental co-product of two types.

\textbf{Bifunctor Implementation}:

\begin{Shaded}
\begin{Highlighting}[]
\KeywordTok{instance} \DataTypeTok{Bifunctor} \DataTypeTok{Either} \KeywordTok{where}
\NormalTok{    bimap f \_ (}\DataTypeTok{Left}\NormalTok{ a)  }\OtherTok{=} \DataTypeTok{Left}\NormalTok{ (f a)}
\NormalTok{    bimap \_ g (}\DataTypeTok{Right}\NormalTok{ b) }\OtherTok{=} \DataTypeTok{Right}\NormalTok{ (g b)}
\end{Highlighting}
\end{Shaded}

\textbf{The ``Why''}: \texttt{Either} encapsulates a choice. If the
constructor contains an \texttt{a} (\texttt{Left}), we are forced to
apply \texttt{f} to obtain a \texttt{c}. If it contains a \texttt{b}
(\texttt{Right}), we are forced to apply \texttt{g} to obtain a
\texttt{d}.

\textbf{Law Verification} (The Developer's Responsibility!):

It is crucial to remember that the Haskell compiler \textbf{only checks
types, not math}. It will perfectly compile a \texttt{Bifunctor}
instance as long as the type signatures align, even if it completely
violates the Identity and Composition laws! You, the developer, are
solely responsible for ensuring your instance mathematically preserves
the shape of your data.

While we can easily prove these properties mathematically by hand for
simple types (as shown below), in Haskell we can actually automate this
verification! Using property testing libraries like
\texttt{tasty-quickcheck} (and typeclass rule validators like
\texttt{quickcheck-classes}), we can generate thousands of random
instances to guarantee our Bifunctor truly behaves correctly.

A test suite verifying \texttt{Either} can be reduced to one simple
line:

\begin{Shaded}
\begin{Highlighting}[]
\KeywordTok{import} \DataTypeTok{Test.QuickCheck.Classes}

\CommentTok{{-}{-} Automatically tests both Identity and Composition!}
\NormalTok{testProperties }\StringTok{"Either Bifunctor"} \OperatorTok{$}\NormalTok{ bifunctor (}\DataTypeTok{Proxy}\OtherTok{ ::} \DataTypeTok{Proxy} \DataTypeTok{Either}\NormalTok{)}
\end{Highlighting}
\end{Shaded}

This ensures we never break the two fundamental rules: *
\emph{Identity}:
\texttt{haskell\ \ \ \ \ bimap\ id\ id\ (Left\ a)\ ==\ Left\ (id\ a)\ ==\ Left\ a\ ==\ id\ (Left\ a)\ \ \ \ \ bimap\ id\ id\ (Right\ b)\ ==\ Right\ (id\ b)\ ==\ Right\ b\ ==\ id\ (Right\ b)}
* \emph{Composition}:
\texttt{haskell\ \ \ \ \ bimap\ (f\ .\ g)\ (h\ .\ i)\ (Left\ a)\ ==\ Left\ ((f\ .\ g)\ a)\ ==\ Left\ (f\ (g\ a))\ ==\ bimap\ f\ h\ (Left\ (g\ a))\ ==\ bimap\ f\ h\ (bimap\ g\ i\ (Left\ a))\ \ \ \ \ bimap\ (f\ .\ g)\ (h\ .\ i)\ (Right\ b)\ ==\ Right\ ((h\ .\ i)\ b)\ ==\ Right\ (h\ (i\ b))\ ==\ bimap\ f\ h\ (Right\ (i\ b))\ ==\ bimap\ f\ h\ (bimap\ g\ i\ (Right\ b))}

\subparagraph{\texorpdfstring{5. The Product Molecule:
\texttt{(,)}}{5. The Product Molecule: (,)}}\label{the-product-molecule}

The fundamental product of two types.

\textbf{Bifunctor Implementation}:

\begin{Shaded}
\begin{Highlighting}[]
\KeywordTok{instance} \DataTypeTok{Bifunctor}\NormalTok{ (,) }\KeywordTok{where}
\NormalTok{    bimap f g (a, b) }\OtherTok{=}\NormalTok{ (f a, g b)}
\end{Highlighting}
\end{Shaded}

\textbf{The ``Why''}: A Tuple constructor definitively contains both an
\texttt{a} \emph{and} a \texttt{b}. To produce a tuple of type
\texttt{(c,\ d)}, we must apply \texttt{f} to the left element and
\texttt{g} to the right element.

\textbf{Law Verification}: * \emph{Identity}:
\texttt{bimap\ id\ id\ (a,\ b)\ ==\ (id\ a,\ id\ b)\ ==\ (a,\ b)\ ==\ id\ (a,\ b)}
* \emph{Composition}:
\texttt{bimap\ (f\ .\ g)\ (h\ .\ i)\ (a,\ b)\ ==\ ((f\ .\ g)\ a,\ (h\ .\ i)\ b)\ ==\ (f\ (g\ a),\ h\ (i\ b))\ ==\ bimap\ f\ h\ (g\ a,\ i\ b)\ ==\ bimap\ f\ h\ (bimap\ g\ i\ (a,\ b))}

\subparagraph{\texorpdfstring{6. The Dual Exponential:
\texttt{BiReader\ r}}{6. The Dual Exponential: BiReader r}}\label{the-dual-exponential-bireader-r}

Just as we saw functions pull us out of polynomial algebras at the 1D
Functor level, an exponential delays computation at the 2D Bifunctor
level. Mathematically, it is \((A \times B)^R\).

\textbf{Bifunctor Implementation}:

\begin{Shaded}
\begin{Highlighting}[]
\KeywordTok{newtype} \DataTypeTok{BiReader}\NormalTok{ r a b }\OtherTok{=} \DataTypeTok{BiReader}\NormalTok{ (r }\OtherTok{{-}\textgreater{}}\NormalTok{ (a, b))}

\KeywordTok{instance} \DataTypeTok{Bifunctor}\NormalTok{ (}\DataTypeTok{BiReader}\NormalTok{ r) }\KeywordTok{where}
\NormalTok{    bimap f g (}\DataTypeTok{BiReader}\NormalTok{ h) }\OtherTok{=} \DataTypeTok{BiReader} \OperatorTok{$}\NormalTok{ \textbackslash{}r }\OtherTok{{-}\textgreater{}} 
        \KeywordTok{let}\NormalTok{ (a, b) }\OtherTok{=}\NormalTok{ h r }
        \KeywordTok{in}\NormalTok{ (f a, g b)}
\end{Highlighting}
\end{Shaded}

\textbf{The ``Why''}: We are returning a delayed computation of a tuple.
We possess a function \texttt{h\ ::\ r\ -\textgreater{}\ (a,\ b)}. We
are given two mapping functions \texttt{f\ ::\ a\ -\textgreater{}\ c}
and \texttt{g\ ::\ b\ -\textgreater{}\ d}. The only legal mathematical
move is to intercept the environment \texttt{r} the moment it arrives,
feed it to \texttt{h} to obtain our \texttt{a} and \texttt{b}, apply
\texttt{f} to \texttt{a}, apply \texttt{g} to \texttt{b}, and return the
newly bundled tuple. The entire pipeline is rigidly defined by the types
involved.

\textbf{Law Verification}: * \emph{Identity}:
\texttt{haskell\ \ \ \ \ bimap\ id\ id\ (BiReader\ h)\ \ \ \ \ \ ==\ BiReader\ (\textbackslash{}r\ -\textgreater{}\ let\ (a,\ b)\ =\ h\ r\ in\ (id\ a,\ id\ b))\ \ \ \ \ ==\ BiReader\ (\textbackslash{}r\ -\textgreater{}\ h\ r)\ \ \ \ \ ==\ BiReader\ h}
* \emph{Composition}:
\texttt{haskell\ \ \ \ \ bimap\ (f\ .\ j)\ (g\ .\ k)\ (BiReader\ h)\ \ \ \ \ ==\ BiReader\ (\textbackslash{}r\ -\textgreater{}\ let\ (a,\ b)\ =\ h\ r\ in\ ((f\ .\ j)\ a,\ (g\ .\ k)\ b))\ \ \ \ \ ==\ BiReader\ (\textbackslash{}r\ -\textgreater{}\ let\ (a,\ b)\ =\ h\ r\ in\ (f\ (j\ a),\ g\ (k\ b)))\ \ \ \ \ -\/-\ Which\ is\ equivalent\ to:\ \ \ \ \ ==\ bimap\ f\ g\ (BiReader\ (\textbackslash{}r\ -\textgreater{}\ let\ (a,\ b)\ =\ h\ r\ in\ (j\ a,\ k\ b)))\ \ \ \ \ ==\ bimap\ f\ g\ (bimap\ j\ k\ (BiReader\ h))}

\subsubsection{Section 1.4: Bifunctors as Binary Operations on
Functors}\label{section-1.4-bifunctors-as-binary-operations-on-functors}

Because a Bifunctor maps two types into a new type, we can think of it
mathematically as a \textbf{binary operator} on the category of
Functors! By taking two existing Functors, \(F\) and \(G\), and
combining them using a Bifunctor operator \(B\), we generate an entirely
new Functor: \(H(x) = B(F(x), G(x))\).

Let's explore this using our minimal atomic functors (\texttt{Zero} and
\texttt{Proxy}) and our fundamental binary operators: Sum
(\texttt{Either} or \(+\)) and Product (\texttt{(,)} or \(\times\)). By
interacting them, we see the algebra mirror elementary arithmetic
perfectly:

\paragraph{1. Zero + Proxy = Proxy}\label{zero-proxy-proxy}

\textbf{Math}: \(0 + 1 = 1\). \textbf{Haskell}:
\texttt{Either\ (Zero\ a)\ (Proxy\ a)}. Since \texttt{Zero} is
mathematically uninhabited, it is impossible to construct the
\texttt{Left} side of the \texttt{Either}. Therefore, the only possible
inhabited value of this structure is \texttt{Right\ Proxy}. Because
there is exactly 1 state, it holds zero computational data and precisely
zero \emph{bits} of contextual data. It is perfectly isomorphic to
\texttt{Proxy}.

\paragraph{2. Zero * Proxy = Zero}\label{zero-proxy-zero}

\textbf{Math}: \(0 \times 1 = 0\). \textbf{Haskell}:
\texttt{(Zero\ a,\ Proxy\ a)}. To construct a tuple, you MUST provide
both the left and right sides. Because we can never construct a
\texttt{Zero}, it becomes impossible to \emph{ever} construct the tuple
as a whole. The type is uninhabited, making it perfectly isomorphic to
\texttt{Zero}.

\paragraph{3. Proxy + Proxy = Const Bool}\label{proxy-proxy-const-bool}

\textbf{Math}: \(1 + 1 = 2\). \textbf{Haskell}:
\texttt{Either\ (Proxy\ a)\ (Proxy\ a)}. Since \texttt{Proxy} on both
sides is an empty box, this structure holds absolutely no computational
data \texttt{a}. However, it \emph{does} hold exactly 1 bit of
information: whether it is the \texttt{Left} empty box or the
\texttt{Right} empty box! Because a Bool has exactly 2 states
(True/False), this structure is isomorphic to \texttt{Const\ Bool\ a}.
\(1 + 1\) successfully yielded \(2\)!

\paragraph{4. Proxy * Proxy = Proxy}\label{proxy-proxy-proxy}

\textbf{Math}: \(1 \times 1 = 1\). \textbf{Haskell}:
\texttt{(Proxy\ a,\ Proxy\ a)}. We must provide an empty box for the
left side and an empty box for the right side. The state
\texttt{(Proxy,\ Proxy)} is the \emph{only} possible state this
structure can ever be in. Since it has only one state, it yields zero
bits of contextual information and holds zero data, bringing us right
back to 1. It is isomorphic to \texttt{Proxy}.

\paragraph{5. Proxy * Identity =
Identity}\label{proxy-identity-identity}

\textbf{Math}: \(1 \times X = X\). \textbf{Haskell}:
\texttt{(Proxy\ a,\ Identity\ a)}. A tuple containing an empty box and a
single \texttt{a}. The left side adds no data and has no alternative
states. The entire structure simply holds precisely one \texttt{a},
making it perfectly isomorphic to \texttt{Identity\ a}.

\emph{(Notice that all the examples in this section were specifically
chosen to demonstrate mathematical relations between the exact minimal
Functors we have already defined. We haven't built any ``new'' ADTs
yet!)}

\paragraph{6. Constant Functors as The
Ordinals}\label{constant-functors-as-the-ordinals}

Now that we have seen how \texttt{+} and \texttt{\textbackslash{}times}
interact with our minimal atoms, we can finally understand a profound
property of the \texttt{Const\ r\ a} functor. By changing the embedded
type \texttt{r}, \texttt{Const} mathematically represents the discrete
numbers (Ordinals) based solely on the number of inhabited states of
\texttt{r}: * \textbf{\(0\)}: \texttt{Const\ Void} (zero inhabitants,
isomorphic to \texttt{Zero}) * \textbf{\(1\)}: \texttt{Const\ ()} (one
inhabitant, isomorphic to \texttt{Proxy}) * \textbf{\(2\)}:
\texttt{Const\ Bool} (two inhabitants, exactly as derived by \(1 + 1\))
* \textbf{\(3\)}: \texttt{Const\ Ordering} (three inhabitants:
\texttt{LT}, \texttt{EQ}, \texttt{GT}, exactly matching \(1 + 1 + 1\)) *
\textbf{\(4\)}: \texttt{Const\ (Bool,\ Bool)} (four inhabitants, exactly
matching \(2 \times 2\)) * \textbf{\(5\)}:
\texttt{Const\ (Either\ Bool\ Ordering)} (five inhabitants, exactly
matching \(2 + 3\)) * \textbf{\(6\)}: \texttt{Const\ (Bool,\ Ordering)}
(six inhabitants, exactly matching \(2 \times 3\)) * \textbf{\(7\)}:
\texttt{Const\ (Either\ (Bool,\ Bool)\ Ordering)} (seven inhabitants,
exactly matching \(4 + 3\)) * \textbf{\(n\)}: Any \texttt{Const\ r}
where \texttt{r} is a finite enum with \(n\) states\ldots{}

This conceptually proves why \texttt{Const\ Void} acts as the true
algebraic identity for Sum (\(0\)), and \texttt{Const\ ()} acts as the
true algebraic identity for Product (\(1\)) when subjected to actual
Bifunctor addition and multiplication!

By treating Bifunctors as binary operators running on simple atomic
Functors, we observe the foundation of Algebraic Data Types emerging
exactly like fundamental school arithmetic.

\subsubsection{Section 1.5: Deriving the Atoms from
Bifunctors}\label{section-1.5-deriving-the-atoms-from-bifunctors}

In mathematical systems, we often don't just invent the ``atomic''
elements out of thin air. We derive them from the operations themselves.
Here, we are deeply interested in extracting ``natural'' atomic Functors
directly out of our foundational Bifunctors.

\textbf{The Big Picture}: Our grand architectural goal is to select a
minimal set of fundamental Bifunctor binary operations (like \texttt{+}
and \texttt{*}). From this selected set of Bifunctors, we want to
``naturally'' extract simple, atomic Functors (like \(0\) and \(1\)).
Once we have derived these foundational atoms, we can combine them
iteratively with our Bifunctors to form their mathematical
\emph{closure}. This exact generative process---using Bifunctor
operations to compose simple extracted atoms---is the traditional
mathematical mechanism for defining entire sub-categories of Functors.
This is exactly how we generate the infinitely rich families of everyday
Algebraic Data Types we use in programming! In particular, we are deeply
interested in extracting atomic Functors that perfectly preserve
\emph{parametricity}. By doing so, the type system strictly forces our
hand to yield a single, mathematically unique, ``correct by
construction'' implementation for each structure---a profound
intellectual economy that we will explore below.

Let's break down exactly how this natural extraction works. \#\#\#\# 1.
Extracting a Functor from a Bifunctor

How do we extract a standard Functor out of a generic Bifunctor?
Technically, we can \emph{always} extract a Functor simply by fixing one
of the two type arguments to an arbitrary type \(T\) (so
\(F(A) = B(T, A)\)). This is mathematically just partial application!

For example: * Instead of \(A + B\), we fix the left side to
\texttt{String}: \texttt{Either\ String\ a}. This yields a Functor
representing a computation that either succeeds with an \texttt{a} or
fails with a \texttt{String} error. * Instead of \(A \times B\), we fix
the left side to \texttt{Int}: \texttt{(Int,\ a)}. This yields a Functor
that simply packages an arbitrary integer alongside an \texttt{a}.

\emph{(Note: Because a true Bifunctor is mathematically covariant in
both arguments, fixing either the left side \(B(T, A)\) or the right
side \(B(A, T)\) yields a perfectly valid Functor! However, in Haskell,
type lambdas are partially applied left-to-right, making fixing the left
side the native default syntax).}

However, making a random, arbitrary choice of \(T\) (like picking
\texttt{String} or \texttt{Int} out of millions of possible types) is
not a ``natural'' mathematical progression. When you arbitrarily choose
a type \(T\) to partially apply, you are making an ad-hoc, manual
decision. There are infinite possible choices, and none of them are
mathematically ``more correct'' than the others.

Crucially, \textbf{this breaks parametricity if we try to extract the
inner data!} Because \texttt{String} contains actual data, we cannot
write a parametrically polymorphic, total function to extract \texttt{a}
from \texttt{Either\ String\ a} without either handling the string
(which requires specific knowledge of \texttt{String}) or crashing. We
lose the ability to generically and losslessly map our structure.

For a completely generic Bifunctor with no special algebraic properties,
making an arbitrary choice like this might be the only way to extract a
Functor.

\paragraph{2. Bifunctors with Identity
(``Naturality'')}\label{bifunctors-with-identity-naturality}

But if the Bifunctor has a special structural property---such as
possessing a left and/or right identity element---then it is better to
find a more natural way to extract a Functor!

At its absolute bare minimum, we just need a \textbf{left identity} or a
\textbf{right identity}. What does this actually mean mathematically? It
means there must exist a specific type \(I\) along with a perfect
two-way mapping---a structural isomorphism---that proves combining \(I\)
with any type \(A\) leaves \(A\) completely unchanged (neither losing
nor inventing any data): * \textbf{A Left Identity} requires a
structural isomorphism known as the \textbf{Left Unitor} (often denoted
\(\lambda\)): proving \(B(I, A) \cong A\). * \textbf{A Right Identity}
requires a structural isomorphism known as the \textbf{Right Unitor}
(often denoted \(\rho\)): proving \(B(A, I) \cong A\).

\textbf{Crucial Distinction}: Do not confuse these properties with the
\texttt{Bifunctor} laws! The Functor/Bifunctor laws (Identity and
Composition) govern the \emph{behavior of mapping functions} and must
hold via \textbf{strict equality} (e.g., \texttt{fmap\ id\ ==\ id}). In
contrast, possessing a Left or Right Identity type is a property of the
\emph{data structure itself}, proven via \textbf{structural isomorphism}
(\(\cong\), meaning the shapes can losslessly map to each other even if
they aren't strictly identical types).

\emph{(Note: If a binary operation has both, math dictates they must be
the identical type \(I\). See the
\hyperref[proof-of-identity-uniqueness]{Annex: Proof of Identity
Uniqueness} for the derivation!)}

When you use the identity \(I\) to perform your partial application, the
choice is no longer yours---the inherent structure of the Bifunctor
\emph{forces} its own unique canonical choice onto you! That uniqueness
is exactly what ``naturality'' refers to in this context: it arises
purely from the structure itself, independent of arbitrary external
choices.

This \textbf{``forced hand''} is exactly what we are aiming for. In
functional programming, we are deeply interested in this kind of
\textbf{intellectual economy}: we want to identify and produce
foundational Functors that have exactly \emph{one} mathematically unique
implementation. By relying on naturality, we eliminate arbitrary
decisions and derive primitive structures that are completely ``correct
by construction.''

By taking that uniquely canonical identity \(I\) and turning it into a
constant mapping, we establish the fundamental ``Atomic'' Functor for
that operation naturally. We create a Constant Functor \(C(A) = I\).

\emph{(Technical Note: In Haskell, a ``Natural Transformation'' between
two Functors \texttt{f} and \texttt{g} is exactly the type signature
\texttt{forall\ a.\ f\ a\ -\textgreater{}\ g\ a}. When we say
\(B(I, A)\) naturally resolves to \(A\), it means we can write a perfect
Natural Transformation mapping without losing or inventing data. For
example, for the Product \texttt{(,)} with identity \texttt{()}, the
natural transformation to \texttt{Identity} is literally just
\texttt{snd\ ::\ forall\ a.\ ((),\ a)\ -\textgreater{}\ a}! For Sum
\texttt{Either} with identity \texttt{Void}, it is
\texttt{extract\ ::\ forall\ a.\ Either\ Void\ a\ -\textgreater{}\ a}
via absurd. It is a mathematical guarantee encoded seamlessly into the
language.)}

Let's classify the ``zoo'' of Bifunctors we have seen so far based on
this profound property: * \textbf{No Identity}: Bifunctors like
\texttt{BiProxy} or \texttt{ConstContext} have neither a left nor a
right identity. To extract a Functor from them, you are forced to make
an arbitrary, non-natural choice! \emph{(Why? Because if a left identity
\texttt{I} existed, then \texttt{BiProxy\ I\ Bool} must be perfectly
isomorphic to \texttt{Bool}. But \texttt{BiProxy} always has exactly 1
inhabitant, which makes it mathematically impossible to form a two-way
mapping with \texttt{Bool}'s 2 inhabitants!)} * \textbf{Left Identity
Only}: The function arrow \texttt{(-\textgreater{})} is a profound
binary operation. It only possesses a left identity \texttt{()} (since
\texttt{()\ -\textgreater{}\ a} is isomorphic to exactly one \texttt{a},
but \texttt{a\ -\textgreater{}\ ()} is not \texttt{a}). * \textbf{Full
Identity}: * The \textbf{Sum Bifunctor} (\texttt{Either} or \(+\)) has
the two-sided mathematical identity \(0\) (the \texttt{Void} type, since
\(A + 0 \cong A\)). From this, we gracefully extract the constant
functor \texttt{Const\ Void} (or \texttt{Zero}). * The \textbf{Product
Bifunctor} (\texttt{(,)} or \(\times\)) has the two-sided mathematical
identity \(1\) (the \texttt{()} type, since \(A \times 1 \cong A\)).
From this, we extract the constant functor \texttt{Const\ ()} (or
\texttt{Proxy}).

\emph{(Note: This means mathematically, \texttt{Proxy} is not truly the
``simplest''---it is simply \(1\). \texttt{Const\ Void} is strictly
smaller as it is exactly \(0\)!)}

\paragraph{3. The Power of ``Families'' (Sub-Category
Closures)}\label{the-power-of-families-sub-category-closures}

What happens if we iteratively apply a Bifunctor and its identity? By
definition, if we only take a single Bifunctor (like \(\times\)) and its
identity (\(1\)), the mathematical closure is fairly trivial. We can
only generate structures like \(1\), \(1 \times 1\), \(1 \times A\),
\(A \times A\), etc. This forms a flat lineage (just tuples of identical
shape or empty structures). If we just take the closure of the identity
itself with \(A\), we trivially just get the Identity functor.

\paragraph{4. The Magic of Polynomial
Functors}\label{the-magic-of-polynomial-functors}

However, things get deeply interesting when we take a \emph{set} of two
orthogonal interacting Bifunctors---like \(+\) and \(\times\)---and
their respective identities. By mixing Sums, Products, Zeros, and Ones,
we generate an infinitely rich family of structures. This exact closure
is the \textbf{Category of Polynomial Functors} (e.g.,
\(1 + A + A \times A...\)). This interplay is what allows us to define
lists, trees, and essentially every Algebraic Data Type (ADT) in
programming.

\paragraph{5. Is an Identity strictly
required?}\label{is-an-identity-strictly-required}

Must every Bifunctor in our set have an identity? Not necessarily! It is
mathematically perfectly valid to consider a set of Bifunctors where
only some (or none) have identities (this essentially forms a non-unital
algebraic structure).

But does this restricted set generate an \emph{interesting} subcategory
of functors? Absolutely! Let's say we have our two fundamental
bifunctors (\(+\) and \(\times\)). Let's assume we possess the Sum
Identity \(0\) (the \texttt{Void} type) but we \textbf{do not possess}
the Product Identity \(1\) (the \texttt{()} type/\texttt{Proxy}).

By missing \(1\), we can never create a ``Nil'' or an ``Empty''
constructor to terminate our recursive shapes. As a profound result, the
closure of our variables with merely \(\{+, \times, 0\}\) mathematically
generates the incredibly restrictive \emph{Subcategory of Non-Empty Data
Structures}: * \textbf{The Non-Empty List}:
\(NEL(A) = A + A \times NEL(A)\). (Haskell's
\texttt{Data.List.NonEmpty}). * \textbf{The Un-emptyable Tree}:
\(Tree(A) = A + Tree(A) \times Tree(A)\). (A tree where every leaf must
have a value).

This subcategory guarantees---at the compiler level---that every single
structure geometrically contains at least one \(A\). The absence of the
mathematical \(1\) identity is exactly what powers this profound
property!

However, to form the full ``Polynomial'' category that exactly matches
the power of general computer science ADTs, \emph{both} of our
fundamental operations (\(+\) and \(\times\)) require their natural
identities (\(0\) and \(1\)) to terminate data structures (like using
\(1\) as the empty \texttt{Nil} constructor ending a \texttt{List}).

\paragraph{6. Examples of Deriving
Compounds}\label{examples-of-deriving-compounds}

By leveraging combinations of our extracted identities (\texttt{Zero},
\texttt{Proxy}) and fundamental functors (\texttt{Identity}), we
systematically generate powerful structures using the Bifunctor
operations. * \textbf{Optional Data}: \(1 + X\). Using Sum:
\texttt{Either\ (Proxy\ a)\ (Identity\ a)} is isomorphic to
\texttt{Maybe\ a}. * \textbf{Error Context}: \(E + X\).
\texttt{Either\ (Const\ e\ a)\ (Identity\ a)} gives us a computation
that succeeds with an \texttt{a} or fails with an error \texttt{e}. *
\textbf{Logging Context}: \(E \times X\).
\texttt{(Const\ e\ a,\ Identity\ a)} perfectly mirrors a \texttt{Writer}
log context bundled with an \texttt{a}.

\paragraph{The Ultimate Closure: Bicartesian Closed Categories
(BCC)}\label{the-ultimate-closure-bicartesian-closed-categories-bcc}

So, we have established our two algebraic bifunctors (Sum and Product)
and derived their natural identity atoms (\(0\) and \(1\)). What happens
if we take exactly these, and add our third non-algebraic bifunctor: the
\textbf{Exponential} (\texttt{-\textgreater{}})?

If a category contains exactly those three foundational Bifunctor
operations (\texttt{Either}, \texttt{(,)}, and \texttt{-\textgreater{}})
along with their identities (\texttt{Void} and \texttt{()}), it fulfills
the mathematical requirements to be called a \textbf{Bicartesian Closed
Category} (BCC).

\begin{itemize}
\tightlist
\item
  \textbf{``Cartesian''}: The category possesses Products (\(\times\))
  and a Terminal Object (\(1\)).
\item
  \textbf{``Bi-''}: The category \emph{also} possesses Coproducts
  (\(+\)) and an Initial Object (\(0\)).
\item
  \textbf{``Closed''}: The category possesses Exponentials
  (\texttt{-\textgreater{}}), allowing functions to be treated as values
  and evaluated.
\end{itemize}

This completely ``closed'' loop of operations is extraordinarily
profound. According to the Curry-Howard isomorphism, a Bicartesian
Closed Category is the exact mathematical equivalent of \textbf{Simply
Typed Lambda Calculus}, the theoretical foundation of intuitionistic
propositional logic.

The closure built by these three simple Bifunctors creates the entire
logical framework that strongly typed programming languages like Haskell
rely on!

\subsubsection{Section 1.6: Generating Functor Subcategories (The
Algebra as a Special
Case)}\label{section-1.6-generating-functor-subcategories-the-algebra-as-a-special-case}

\emph{(Note on Terminology: When mathematicians or Haskell programmers
say a structure is ``algebraic'' --- as in Algebraic Data Types or ADTs
--- they mean it is constructed strictly using only polynomial
combinations: Sums \texttt{+} and Products \texttt{*}. Function arrows
\texttt{-\textgreater{}} represent Exponentials, which are conceptually
a tier ``above'' simple algebra! To make this concrete: }
\textbf{Algebraic}: Things defined exclusively by values and their
geometry. This includes types like \texttt{Bool} (\(1 + 1\)),
\texttt{Maybe} (\(1 + X\)), \texttt{List}, and \texttt{Tree}, as well as
mathematical structures like \textbf{Monoids} and \textbf{Groups}. *
\textbf{Non-Algebraic (Exponentials)}: Things that require an execution
environment or delayed computation (\texttt{-\textgreater{}}). This
includes types like the \texttt{Reader} (\(A^R\)), \texttt{State}, and
\texttt{Cont}, which are structurally higher-order).*

\paragraph{1. The Algebra of Functors}\label{the-algebra-of-functors}

When you build an algebraic equation in mathematics, like
\(f(x) = 2x + 1\), you only need two foundational components to start
building: your numbers (constants like 1, 2) and your variable (\(x\)).

For standard Endofunctors (\texttt{Type\ -\textgreater{}\ Type}), it is
incredibly obvious what our two ``atomic'' building blocks must
therefore be: 1. \textbf{The Constants (\(C\))}: \texttt{Const\ r}
represents any constant value independent of \texttt{x}. At its absolute
simplest scale, \texttt{Proxy} (or \texttt{Const\ ()}) represents the
mathematical constant \(1\). 2. \textbf{The Single Variable (\(X\))}:
\texttt{Identity} rigidly represents the single parameter/variable \(x\)
itself.

Every other single-variable algebraic data type in Haskell can be built
by taking these primitives, \textbf{summing} them (using Alternative
constructors, representing \(+\)), and \textbf{multiplying} them (using
Multiple fields, representing \(\times\))!

But are Sums and Products Functors themselves? Yes! In Category Theory,
operations like Sum (\(+\)) and Product (\(\times\)) are specifically
known as \textbf{Bifunctors} because they map \emph{two} categories (or
a product of categories) into one. In Haskell, these are represented by
\texttt{Either} (Sum) and \texttt{(,)} (Product).

Because they are Bifunctors, if you fix one of their arguments, they
immediately become standard Endofunctors
(\texttt{Type\ -\textgreater{}\ Type}). Furthermore, the category of
Functors is closed over these operations: the sum or product of two
Functors is inherently a Functor (like \texttt{Data.Functor.Sum} and
\texttt{Data.Functor.Product}).

\emph{(Note: The formal laws governing how these products and sums
associate and interact are a bit more complex, requiring them to verify
the \textbf{pentagon} and \textbf{triangle} laws from Monoidal
Categories. We will refer to the details of these laws in
\hyperref[chapter-5-monoidal-categories]{Chapter 5} at the end of this
journey).}

\textbf{Functors entirely out of Proxy:} To see these Bifunctors in
action with our simplest atomic functor, \texttt{Proxy}: * \textbf{Proxy
+ Proxy = Const Bool}: Summing two Proxies creates two possible empty
states. \texttt{Either\ ()\ ()} is isomorphic to a Boolean.
Mathematically: \(1 + 1 = 2\). * \textbf{Proxy * Proxy = Proxy}: A
product of two empty boxes remains an empty box. Mathematically:
\(1 \times 1 = 1\).

\paragraph{2. The Algebra of
Bifunctors}\label{the-algebra-of-bifunctors}

Is there an algebra for Bifunctors just as there is for standard
Functors? Absolutely! Because the category of Functors is closed over
Products and Sums, we can combine our foundational Bifunctor atoms
exactly the same way to build incredibly complex Bifunctors.

If standard Functors (\texttt{Type\ -\textgreater{}\ Type}) are
single-variable polynomials like \(f(x) = x^2 + 1\), then Bifunctors
(\texttt{Type\ -\textgreater{}\ Type\ -\textgreater{}\ Type}) are simply
two-variable polynomials like \(f(a, b) = a \times b + a\).

This means it becomes very obvious what our two ``atomic variables''
are: * \textbf{The First Variable (\(A\))}:
\texttt{ConstLeft\ a\ b\ =\ ConstLeft\ a} (ignoring the right). *
\textbf{The Second Variable (\(B\))}:
\texttt{ConstRight\ a\ b\ =\ ConstRight\ b} (ignoring the left).

Equipped with our two atomic variables, we can perform any algebraic
operation: * \textbf{Bifunctor Sums (\(+\))}: We can wrap a Bifunctor
inside \texttt{Either} (e.g.~\texttt{Either\ (BiProxy\ a\ b)\ (a,\ b)}).
* \textbf{Bifunctor Products (\(\times\))}: We can tuple Bifunctors
together (e.g.~\texttt{(Either\ a\ b,\ ConstContext\ String\ a\ b)}). *
\textbf{Bifunctor Fixed Points}: Just like \texttt{List} recursively
nests standard Functors, structures like a \texttt{Bifunctor\ Tree} can
recursively nest Bifunctors
(e.g.~\texttt{data\ BiTree\ a\ b\ =\ Leaf\ a\ b\ \textbar{}\ Node\ (BiTree\ a\ b)\ (BiTree\ a\ b)}).

Anything you can do in one dimension
(\texttt{Type\ -\textgreater{}\ Type}), Category Theory allows you to
transparently extend into two dimensions
(\texttt{Type\ -\textgreater{}\ Type\ -\textgreater{}\ Type}) using the
exact same polynomial algebra!

\paragraph{\texorpdfstring{3. Composing Functors into a Bifunctor
(\texttt{Biff})}{3. Composing Functors into a Bifunctor (Biff)}}\label{composing-functors-into-a-bifunctor-biff}

While \texttt{Compose} elegantly handles nesting a Functor inside
another Functor (\texttt{f\ ∘\ g}), what happens when we want to compose
Functors directly into the independent branches of a \textbf{Bifunctor}?

Because a standard Bifunctor \texttt{p} takes exactly two type
arguments, we can mathematically substitute two independent Functors
(\texttt{f} and \texttt{g}) into those dimensional parameters! In
Haskell, this exact compositional bridge is completely formalized by the
\texttt{Biff} operator in \texttt{Data.Bifunctor.Biff}:

\begin{Shaded}
\begin{Highlighting}[]
\CommentTok{{-}{-} \textquotesingle{}p\textquotesingle{} is a Bifunctor (like Either or Pair)}
\CommentTok{{-}{-} \textquotesingle{}f\textquotesingle{} and \textquotesingle{}g\textquotesingle{} are Functors (like List, Maybe)}
\KeywordTok{newtype} \DataTypeTok{Biff}\NormalTok{ p f g a b }\OtherTok{=} \DataTypeTok{Biff}\NormalTok{ (p (f a) (g b))}

\KeywordTok{instance}\NormalTok{ (}\DataTypeTok{Bifunctor}\NormalTok{ p, }\DataTypeTok{Functor}\NormalTok{ f, }\DataTypeTok{Functor}\NormalTok{ g) }\OtherTok{=\textgreater{}} \DataTypeTok{Bifunctor}\NormalTok{ (}\DataTypeTok{Biff}\NormalTok{ p f g) }\KeywordTok{where}
\NormalTok{    bimap f1 f2 (}\DataTypeTok{Biff}\NormalTok{ pfg) }\OtherTok{=} \DataTypeTok{Biff}\NormalTok{ (bimap (}\FunctionTok{fmap}\NormalTok{ f1) (}\FunctionTok{fmap}\NormalTok{ f2) pfg)}
\end{Highlighting}
\end{Shaded}

\texttt{Biff} mathematically proves that if you take a base Bifunctor
(\(p\)) and compose it with two Functors (\(f\) and \(g\)), the
structure is mathematically guaranteed to generate a perfectly lawful,
brand-new \textbf{Bifunctor}!

For example, \texttt{Biff\ Either\ {[}{]}\ Maybe\ a\ b} geometrically
creates \texttt{Either\ {[}a{]}\ (Maybe\ b)}. Because \texttt{Either},
\texttt{List}, and \texttt{Maybe} are completely lawful atoms,
\texttt{Biff} automatically writes \texttt{bimap} for you by natively
mapping the left function over the list and the right function over the
\texttt{Maybe} branch. This flawlessly bridges 1D Functors and 2D
Bifunctors in our mathematical closed algebraic system!

\subsubsection{Section 1.7: Polynomial
Functors}\label{section-1.7-polynomial-functors}

The relationship between Category Theory and Haskell's \textbf{Algebraic
Data Types (ADTs)} is formalized through \textbf{Polynomial Functors}.

If a functor is built solely from: - \textbf{Constants}:
\texttt{Const\ r} (\(C\) or \(1\)) - \textbf{Identity}:
\texttt{Identity} (\(X\)) - \textbf{Sums}: \texttt{Either} (\(+\)) -
\textbf{Products}: Tuples (\(\times\))

\ldots{} it is a \textbf{Polynomial Functor}. Most standard Haskell ADTs
(like \texttt{Maybe}, \texttt{Either}, and non-recursive records) are
polynomial. They are the ``algebra'' of types, where complex structures
are discovered by summing and multiplying simpler ones.

\paragraph{Why the Name ``Polynomial''?}\label{why-the-name-polynomial}

The terminology is beautifully literal. Think about a regular algebraic
polynomial from high school math, like \(F(X) = 1 + Int + X^2\). It is
built using exactly the same operations: * \textbf{\(X\)}: The variable
(The Identity Functor). * \textbf{\(1, Int\)}: Constants (The Constant
Functor). * \textbf{Multiplication (\(X^2 = X \times X\))}: Products
(Tuples \texttt{(a,\ a)}). * \textbf{Addition (\(+\))}: Sums
(\texttt{Either} or alternative constructors).

When we build an Algebraic Data Type (ADT) in Haskell, we are quite
literally writing a polynomial equation. For example, consider this
functor:

\begin{Shaded}
\begin{Highlighting}[]
\KeywordTok{data} \DataTypeTok{Shape}\NormalTok{ a }\OtherTok{=} \DataTypeTok{Empty} \OperatorTok{|} \DataTypeTok{Point} \DataTypeTok{Int} \OperatorTok{|} \DataTypeTok{Line}\NormalTok{ a a}
\end{Highlighting}
\end{Shaded}

If we translate this to algebra using our building blocks: *
\texttt{Empty} has zero parameters: It is \(1\) (a constant,
\texttt{Proxy}). * \texttt{Point\ Int} has an \texttt{Int} but no
parameter \texttt{a}: It is the constant \(Int\). * \texttt{Line\ a\ a}
has two parameters (a pair): It is the product of identity with itself,
\(X \times X = X^2\).

So, the polynomial functor shape for \texttt{Shape\ a} is mathematically
written as: \textbf{\(F(X) = 1 + Int + X^2\)}

\subsubsection{Section 1.8: The Parallel Functor Ecosystem (Solutions
for Restricted
Functors)}\label{section-1.8-the-parallel-functor-ecosystem-solutions-for-restricted-functors}

As we briefly highlighted in Section 2.1, the mathematical definition of
a functor is far broader than Haskell's native \texttt{Functor}
typeclass (which strictly maps \texttt{Type\ -\textgreater{}\ Type}
unconstrained). When structures inevitably violate these two rules, we
do not throw our hands up in defeat!

The Haskell ecosystem simply defines \emph{parallel} typeclasses to
capture these different categorical mappings, allowing us to retain the
exact same structural guarantees.

\paragraph{\texorpdfstring{1. The Too-Wide Functor:
\texttt{Bifunctor}}{1. The Too-Wide Functor: Bifunctor}}\label{the-too-wide-functor-bifunctor}

If a structure has a kind of
\texttt{Type\ -\textgreater{}\ Type\ -\textgreater{}\ Type} (like
\texttt{Either} or \texttt{(,)}), it is a perfectly valid functor
mapping from the product category \(Hask \times Hask \to Hask\). Because
it requires two types, we use \texttt{Data.Bifunctor}:

\begin{Shaded}
\begin{Highlighting}[]
\KeywordTok{class} \DataTypeTok{Bifunctor}\NormalTok{ p }\KeywordTok{where}
\OtherTok{    bimap ::}\NormalTok{ (a }\OtherTok{{-}\textgreater{}}\NormalTok{ b) }\OtherTok{{-}\textgreater{}}\NormalTok{ (c }\OtherTok{{-}\textgreater{}}\NormalTok{ d) }\OtherTok{{-}\textgreater{}}\NormalTok{ p a c }\OtherTok{{-}\textgreater{}}\NormalTok{ p b d}
\end{Highlighting}
\end{Shaded}

\emph{(We will completely deconstruct these in
\hyperref[chapter-4-deep-dive-into-bifunctors]{Chapter 4}).}

\paragraph{\texorpdfstring{2. The Reverse Functor:
\texttt{Contravariant}}{2. The Reverse Functor: Contravariant}}\label{the-reverse-functor-contravariant}

A standard Functor maps ``covariant'' inputs (it \emph{produces}
values). But what if a structure only \emph{consumes} values? This is
mathematically a functor mapped from the opposite category:
\(Hask^{op} \to Hask\).

If you have a \texttt{Predicate\ a} (a wrapper around
\texttt{a\ -\textgreater{}\ Bool}), you can't map its output
(\texttt{Bool}), but you can map its input!

\begin{Shaded}
\begin{Highlighting}[]
\KeywordTok{class} \DataTypeTok{Contravariant}\NormalTok{ f }\KeywordTok{where}
\OtherTok{    contramap ::}\NormalTok{ (a }\OtherTok{{-}\textgreater{}}\NormalTok{ b) }\OtherTok{{-}\textgreater{}}\NormalTok{ f b }\OtherTok{{-}\textgreater{}}\NormalTok{ f a  }\CommentTok{{-}{-} Notice the reversed \textquotesingle{}b\textquotesingle{} and \textquotesingle{}a\textquotesingle{}!}
\end{Highlighting}
\end{Shaded}

\paragraph{\texorpdfstring{3. The Mixed Functor:
\texttt{Profunctor}}{3. The Mixed Functor: Profunctor}}\label{the-mixed-functor-profunctor}

If a Bifunctor maps two covariant types, a \textbf{Profunctor} is a
mapping over one contravariant shape and one covariant shape. The
standard function arrow \texttt{(-\textgreater{})} is a Profunctor.

\begin{Shaded}
\begin{Highlighting}[]
\KeywordTok{class} \DataTypeTok{Profunctor}\NormalTok{ p }\KeywordTok{where}
\OtherTok{    dimap ::}\NormalTok{ (a }\OtherTok{{-}\textgreater{}}\NormalTok{ b) }\OtherTok{{-}\textgreater{}}\NormalTok{ (c }\OtherTok{{-}\textgreater{}}\NormalTok{ d) }\OtherTok{{-}\textgreater{}}\NormalTok{ p b c }\OtherTok{{-}\textgreater{}}\NormalTok{ p a d}
\end{Highlighting}
\end{Shaded}

\texttt{dimap} allows you to simultaneously map the \emph{incoming}
argument (before the function runs) and the \emph{outgoing} result
(after the function runs). They form the categorical backbone of the
\texttt{lens} library.

\paragraph{\texorpdfstring{4. The Constrained Functor:
\texttt{MonoFunctor} (The \texttt{mono-traversable}
library)}{4. The Constrained Functor: MonoFunctor (The mono-traversable library)}}\label{the-constrained-functor-monofunctor-the-mono-traversable-library}

Recall that \texttt{Data.Set} fails to be a \texttt{Functor} because
rebuilding its internal tree requires an \texttt{Ord\ a} constraint on
mapping. It is a ``Restricted Functor'' mapping only onto a subcategory.

Similarly, structures like \texttt{ByteString} or \texttt{Text} aren't
parametric at all (they have kind \texttt{Type}), but logically act
precisely like containers. To solve this, Michael Snoyman's
\texttt{mono-traversable} library created the \texttt{MonoFunctor}
typeclass:

\begin{Shaded}
\begin{Highlighting}[]
\KeywordTok{class} \DataTypeTok{MonoFunctor}\NormalTok{ mono }\KeywordTok{where}
\OtherTok{    omap ::}\NormalTok{ (}\DataTypeTok{Element}\NormalTok{ mono }\OtherTok{{-}\textgreater{}} \DataTypeTok{Element}\NormalTok{ mono) }\OtherTok{{-}\textgreater{}}\NormalTok{ mono }\OtherTok{{-}\textgreater{}}\NormalTok{ mono}
\end{Highlighting}
\end{Shaded}

This allows us to maintain the interface and laws of a Functor over
mathematically restricted or entirely monomorphic structures.

\subsubsection{Section 1.9: Discovering Molecules
(Compounds)}\label{section-1.9-discovering-molecules-compounds}

Using these ``atoms,'' let's see how we can discover the rest of the
Haskell universe.

\paragraph{\texorpdfstring{1. The Sum Molecule:
\texttt{Maybe}}{1. The Sum Molecule: Maybe}}\label{the-sum-molecule-maybe}

If we take the \textbf{Sum} (\texttt{+} in algebra, \texttt{Either} in
Haskell) of \texttt{Proxy} (the number \(1\)) and \texttt{Identity}
(\(X\)), we get the structure for choice or failure:
\texttt{Maybe\ a\ ≅\ Sum\ Proxy\ Identity\ a\ ≅\ Either\ ()\ a}
\textbf{Algebraically}: \(1 + X\)

\paragraph{\texorpdfstring{2. The Product Molecule:
\texttt{Writer}}{2. The Product Molecule: Writer}}\label{the-product-molecule-writer}

If we take the \textbf{Product} (\(\times\) in algebra, a Tuple in
Haskell) of a constant \texttt{Const\ r} and \texttt{Identity} (\(X\)),
we get a structure that carries a ``log'' along with the value:
\texttt{Writer\ r\ a\ ≅\ Product\ (Const\ r)\ Identity\ a\ ≅\ (r,\ a)}
\textbf{Algebraically}: \(r \times X\)

\emph{(Note:
\texttt{Proxy\ *\ Identity\ ≅\ ((),\ a)\ ≅\ a\ ≅\ Identity}. Proxy acts
as the number \(1\) in multiplication).}

\paragraph{\texorpdfstring{3. The Infinite Chain:
\texttt{List}}{3. The Infinite Chain: List}}\label{the-infinite-chain-list}

By using both Sums and Products with \textbf{Recursion}, we can build a
list. A list is either empty (\texttt{Proxy}) OR a head and a tail
(\texttt{Product\ Identity\ List}).
\texttt{List\ a\ ≅\ Sum\ Proxy\ (Product\ Identity\ List)\ a}
\textbf{Algebraically}: \(L(X) = 1 + X \times L(X)\)

\begin{quote}
\textbf{Is \(1 + X \times W = W\) always the case?} Looking at the list
equation, you might ask: ``is it always the case that
\texttt{Sum\ Proxy\ (Product\ Identity\ Whatever)\ =\ Whatever}?'' The
answer is no! The formula \(1 + X \times W\) describes the ``shape'' of
a single layer of a List. When we say \(L(X) = 1 + X \times L(X)\), we
are saying that \texttt{List} is exactly the type that satisfies this
equation (it is the \emph{Fixed Point} of that functor). If
\texttt{Whatever} was a Binary Tree, its shape equation would look
entirely different, such as \(T(X) = 1 + X \times T(X) \times T(X)\).
\end{quote}

\begin{center}\rule{0.5\linewidth}{0.5pt}\end{center}

\subsection{Chapter 2: Foldable (Lossy
Aggregation)}\label{chapter-2-foldable-lossy-aggregation}

While Functors map values and Applicatives/Monads sequence them, a
\texttt{Foldable} is fundamentally about \emph{aggregating} or
destroying a structure down to a summary value.

\subsubsection{Section 7.1: What is a
Foldable?}\label{section-7.1-what-is-a-foldable}

At its core, a \texttt{Foldable} is a typeclass that abstracts the idea
of ``walking through'' a data structure and squashing all of its
elements together.

\paragraph{1. A Well-Kinded Type
Constructor}\label{a-well-kinded-type-constructor}

Before anything else, a type must have the correct ``shape'' to be
Foldable. Mathematically, \texttt{Foldable} is a property of a type
constructor of kind \texttt{Type\ -\textgreater{}\ Type} (like
\texttt{List} or \texttt{Maybe}). It describes a container that holds
some type \texttt{a} (\texttt{t\ a}). Because of this strict kind
signature, absolute atomic concrete types like \texttt{Int},
\texttt{Double}, or the uninhabited type \texttt{Void} (which all
possess kind \texttt{Type}) mathematically cannot be \texttt{Foldable}.
You cannot fold an \texttt{Int} because there's no generic type
parameter \texttt{a} to map over!

\paragraph{\texorpdfstring{2. Unconstrained Morphism Mapping
(\texttt{foldMap})}{2. Unconstrained Morphism Mapping (foldMap)}}\label{unconstrained-morphism-mapping-foldmap}

In Haskell's \texttt{Data.Foldable} class, there are dozens of functions
available (such as \texttt{length}, \texttt{null}, \texttt{toList}, and
\texttt{foldl}). However, to make your type an instance of
\texttt{Foldable}, you mathematically only need to provide exactly
\textbf{one} of two core functions (The Minimal Complete Definition): 1.
\texttt{foldMap} 2. \texttt{foldr}

If you provide just \texttt{foldMap}, Haskell automatically derives
\texttt{foldr} (using the \texttt{Endo} monoid under the hood). Every
single other \texttt{Foldable} function is derived for free from
whichever of those two you choose to implement!

While you can implement \texttt{Foldable} using standard right-folds
(\texttt{foldr}), the most mathematically elegant way to understand it
is through \texttt{foldMap}:

\begin{Shaded}
\begin{Highlighting}[]
\FunctionTok{foldMap}\OtherTok{ ::} \DataTypeTok{Monoid}\NormalTok{ m }\OtherTok{=\textgreater{}}\NormalTok{ (a }\OtherTok{{-}\textgreater{}}\NormalTok{ m) }\OtherTok{{-}\textgreater{}}\NormalTok{ t a }\OtherTok{{-}\textgreater{}}\NormalTok{ m}
\end{Highlighting}
\end{Shaded}

This signature tells a crystal-clear story of \textbf{parametricity}
(The Parametricity Constraints): * The \texttt{m} in this signature is
universally quantified
(\texttt{forall\ m.\ Monoid\ m\ =\textgreater{}\ ...}). * This means
that by the laws of parametricity, your implementation of
\texttt{foldMap} \textbf{cannot possibly know} which specific monoid the
user has chosen. It has no idea if the user is using \texttt{Sum},
\texttt{Product}, \texttt{List}, or \texttt{Any}. * Therefore, your
\texttt{Foldable} instance \emph{must} work blindly and uniformly for
\textbf{any} mathematically valid monoid. The only tools your
implementation is legally allowed to use to collapse the structure are
the monoid's \texttt{mempty} and \texttt{mappend}
(\texttt{\textless{}\textgreater{}}). * Because of this, we are
restricted from inspecting the structure dynamically. If we possessed a
naive signature like
\texttt{(a\ -\textgreater{}\ m)\ -\textgreater{}\ t\ a\ -\textgreater{}\ m}
\emph{without} any external laws, a developer could traverse the
structure backwards, skip every second element, or duplicate elements
randomly, and the compiler would not complain.

\paragraph{3. Mathematical Laws}\label{mathematical-laws-1}

To prevent absolute chaos across these different traversals,
\texttt{Foldable} instances must obey mathematical laws ensuring
consistency across different modes of traversal. While the compiler
cannot enforce these, they are mathematically required.

Unlike \texttt{Functor} or \texttt{Monad} which have rigorous
categorical laws (Identity and Composition), \texttt{Foldable} is
somewhat unique: its laws are primarily \textbf{consistency laws}. There
isn't just one, but a family of required equivalences ensuring that all
the derived folding methods agree with each other.

The primary \textbf{consistency equalities} demand that \texttt{foldMap}
is structurally isomorphic to both sequential folding methods
(\texttt{foldr} and \texttt{foldl}). * \textbf{The Right-Fold Law:}
\texttt{foldMap\ f\ ==\ foldr\ (mappend\ .\ f)\ mempty} * \textbf{The
Left-Fold Law:}
\texttt{foldMap\ f\ ==\ foldl\ (\textbackslash{}acc\ x\ -\textgreater{}\ acc\ \textless{}\textgreater{}\ f\ x)\ mempty}

If you map elements to a monoid and combine them sequentially, it
\emph{must} yield the exact same result as using a right/left-fold that
applies \texttt{f} and strictly \texttt{mappend}s the accumulation.

\paragraph{Category Theory Origin: Destructive
Traversals}\label{category-theory-origin-destructive-traversals}

Categorically, \texttt{Foldable} represents an explicitly \emph{lossy}
operation. Unlike \texttt{Functor} which rigidly preserves the ``shape''
of the data, a Foldable traversal inherently destroys the structural
geometry of the wrapper \texttt{t} and projects the data down onto a
Monoid.

In Category Theory, there is a profound insight lying at the bottom of
the \texttt{Foldable} hierarchy: \texttt{foldMap} is literally just
\texttt{traverse} using the \texttt{Const} Applicative Functor! If you
use \texttt{traverse} with \texttt{Const\ m}, you are running an
applicative computation that strictly ignores the purely computational
\texttt{a} part and only accumulates the contextual \texttt{Monoid\ m}
part. Because you are accumulating the monoid and throwing away the
structure, \texttt{traverse} geometrically degrades into a purely
destructive fold.

\subsubsection{Section 7.2: The Absolute Minimum
Foldable}\label{section-7.2-the-absolute-minimum-foldable}

Before we can conceptually fold a structure using \texttt{foldMap}, we
must supply it with its first argument: a monoidal mapping function
\texttt{(a\ -\textgreater{}\ m)}.

\begin{quote}
{[}!NOTE{]} \textbf{What is \texttt{a\ -\textgreater{}\ m} in Category
Theory?} In abstract algebra and Category Theory,
\texttt{a\ -\textgreater{}\ m} is known as a \textbf{generator map} (or
an insertion mapping) from a simple Set (\texttt{a}) into the underlying
set of a Monoid (\texttt{m}). By the \textbf{Universal Property of the
Free Monoid}, any such simple set-theoretic mapping is mathematically
guaranteed to uniquely extend into a rigorous \textbf{Monoid
Homomorphism} (a structure-preserving aggregation) from the Free Monoid
(\texttt{{[}a{]}}) to \texttt{m}.

The \texttt{foldMap} function is the pure Haskell realization of this
profound mathematical adjunction: it takes your humble generator map
\texttt{a\ -\textgreater{}\ m} and elegantly elevates it into a
universal structural fold!
\end{quote}

What are the absolute simplest mathematical mappings we can create to
define the minimal folding behavior?

\paragraph{\texorpdfstring{The Minimal Monoidal Mappings
(\texttt{a\ -\textgreater{}\ m})}{The Minimal Monoidal Mappings (a -\textgreater{} m)}}\label{the-minimal-monoidal-mappings-a---m}

\begin{enumerate}
\def\labelenumi{\arabic{enumi}.}
\item
  \textbf{The Empty Mapping (\texttt{Void\ -\textgreater{}\ m})}:
  Because \texttt{Void} is the initial object in Haskell, there exists a
  unique, mathematically rigorous function from \texttt{Void} to any
  arbitrarily chosen monoid \texttt{m}:
  \texttt{absurd\ ::\ Void\ -\textgreater{}\ m}. Because a value of
  \texttt{Void} can never be physically constructed, this mapping
  function is never actually executed at runtime. However, it exists
  mathematically to perfectly satisfy the typechecker when we are forced
  to fold over logically empty structures (like \texttt{Proxy\ Void}) or
  branches that have been proven geometrically impossible!
\item
  \textbf{The Trivial Mapping (\texttt{a\ -\textgreater{}\ ()})}: The
  \texttt{()} type is the terminal monoid. If we map
  \texttt{\textbackslash{}x\ -\textgreater{}\ ()} (or equivalently use
  the function \texttt{const\ ()}), we completely erase every element.
  When aggregated together
  (\texttt{()\ \textless{}\textgreater{}\ ()\ \textless{}\textgreater{}\ ...\ \textless{}\textgreater{}\ ()}),
  the result is just \texttt{()}. This is the ultimate ``destroyer of
  information'', completely collapsing both the data and the structural
  shape into the void.
\item
  \textbf{The Constant Counting Mapping
  (\texttt{a\ -\textgreater{}\ Sum\ 1})}: What if we throw away the
  value of the element, but replace it with a mathematical \emph{tick}?
  By using \texttt{\textbackslash{}x\ -\textgreater{}\ Sum\ 1} (or
  \texttt{const\ (Sum\ 1)}), every element becomes a \texttt{1}. When
  aggregated via \texttt{Sum}, this mathematically calculates the
  \textbf{length} of the structure!

\begin{Shaded}
\begin{Highlighting}[]
\CommentTok{{-}{-} An elegant, universal way to calculate length across ANY Foldable!}
\OtherTok{len ::} \DataTypeTok{Foldable}\NormalTok{ t }\OtherTok{=\textgreater{}}\NormalTok{ t a }\OtherTok{{-}\textgreater{}} \DataTypeTok{Int}
\NormalTok{len xs }\OtherTok{=}\NormalTok{ getSum (}\FunctionTok{foldMap}\NormalTok{ (}\FunctionTok{const}\NormalTok{ (}\DataTypeTok{Sum} \DecValTok{1}\NormalTok{)) xs)}
\end{Highlighting}
\end{Shaded}

  \begin{quote}
  {[}!NOTE{]} \textbf{Other Constant Mappings} The power of this pattern
  is that you can swap out the monoid! If you instead mapped
  \texttt{\textbackslash{}x\ -\textgreater{}\ Any\ True}, the
  aggregation calculates whether the structure is non-empty
  (\texttt{not\ .\ null}). If you map
  \texttt{\textbackslash{}x\ -\textgreater{}\ Product\ 2}, it calculates
  \(2^{\text{length}}\) (the number of possible subsets). The behavior
  is purely dictated by the monoid!
  \end{quote}
\item
  \textbf{The Accumulating Mapping (\texttt{a\ -\textgreater{}\ Sum\ a},
  \texttt{a\ -\textgreater{}\ Product\ a}, etc.)}: Unlike the previous
  mappings that forcibly erase the value of the elements, we can
  \emph{pass the data forward} for accumulation by wrapping them in a
  constructor like \texttt{\textbackslash{}x\ -\textgreater{}\ Sum\ x}
  (or simply \texttt{Sum}). While the value \texttt{a} is strictly
  preserved during the mapping step, the subsequent aggregation
  (\texttt{\textless{}\textgreater{}} for \texttt{Sum}) mathematically
  squashes everything together into a single summarized value.

  \begin{quote}
  {[}!NOTE{]} \textbf{Other Accumulating Mappings} This is the exact
  same mechanism used for other simple accumulating monoids! You can
  equivalently swap \texttt{Sum} with \texttt{Product} to multiply all
  elements (\texttt{foldMap\ Product}), or map booleans into
  \texttt{Any} / \texttt{All} to aggregate logical conditions across the
  entire Foldable structure.
  \end{quote}

\begin{Shaded}
\begin{Highlighting}[]
\CommentTok{{-}{-} An elegant, universal way to calculate the sum across ANY Foldable!}
\OtherTok{sumElements ::}\NormalTok{ (}\DataTypeTok{Foldable}\NormalTok{ t, }\DataTypeTok{Num}\NormalTok{ a) }\OtherTok{=\textgreater{}}\NormalTok{ t a }\OtherTok{{-}\textgreater{}}\NormalTok{ a}
\NormalTok{sumElements xs }\OtherTok{=}\NormalTok{ getSum (}\FunctionTok{foldMap} \DataTypeTok{Sum}\NormalTok{ xs)}
\end{Highlighting}
\end{Shaded}
\item
  \textbf{The True Data-Preserving Mapping
  (\texttt{a\ -\textgreater{}\ {[}a{]}})}: What if we want to preserve
  \emph{everything} (the exact values and their sequential order) and
  lose \emph{only} the structural container geometry? We map every
  element directly into the Free Monoid! By mapping
  \texttt{\textbackslash{}x\ -\textgreater{}\ {[}x{]}} (or
  \texttt{(:{[}{]})}), we preserve each value perfectly in its own
  isolated List. When the monoid aggregates them using \texttt{++}, we
  get a perfect log of all the elements!

\begin{Shaded}
\begin{Highlighting}[]
\CommentTok{{-}{-} This is the exact mathematical definition of \textasciigrave{}toList\textasciigrave{}!}
\OtherTok{toList\textquotesingle{} ::} \DataTypeTok{Foldable}\NormalTok{ t }\OtherTok{=\textgreater{}}\NormalTok{ t a }\OtherTok{{-}\textgreater{}}\NormalTok{ [a]}
\NormalTok{toList\textquotesingle{} xs }\OtherTok{=} \FunctionTok{foldMap}\NormalTok{ (}\OperatorTok{:}\NormalTok{[]) xs}
\end{Highlighting}
\end{Shaded}
\end{enumerate}

These minimal mappings serve as the absolute bedrock of data
aggregation. By simply swapping out the \texttt{a\ -\textgreater{}\ m}
mapping, \texttt{foldMap} elegantly shifts from forgetting data, to
counting data, to aggregating data!

Now that we understand how elements are mapped into Monoids, let's
explore the absolute minimal structural implementations of
\texttt{Foldable} (\texttt{t\ a}), and rigorously verify that they
fulfill the folding laws.

\paragraph{\texorpdfstring{1. The Mathematically Unreachable Foldable
(\texttt{Zero})}{1. The Mathematically Unreachable Foldable (Zero)}}\label{the-mathematically-unreachable-foldable-zero}

Before we hit \texttt{Proxy}, we should mathematically ask: can
\texttt{Zero} be Foldable? \texttt{Zero} is an uninhabited type (it has
exactly zero constructors).

\begin{Shaded}
\begin{Highlighting}[]
\KeywordTok{data} \DataTypeTok{Zero}\NormalTok{ a }\CommentTok{{-}{-} mathematically empty!}

\KeywordTok{instance} \DataTypeTok{Foldable} \DataTypeTok{Zero} \KeywordTok{where}
    \CommentTok{{-}{-} foldMap :: Monoid m =\textgreater{} (a {-}\textgreater{} m) {-}\textgreater{} Zero a {-}\textgreater{} m}
    \FunctionTok{foldMap}\NormalTok{ \_ \_ }\OtherTok{=} \FunctionTok{mempty}
\end{Highlighting}
\end{Shaded}

Since \texttt{Zero} has zero constructors, it holds zero values of type
\texttt{a}. Because we do not possess an \texttt{a}, we cannot use our
\texttt{(a\ -\textgreater{}\ m)} mapping function. However, to satisfy
\texttt{foldMap}, we \emph{must} produce a Monoid \texttt{m}. How do we
conjure one?

The exact way we implement it is simply by reaching for the Monoid
constraint to pull out its identity element: \texttt{mempty}! Because
\texttt{Zero} contains no data, we ignore the inputs (\texttt{\_\ \_})
and simply yield an empty Monoid.

\begin{quote}
{[}!NOTE{]} \textbf{The Case of Absurdity} Because \texttt{Zero} is
uninhabited and impossible to instantiate at runtime, there is actually
a second, much creepier way to implement this in Haskell without even
using \texttt{mempty}. If you pattern match on an impossible value
(\texttt{foldMap\ \_\ z\ =\ case\ z\ of\ \{\}}), the GHC compiler uses
the logical principle of explosion (or ``absurdity'') to vacuously
satisfy the return type! But \texttt{\_\ \_\ =\ mempty} is far more
readable and idiomatic for a programmer.
\end{quote}

\textbf{Verifying the Law:} We still have to check that the consistency
laws are fulfilled! Does
\texttt{foldMap\ f\ ==\ foldr\ (mappend\ .\ f)\ mempty} hold? Yes. Since
there are absolutely zero elements to fold over, a sequential right-fold
inherently falls back to its base case (\texttt{mempty}). Thus,
\texttt{mempty\ ==\ mempty}. The law holds!

However, because we can never actually instantiate it to hold data, we
must move exactly one step up to find our first \emph{usable} minimal:

\paragraph{\texorpdfstring{2. The Empty Foldable
(\texttt{Proxy})}{2. The Empty Foldable (Proxy)}}\label{the-empty-foldable-proxy}

How do you fold an explicitly empty structure (\texttt{Proxy}) that you
\emph{can} instantiate?

\begin{Shaded}
\begin{Highlighting}[]
\KeywordTok{instance} \DataTypeTok{Foldable} \DataTypeTok{Proxy} \KeywordTok{where}
    \FunctionTok{foldMap}\NormalTok{ \_ \_ }\OtherTok{=} \FunctionTok{mempty}
\end{Highlighting}
\end{Shaded}

Because \texttt{Proxy} holds zero values of type \texttt{a}, it is
impossible to apply our mapping function
\texttt{(a\ -\textgreater{}\ m)}. The type system rigorously enforces
that we must return an \texttt{m} (which is constrained to be a
\texttt{Monoid}). The \emph{only} mathematically sound way to conjure an
\texttt{m} out of thin air without possessing an \texttt{a} is to use
the monoid's identity element: \texttt{mempty} (and just like you noted,
\texttt{\_\ \_\ =\ mempty} is perfectly valid here instead of
\texttt{\_\ Proxy\ =\ mempty} because we don't care about evaluating the
Proxy value itself!).

\textbf{Verifying the Law:} Identical to \texttt{Zero}. There are zero
elements inside the \texttt{Proxy}, so the strictly sequential
right-fold \texttt{foldr\ (mappend\ .\ f)\ mempty} immediately
circumvents the list and returns its initial accumulator
\texttt{mempty}. Our \texttt{foldMap} returns \texttt{mempty}. The law
is perfectly fulfilled \texttt{mempty\ ==\ mempty}.

\paragraph{\texorpdfstring{3. The Single-Element Foldable
(\texttt{Identity})}{3. The Single-Element Foldable (Identity)}}\label{the-single-element-foldable-identity}

How do you fold a structure containing exactly one element?

By the strict rules of parametricity, we know absolutely nothing about
the internal data \texttt{x} nor the resulting Monoid \texttt{m} at
compile-time. The typechecker only guarantees that \texttt{m} possesses
a Monoid constraint and that we have a mapping function
\texttt{(a\ -\textgreater{}\ m)}. Therefore, parametricity logically
dictates there are exactly two structural implementations that compile:

\begin{itemize}
\tightlist
\item
  \textbf{Candidate 1 (The Ignore Cheat):} Ignore the data entirely and
  pull the identity element out of thin air:
  \texttt{foldMap\ \_\ \_\ =\ mempty}.
\item
  \textbf{Candidate 2 (The Application):} Actually use the mapping
  function on our single piece of data:
  \texttt{foldMap\ f\ (Identity\ x)\ =\ f\ x}.
\end{itemize}

Just like we saw with Functors, this is exactly where the mathematical
laws execute a ``Forced Hand''! We must verify both candidates against
the consistency equality:
\texttt{foldMap\ f\ ==\ foldr\ (mappend\ .\ f)\ mempty}.

\textbf{Testing Candidate 1 (\texttt{mempty}):} For exactly one element
\texttt{x}, a rigorous sequential right-fold (\texttt{foldr}) traverses
the structure, applies \texttt{f} to \texttt{x}, and strictly combines
it with the base case: \texttt{f\ x\ \textless{}\textgreater{}\ mempty}.
By the absolute Right Identity Law of Monoids,
\texttt{f\ x\ \textless{}\textgreater{}\ mempty} perfectly simplifies to
\texttt{f\ x}. If we use Candidate 1, the Foldable consistency law
(\texttt{foldMap\ f\ ==\ foldr\ (mappend\ .\ f)\ mempty}) demands that
our \texttt{foldMap} result (\texttt{mempty}) must exactly equal the
fully evaluated right-fold result (\texttt{f\ x}). Because \texttt{f\ x}
could mathematically evaluate to anything (e.g., \texttt{Sum\ 5}),
proposing that \texttt{mempty\ ==\ Sum\ 5} is a blatant universal
contradiction. While we used the Monoid law to simplify the right side
of the equation, it is ultimately the Foldable consistency law that
violently excludes Candidate 1!

\textbf{Testing Candidate 2 (\texttt{f\ x}):} If we use Candidate 2, our
\texttt{foldMap} yields \texttt{f\ x}. Does this equal our
mathematically evaluated right-fold result of
\texttt{f\ x\ \textless{}\textgreater{}\ mempty} (which simplifies to
\texttt{f\ x} via the Monoid law)? Yes, \texttt{f\ x\ ==\ f\ x}.
Candidate 2 perfectly respects the Foldable consistency laws and
leverages the Monoid laws correctly!

Therefore, parametricity establishes the only two possible paths in the
universe, and the Foldable consistency laws strictly \textbf{force our
hand} to choose Candidate 2. Assessed together, there is exactly one
mathematically legal implementation!

\begin{Shaded}
\begin{Highlighting}[]
\KeywordTok{instance} \DataTypeTok{Foldable} \DataTypeTok{Identity} \KeywordTok{where}
    \FunctionTok{foldMap}\NormalTok{ f (}\DataTypeTok{Identity}\NormalTok{ x) }\OtherTok{=}\NormalTok{ f x}
\end{Highlighting}
\end{Shaded}

\paragraph{\texorpdfstring{4. The Ghost Data
(\texttt{Const\ r})}{4. The Ghost Data (Const r)}}\label{the-ghost-data-const-r}

What if your data structure physically holds data in memory, but it's
the \emph{wrong} type parameter?

\begin{Shaded}
\begin{Highlighting}[]
\CommentTok{{-}{-} Notice the kind is \textasciigrave{}Type {-}\textgreater{} Type {-}\textgreater{} Type\textasciigrave{}. We fold over the SECOND parameter \textquotesingle{}a\textquotesingle{}}
\KeywordTok{data} \DataTypeTok{Const}\NormalTok{ r a }\OtherTok{=} \DataTypeTok{Const}\NormalTok{ r}

\KeywordTok{instance} \DataTypeTok{Foldable}\NormalTok{ (}\DataTypeTok{Const}\NormalTok{ r) }\KeywordTok{where}
    \FunctionTok{foldMap}\NormalTok{ \_ \_ }\OtherTok{=} \FunctionTok{mempty}
\end{Highlighting}
\end{Shaded}

Even though \texttt{Const} physically holds data (a value of type
\texttt{r}), it holds exactly \emph{zero} values of the generic type
\texttt{a} we are folding over. Therefore, relative to our mapping
function \texttt{(a\ -\textgreater{}\ m)}, the structure is effectively
empty! Just like \texttt{Zero} and \texttt{Proxy}, we cannot call
\texttt{f}, and we are mathematically forced to pull \texttt{mempty} out
of thin air to satisfy the return constraint.

\paragraph{\texorpdfstring{5. The Static Pairing
(\texttt{(e,\ a)})}{5. The Static Pairing ((e, a))}}\label{the-static-pairing-e-a}

What if our structure holds both the wrong type AND the right type?

\begin{Shaded}
\begin{Highlighting}[]
\KeywordTok{instance} \DataTypeTok{Foldable}\NormalTok{ ((,) e) }\KeywordTok{where}
    \FunctionTok{foldMap}\NormalTok{ f (\_, x) }\OtherTok{=}\NormalTok{ f x}
\end{Highlighting}
\end{Shaded}

By parametricity, our mapping function \texttt{f} only operates on
\texttt{a}, making the \texttt{e} part of the tuple (the left side)
mathematically useless to our fold. We simply throw it away. Because we
possess exactly one \texttt{x}, the Foldable laws execute the exact same
``Forced Hand'' we proved for \texttt{Identity}, strictly locking us
into returning \texttt{f\ x}.

\textbf{Verifying the Law:} Just like \texttt{Identity}, a rigorous
sequential right-fold traverses the tuple, applies \texttt{f} to the
single valid parameter \texttt{x}, and strictly combines it with the
base case: \texttt{f\ x\ \textless{}\textgreater{}\ mempty}. By the
absolute identity laws of Monoids, this effortlessly simplifies down to
precisely \texttt{f\ x}. If we attempted to cheat and blindly return
\texttt{mempty} instead of \texttt{f\ x}, the consistency law would
violently fail as \texttt{mempty\ ==\ f\ x} poses a contradiction.
Therefore, parametricity combined with the Foldable law strictly
enforces \texttt{f\ x}.

\paragraph{\texorpdfstring{6. The Branching Possibility
(\texttt{Either\ e})}{6. The Branching Possibility (Either e)}}\label{the-branching-possibility-either-e}

Finally, what if we have a structure that \emph{sometimes} has an
\texttt{a} (like \texttt{Identity}), and \emph{sometimes} doesn't (like
\texttt{Const})?

\begin{Shaded}
\begin{Highlighting}[]
\KeywordTok{instance} \DataTypeTok{Foldable}\NormalTok{ (}\DataTypeTok{Either}\NormalTok{ e) }\KeywordTok{where}
    \FunctionTok{foldMap}\NormalTok{ f (}\DataTypeTok{Right}\NormalTok{ x) }\OtherTok{=}\NormalTok{ f x}
    \FunctionTok{foldMap}\NormalTok{ \_ (}\DataTypeTok{Left}\NormalTok{ \_)  }\OtherTok{=} \FunctionTok{mempty}
\end{Highlighting}
\end{Shaded}

Parametricity and the laws perfectly fuse our previous proofs together
based entirely on the shape of the branch! * Within the \texttt{Right}
branch, we possess exactly one \texttt{a}. Like \texttt{Identity}, the
consistency laws force us to return \texttt{f\ x}. * Within the
\texttt{Left} branch, we possess exactly zero \texttt{a}s. Like
\texttt{Const}, the mathematical vacuum forces us to return
\texttt{mempty}!

\paragraph{\texorpdfstring{7. The Homomorphism (The Essence of
\texttt{toList})}{7. The Homomorphism (The Essence of toList)}}\label{the-homomorphism-the-essence-of-tolist}

If you can aggressively fold any structure down into a list, you can
fold it. The \texttt{Foldable} class essentially guarantees that your
structure can be flattened into a standard list via \texttt{toList}. In
mathematics, this means there is a \textbf{homomorphism} (a
structure-preserving map) from your specific type to \texttt{{[}a{]}}.

But what exactly is the structure being preserved if \texttt{Foldable}
\emph{destroys} the original geometry (like flattening a tree into a
line)?

The structure being preserved is the \textbf{sequential monoidal
composition}! A homomorphism strictly guarantees that the following
mathematical equality unconditionally holds:

\begin{Shaded}
\begin{Highlighting}[]
\CommentTok{{-}{-} Folding the raw structure is EXACTLY equivalent to }
\CommentTok{{-}{-} folding its flattened list representation!}
\FunctionTok{foldMap}\NormalTok{ f xs }\OperatorTok{==} \FunctionTok{foldMap}\NormalTok{ f (toList xs)}
\end{Highlighting}
\end{Shaded}

This structural proof reveals a profound secret about Haskell:
\textbf{every \texttt{Foldable} is literally just a List in disguise} as
far as aggregation is concerned! The \texttt{toList} function acts as
the universal algebraic projector. It mathematically maps any exotic
geometric data structure directly onto the ``free monoid'' (a sequential
list \texttt{{[}a{]}}).

Once your elements are mathematically aligned into a list, folding them
reduces to nothing more than inserting
\texttt{\textless{}\textgreater{}} directly between each adjacent
element!

\begin{quote}
{[}!NOTE{]} \textbf{Wait, is any \texttt{Foldable\ t} a Free Monoid?}
No! In both mathematics and Haskell, the title of ``Free Monoid'' is
strictly reserved for the List (\texttt{{[}a{]}}). A ``free'' object
over a set of generators \texttt{a} structures the elements with
\emph{absolutely no other constraints or laws} other than the required
category operations (associativity and identity). If we defined a Monoid
using \texttt{Set\ a} and unions, it would not be ``free'' because it
enforces \(x \cup x = x\) (idempotency) and \(x \cup y = y \cup x\)
(commutativity). The List enforces \emph{nothing} except the exact
sequenced order you provided.
\texttt{{[}1{]}\ \textless{}\textgreater{}\ {[}2{]}\ \textless{}\textgreater{}\ {[}1{]}}
is just \texttt{{[}1,\ 2,\ 1{]}}.

A \texttt{Foldable\ t} is simply any data structure that possesses a
natural transformation down into the Free Monoid. The mathematical
Universal Property of the Free Monoid states that for any mapping
\texttt{a\ -\textgreater{}\ m}, there is a unique monoid homomorphism
from \texttt{{[}a{]}\ -\textgreater{}\ m}. When you call
\texttt{foldMap\ f} on a generic \texttt{Foldable\ t}, you are
mathematically flattening your structure into the Free Monoid
(\texttt{toList}), and then immediately using its universal property to
compute the final \texttt{m} (\texttt{foldMap\_List\ f\ .\ toList}).
\end{quote}

\subsubsection{Section 7.3: The Algebra of
Foldables}\label{section-7.3-the-algebra-of-foldables}

Just like Functors and Bifunctors, the \texttt{Foldable} typeclass
strictly shares the same structural shape
(\texttt{Type\ -\textgreater{}\ Type}). Because of this, it inherently
possesses the exact same magnificent algebraic composition rules!

We can mathematically prove that if \texttt{f} and \texttt{g} are both
valid \texttt{Foldable} structures, then their \textbf{Sum},
\textbf{Product}, and \textbf{Composition} are also mathematically
guaranteed to be seamlessly \texttt{Foldable}.

\paragraph{\texorpdfstring{1. Foldable Sums
(\texttt{f\ +\ g})}{1. Foldable Sums (f + g)}}\label{foldable-sums-f-g}

A Sum means you either provide the \texttt{f} structure or the
\texttt{g} structure. To fold it, we simply delegate the fold to
whichever structure was provided:

\begin{Shaded}
\begin{Highlighting}[]
\KeywordTok{data} \DataTypeTok{Sum}\NormalTok{ f g a }\OtherTok{=} \DataTypeTok{InL}\NormalTok{ (f a) }\OperatorTok{|} \DataTypeTok{InR}\NormalTok{ (g a)}

\KeywordTok{instance}\NormalTok{ (}\DataTypeTok{Foldable}\NormalTok{ f, }\DataTypeTok{Foldable}\NormalTok{ g) }\OtherTok{=\textgreater{}} \DataTypeTok{Foldable}\NormalTok{ (}\DataTypeTok{Sum}\NormalTok{ f g) }\KeywordTok{where}
    \FunctionTok{foldMap}\NormalTok{ f (}\DataTypeTok{InL}\NormalTok{ fa) }\OtherTok{=} \FunctionTok{foldMap}\NormalTok{ f fa}
    \FunctionTok{foldMap}\NormalTok{ f (}\DataTypeTok{InR}\NormalTok{ ga) }\OtherTok{=} \FunctionTok{foldMap}\NormalTok{ f ga}
\end{Highlighting}
\end{Shaded}

\begin{quote}
{[}!TIP{]} \textbf{What is the Identity Element (Atom) for Sums?} In
abstract algebra, if you have an addition operation (+), you rigorously
require an identity element (0) such that \texttt{X\ +\ 0\ =\ X}.
Because \texttt{Sum} is an addition operator over Foldables, it
mathematically implies the existence of a ``Zero'' Endofunctor atom!
This atom is precisely \texttt{Zero} (or \texttt{Proxy} if we restrict
ourselves to instantiable types).

\textbf{How exactly do you get \texttt{f} back out of
\texttt{Sum\ f\ Proxy}?} Mathematically, \texttt{Sum\ f\ Proxy} states
that you either possess the structure \texttt{f} (in the \texttt{InL}
branch) OR you possess a \texttt{Proxy} (in the \texttt{InR} branch).
But remember from Section 7.2: a \texttt{Proxy} holds exactly zero
values of our generic parameter \texttt{a}! Because the \texttt{InR}
branch structurally contains absolutely zero data to fold over, the
mathematical ``weight'' of the data strictly remains entirely in the
\texttt{InL} branch. If you iterate over \texttt{Sum\ f\ Proxy}, any
execution path that enters \texttt{InR} instantly returns an empty
Monoid (\texttt{mempty}), which perfectly vanishes during accumulation
(\texttt{\textless{}\textgreater{}\ mempty}). Therefore, algorithmically
and structurally: \texttt{Sum\ f\ Proxy\ ≅\ f}. If you take the Sum of
any structure \texttt{f} and an empty \texttt{Proxy}, it is structurally
isomorphic to just possessing the structure \texttt{f}, perfectly
proving that \texttt{Proxy} acts as the definitive Identity Atom for
Foldable Sums!
\end{quote}

\paragraph{\texorpdfstring{2. Foldable Products
(\texttt{f\ *\ g})}{2. Foldable Products (f * g)}}\label{foldable-products-f-g}

A Product means you possess both structures at the exact same time. To
perform a uniform fold over both spaces simultaneously, you aggressively
\texttt{foldMap} the \texttt{f} structure, aggressively \texttt{foldMap}
the \texttt{g} structure, and then violently smash their resulting
Monoids together using the mathematically required
\texttt{\textless{}\textgreater{}} operator!

\begin{Shaded}
\begin{Highlighting}[]
\KeywordTok{data} \DataTypeTok{Product}\NormalTok{ f g a }\OtherTok{=} \DataTypeTok{Pair}\NormalTok{ (f a) (g a)}

\KeywordTok{instance}\NormalTok{ (}\DataTypeTok{Foldable}\NormalTok{ f, }\DataTypeTok{Foldable}\NormalTok{ g) }\OtherTok{=\textgreater{}} \DataTypeTok{Foldable}\NormalTok{ (}\DataTypeTok{Product}\NormalTok{ f g) }\KeywordTok{where}
    \FunctionTok{foldMap}\NormalTok{ f (}\DataTypeTok{Pair}\NormalTok{ fa ga) }\OtherTok{=} \FunctionTok{foldMap}\NormalTok{ f fa }\OperatorTok{\textless{}\textgreater{}} \FunctionTok{foldMap}\NormalTok{ f ga}
\end{Highlighting}
\end{Shaded}

\begin{quote}
{[}!TIP{]} \textbf{What is the Identity Element (Atom) for Products?} If
you have a multiplication operation (*), you rigorously require an
identity element (1) such that \texttt{X\ *\ 1\ =\ X}. Because
\texttt{Product} is a multiplication operator over Foldables, it
mathematically implies the existence of a ``One'' Endofunctor atom! This
atom is precisely the \texttt{Identity} Functor. Mathematically:
\texttt{Product\ f\ Identity\ ≅\ f}. Having a pair of \texttt{f\ a} and
exactly one \texttt{a} is fundamentally the same operation as iterating
through \texttt{f}, meaning \texttt{Identity} flawlessly acts as the
``1'' Atom!
\end{quote}

\paragraph{The Algebra of Polynomial Functors (No Compose
needed!)}\label{the-algebra-of-polynomial-functors-no-compose-needed}

What happens if we strictly restrict our mathematical toolbox to just
\textbf{Sum (+)} and **Product (*)\textbf{, along with their respective
identity atoms }Proxy (0)\textbf{ and }Identity (1)**, and completely
ignore \texttt{Compose}?

In abstract algebra, any system possessing addition, multiplication, 0,
and 1 (where multiplication distributes over addition) forms a
\textbf{Semiring}. When applied to data structures, this algebraically
generates the magnificent class of \textbf{Polynomial Functors}!

Any data structure built from just Sums, Products, 1s
(\texttt{Identity}), and 0s (\texttt{Proxy}) mathematically corresponds
to a standard polynomial equation with non-negative coefficients. For
example: * \texttt{Identity} is \(x\) *
\texttt{Product\ Identity\ Identity} (a Pair \texttt{(a,\ a)}) is
\(x^2\) * \texttt{Sum\ Identity\ Identity} (an \texttt{Either\ a\ a}) is
\(2x\)

Therefore, a type like
\texttt{data\ BinaryTreeElem\ a\ =\ Empty\ \textbar{}\ Node\ a\ a} is
algebraically represented as the polynomial \(1 + x^2\), composed
entirely without ever requiring function composition.

\textbf{The Profound Realization: Both Functor AND Foldable!} Because we
have mathematically proven that our atoms (\texttt{Proxy},
\texttt{Identity}) and our combinators (\texttt{Sum}, \texttt{Product})
natively provide perfect, lawful implementations for \emph{both}
\texttt{Functor} and \texttt{Foldable} simultaneously\ldots{} this
algebra enforces a breathtaking guarantee: \textbf{Every single
Polynomial Data Type is mathematically guaranteed to be BOTH a lawful
Functor and a lawful Foldable.} You can map over them, and you can
aggregate them. This precise abstract algebra is exactly how the GHC
compiler's \texttt{DeriveFunctor} and \texttt{DeriveFoldable} extensions
work under the hood! They just blindly generate the \texttt{Sum} and
\texttt{Product} instances!

\paragraph{\texorpdfstring{3. Foldable Composition
(\texttt{f\ ∘\ g})}{3. Foldable Composition (f ∘ g)}}\label{foldable-composition-f-g}

Composition means a Foldable deeply nested inside another Foldable (for
example, a \texttt{List} of \texttt{Maybe}s, or a \texttt{Tree} of
\texttt{List}s). To fold it, we elegantly map the inner folding
operation over the entire outer folding operation!

\begin{Shaded}
\begin{Highlighting}[]
\KeywordTok{newtype} \DataTypeTok{Compose}\NormalTok{ f g a }\OtherTok{=} \DataTypeTok{Compose}\NormalTok{ (f (g a))}

\KeywordTok{instance}\NormalTok{ (}\DataTypeTok{Foldable}\NormalTok{ f, }\DataTypeTok{Foldable}\NormalTok{ g) }\OtherTok{=\textgreater{}} \DataTypeTok{Foldable}\NormalTok{ (}\DataTypeTok{Compose}\NormalTok{ f g) }\KeywordTok{where}
    \FunctionTok{foldMap}\NormalTok{ f (}\DataTypeTok{Compose}\NormalTok{ fga) }\OtherTok{=} \FunctionTok{foldMap}\NormalTok{ (}\FunctionTok{foldMap}\NormalTok{ f) fga}
\end{Highlighting}
\end{Shaded}

\begin{quote}
{[}!TIP{]} \textbf{What is the Identity Element (Atom) for Composition?}
For function composition (\(\circ\)), the identity is the structure that
sits invisibly inside or outside another without changing it:
\texttt{X\ ∘\ 1\ =\ X}. Once again, the exact same \texttt{Identity}
Functor mathematically perfectly satisfies this! Mathematically:
\texttt{Compose\ f\ Identity\ ≅\ f} and
\texttt{Compose\ Identity\ f\ ≅\ f}. Composing an \texttt{Identity}
inside or outside of \texttt{f} preserves \texttt{f} exactly. Therefore,
\texttt{Identity} is not just the multiplicative atom (``1''), but also
the compositional atom!
\end{quote}

\paragraph{The Algebra of Monomial Functors (No Sum
needed!)}\label{the-algebra-of-monomial-functors-no-sum-needed}

What happens if we deliberately throw away \textbf{Sum (+)} and the
\textbf{Proxy (0)} zero-atom, and strictly restrict our universe
exclusively to **Product (*)\textbf{, }Compose (\(\circ\))\textbf{, and
the }Identity (1)** atom?

Because we have discarded \texttt{Sum}, we have mathematically banned
``branching'' or ``choice''. We cannot possess a mathematical data type
with multiple constructors (which instantly eliminates \texttt{Maybe},
\texttt{Either}, \texttt{Bool}, and traditional ADTs). Every single data
structure must possess exactly one constructor, and every piece of data
must strictly exist.

If we generate structures strictly using \texttt{Product} and
\texttt{Compose}, we get mathematically perfect \textbf{Monomials}: *
\texttt{Identity} is \(x^1\) * \texttt{Product\ Identity\ Identity} (a
Pair \texttt{(a,\ a)}) is \(x \times x = x^2\) *
\texttt{Product\ Identity\ (Product\ Identity\ Identity)} (a Triple
\texttt{(a,\ a,\ a)}) is \(x \times x^2 = x^3\) *
\texttt{Compose\ (Product\ Identity\ Identity)\ (Product\ Identity\ Identity)}
(a Quadruple or a \texttt{Pair} of \texttt{Pairs}) is \((x^2)^2 = x^4\)

This mathematical closure flawlessly defines the universe of
\textbf{Homogeneous Tuples} (strictly sized n-dimensional vectors) of
exact, fixed size \(N\).

\paragraph{The Algebra of Linear Functors (No Product
needed!)}\label{the-algebra-of-linear-functors-no-product-needed}

What if we perform the exact opposite mathematical experiment: we
entirely throw away **Product (*)\textbf{ and the }Identity (1)\textbf{
multiplicative atom, and strictly restrict our universe to }Sum
(+)\textbf{, }Compose (\(\circ\))\textbf{, and the }Proxy (0)**
zero-atom?

Because we discarded \texttt{Product}, we mathematically banned
``Pairing''. We physically cannot possess a single data constructor that
holds more than one generic \texttt{a} simultaneously. No Tuples
\texttt{(a,\ a)}, no standard \texttt{Tree} nodes holding branches of
\texttt{a} and \texttt{a}.

If we generate structures strictly using \texttt{Sum} and
\texttt{Compose} over \texttt{Identity} (\(x\)), we generate
mathematically perfect \textbf{Linear Polynomials}: * \texttt{Identity}
is \(x\) * \texttt{Sum\ Identity\ Identity} is \(x + x = 2x\) *
\texttt{Compose\ (Sum\ Identity\ Identity)\ (Sum\ Identity\ Identity)}
is \((2x) \circ (2x) = 4x\)

This strictly mathematical closure flawlessly defines the universe of
\textbf{Linear Choices}. Every single data structure built in this
universe represents an initial choice between exactly \(N\) distinct
pathways (like an Enum), but whichever specific path you take, you will
blindly hold exactly \textbf{one} (or zero) elements of type \texttt{a}
at the absolute end of the path!

\begin{quote}
{[}!NOTE{]} \textbf{What if you add \texttt{Fix} to this?} If you
attempt to apply the \textbf{Fixed Point} (\texttt{Fix}) combinator to a
Product calculus (e.g., \texttt{data\ StreamF\ a\ r\ =\ Cons\ a\ r}),
because you mathematically banned \texttt{Sum}, you physically cannot
formulate a structural base case (like \texttt{Nil} or \texttt{Proxy})
to escape the recursion! This uniquely generates exactly one extreme
type of mathematically pure structure: \textbf{Infinite Streams} (data
structures that literally never end).
\end{quote}

\paragraph{\texorpdfstring{Deriving \texttt{Maybe} from the Absolute
Minimals}{Deriving Maybe from the Absolute Minimals}}\label{deriving-maybe-from-the-absolute-minimals}

Using this algebraic closure, we can mathematically derive the standard
\texttt{Maybe} type purely from our foundational building blocks.

Algebraically, \texttt{Maybe} is practically identical to the structural
\textbf{Sum} of our two absolute minimals (\texttt{Proxy} and
\texttt{Identity}): \texttt{Maybe\ a\ \ ≅\ \ Sum\ Proxy\ Identity\ a}

\begin{itemize}
\tightlist
\item
  \texttt{Nothing} corresponds exactly to the uninhabited
  \texttt{InL\ (Proxy)}.
\item
  \texttt{Just\ x} corresponds exactly to the populated
  \texttt{InR\ (Identity\ x)}.
\end{itemize}

Because we mathematically proved that \texttt{Proxy} flawlessly yields
\texttt{mempty}, \texttt{Identity} flawlessly yields \texttt{f\ x}, and
\texttt{Sum} rightfully delegates to the provided branch, our abstract
algebra perfectly mirrors the definitive \texttt{Foldable} instance for
\texttt{Maybe}:

\begin{Shaded}
\begin{Highlighting}[]
\KeywordTok{instance} \DataTypeTok{Foldable} \DataTypeTok{Maybe} \KeywordTok{where}
    \FunctionTok{foldMap}\NormalTok{ \_ }\DataTypeTok{Nothing}  \OtherTok{=} \FunctionTok{mempty}
    \FunctionTok{foldMap}\NormalTok{ f (}\DataTypeTok{Just}\NormalTok{ x) }\OtherTok{=}\NormalTok{ f x}
\end{Highlighting}
\end{Shaded}

\paragraph{\texorpdfstring{5. The Fixed Point (\texttt{Fix}) and
Conjuring
\texttt{List}}{5. The Fixed Point (Fix) and Conjuring List}}\label{the-fixed-point-fix-and-conjuring-list}

To algebraically conjure an infinitely spanning recursive structure like
\texttt{List}, simple Sums and Products of our minimal atoms
(\texttt{Proxy} and \texttt{Identity}) are not enough. We must tie the
recursive knot using the \textbf{Fixed Point} Combinator (\texttt{Fix}).

A List mathematically represents the explicit algebraic polynomial
\texttt{L(a)\ =\ 1\ +\ a\ *\ L(a)}. Using our combinators, we can
express a single non-recursive structural layer of this as a base
Functor (often called \texttt{ListF}):

\begin{Shaded}
\begin{Highlighting}[]
\KeywordTok{data} \DataTypeTok{ListF}\NormalTok{ a r }\OtherTok{=} \DataTypeTok{Nil} \OperatorTok{|} \DataTypeTok{Cons}\NormalTok{ a r   }\CommentTok{{-}{-} \textquotesingle{}r\textquotesingle{} is the recursion parameter}
\CommentTok{{-}{-} Algebraically: Sum Proxy (Product Identity r)}
\end{Highlighting}
\end{Shaded}

\begin{quote}
{[}!NOTE{]} \textbf{Which Bifoldable generates a List?} The structure
\texttt{ListF\ a\ r} physically takes two type parameters, meaning it
mathematically forms a \textbf{Bifoldable} (and Bifunctor) rather than a
simple Foldable! Precisely, \texttt{ListF} is structurally isomorphic to
the Bifunctor/Bifoldable composition \texttt{Either\ ()\ (a,\ r)}. It
uses the \texttt{Either} Bifoldable for the sum (branching the choice
between \texttt{Nil} and \texttt{Cons}) and the \texttt{(,)} Pair
Bifoldable for the product (holding the \texttt{a} and the \texttt{r}).
By mapping over both its left parameter \texttt{a} and right parameter
\texttt{r} via the \texttt{bifoldMap} operation, we mathematically
prepare the perfect aggregation foundational layer.
\end{quote}

To permanently lock this into an infinite recursive \texttt{List\ a}, we
rigidly apply the type-level \texttt{Fix} combinator (which
mathematically plugs the entire structure back into its own \texttt{r}
parameter infinitely): \texttt{List\ a\ ≅\ Fix\ (ListF\ a)}

When you fold over a \texttt{Fix} structure, what you are essentially
executing is a mathematically pure \textbf{Catamorphism}. The recursion
simply repeatedly delegates strictly to the \texttt{Foldable} instance
of the inner \texttt{Sum} and \texttt{Product} layers. This drills down
the tree until it violently hits the \texttt{Proxy} base case
(\texttt{Nil}), which seamlessly evaluates to \texttt{mempty}, and then
elegantly recursively bubbles all the \texttt{mappend} operations back
up the execution tree to form a single value!

\begin{quote}
{[}!NOTE{]} \textbf{Wait, why didn't we explicitly do this for
Functors?} Actually, we completely \emph{did}, just without using the
explicit algebraic terminology! Functors mathematically possess the
exact same compositional closures.

Under the hood, Functor Sums and Products are quite literally
implemented using the foundational \textbf{Bifunctors} we saw in Chapter
1: \texttt{Either} (\(+\)) and \texttt{(,)} (\(\times\)). The
\texttt{Data.Functor.Sum} type physically just wraps the Bifunctor
\texttt{Either\ (f\ a)\ (g\ a)}!

And what about \texttt{Compose}? Can you compose Functors into a
Bifunctor? Yes! Haskell explicitly provides the \texttt{Biff} operator
in \texttt{Data.Bifunctor.Biff}:
\(Biff \ p \ f \ g \ a \ b = p \ (f \ a) \ (g \ b)\). \texttt{Biff}
mathematically proves that if you take a base Bifunctor (\(p\)) and
substitute two Functors into its parameters (\(f, g\)), the result is
mathematically guaranteed to be a perfectly lawful Bifunctor!

Furthermore, recursive Functors like \texttt{List} are algebraically
constructed by applying \texttt{Fix} to a base Functor (just as we
proved). Everything we just mathematically proved for \texttt{Foldable}
applies universally and flawlessly to \texttt{Functor} and
\texttt{Bifunctor}!
\end{quote}

\paragraph{Do the Algebraic Combinators Require New
Laws?}\label{do-the-algebraic-combinators-require-new-laws}

For these compound structures (\texttt{Sum}, \texttt{Product},
\texttt{Compose}, and \texttt{Fix}) to be mathematically valid
Foldables, they must inherently respect the Foldable consistency laws
(e.g., \texttt{foldMap\ f\ ==\ foldr\ (mappend\ .\ f)\ mempty}).

Do we need to explicitly prove new laws for them, like we did for
associativity and identity in Bifunctors?

No! We receive a massive mathematical freebie. Because our combinators
are defined \emph{strictly} using the underlying Base \texttt{foldMap}
operations and the Monoid \texttt{\textless{}\textgreater{}} operator,
\textbf{the abstract algebra automatically guarantees the compound
structures obey the laws}. Assuming the base atoms (like \texttt{Proxy}
and \texttt{Identity}) are valid, the rigid associativity of the Monoid
(\texttt{\textless{}\textgreater{}}) flawlessly ensures that whether you
fold completely sequentially (\texttt{foldr}), or smash nested
structures together hierarchically (\texttt{Compose} / \texttt{Product}
/ \texttt{Fix}), the final aggregated value will unequivocally evaluate
to the exact same monoidal mathematical truth!

\subsection{Chapter 3: Traversable (Effectful
Folding)}\label{chapter-3-traversable-effectful-folding}

If \texttt{Functor} is about \textbf{Shape Preservation} (mapping
functions over data without altering the container) and
\texttt{Foldable} is about \textbf{Aggregation} (destroying the shape to
fold its elements into a single Monoid), then \texttt{Traversable}
represents the final pillar of this mathematical trinity:
\textbf{Effectful Sequencing}.

\texttt{Traversable} allows you to navigate the shape from left to right
while performing an \texttt{Applicative} effect on every element, and
finally sequence all those effects into a single overarching context
that rebuilds the exact original shape inside!

\subsubsection{Section 8.1: What is
Traversable?}\label{section-8.1-what-is-traversable}

The foundational method of \texttt{Traversable} is \texttt{traverse}:

\begin{Shaded}
\begin{Highlighting}[]
\KeywordTok{class}\NormalTok{ (}\DataTypeTok{Functor}\NormalTok{ t, }\DataTypeTok{Foldable}\NormalTok{ t) }\OtherTok{=\textgreater{}} \DataTypeTok{Traversable}\NormalTok{ t }\KeywordTok{where}
\OtherTok{    traverse ::} \DataTypeTok{Applicative}\NormalTok{ f }\OtherTok{=\textgreater{}}\NormalTok{ (a }\OtherTok{{-}\textgreater{}}\NormalTok{ f b) }\OtherTok{{-}\textgreater{}}\NormalTok{ t a }\OtherTok{{-}\textgreater{}}\NormalTok{ f (t b)}
\end{Highlighting}
\end{Shaded}

Look closely at the signature: it requires \texttt{Functor} and
\texttt{Foldable} as prerequisites. This is because traversing
fundamentally requires walking through the entire structure (like
\texttt{Foldable}) and physically rebuilding the identical shape of the
container at the end (like \texttt{Functor}).

While \texttt{Foldable} aggressively tears down the structure using
\texttt{\textless{}\textgreater{}} into a single \texttt{mempty},
\texttt{Traversable} delicately sequences evaluations using
\texttt{\textless{}*\textgreater{}} to yield an \texttt{Applicative}
effect \texttt{f} that flawlessly holds a rebuilt structure
\texttt{t\ b}.

\begin{quote}
{[}!NOTE{]} \textbf{Wait, didn't we say \texttt{foldMap} was just
\texttt{traverse} hiding inside the \texttt{Const} Functor?} Yes! Now
that we've defined \texttt{traverse}, we can prove it. If we choose our
\texttt{Applicative} to be the \texttt{Const} Functor (which ignores the
value update and merely accumulates the monoidal tags via
\texttt{\textless{}*\textgreater{}}!), \texttt{traverse} functionally
devolves directly into \texttt{foldMap}!

\begin{Shaded}
\begin{Highlighting}[]
\CommentTok{{-}{-} An elegant, universal way to calculate length across ANY Foldable!}
\OtherTok{len ::} \DataTypeTok{Foldable}\NormalTok{ t }\OtherTok{=\textgreater{}}\NormalTok{ t a }\OtherTok{{-}\textgreater{}} \DataTypeTok{Int}
\NormalTok{len xs }\OtherTok{=}\NormalTok{ getSum (}\FunctionTok{foldMap}\NormalTok{ (}\FunctionTok{const}\NormalTok{ (}\DataTypeTok{Sum} \DecValTok{1}\NormalTok{)) xs)}

\CommentTok{{-}{-} Equivalently, via a destructive Traversable using the \textasciigrave{}Const\textasciigrave{} Functor:}
\OtherTok{lenT ::} \DataTypeTok{Traversable}\NormalTok{ t }\OtherTok{=\textgreater{}}\NormalTok{ t a }\OtherTok{{-}\textgreater{}} \DataTypeTok{Int}
\NormalTok{lenT xs }\OtherTok{=}\NormalTok{ getSum (getConst (}\FunctionTok{traverse}\NormalTok{ (}\DataTypeTok{Const} \OperatorTok{.} \FunctionTok{const}\NormalTok{ (}\DataTypeTok{Sum} \DecValTok{1}\NormalTok{)) xs))}
\end{Highlighting}
\end{Shaded}
\end{quote}

The most beautiful revelation is that \texttt{Traversable} shares
\textbf{the exact same polynomial algebra, atoms, and operations} that
we rigorously defined for Functors and Foldables. Let's prove it by
reconstructing \texttt{Traversable} from the mathematical substrate up.

\subsubsection{Section 8.2: The Absolute Minimum Traversable
Atoms}\label{section-8.2-the-absolute-minimum-traversable-atoms}

Because Traversable relies on the same polynomial closure, we begin with
our trusted atoms:

\paragraph{\texorpdfstring{1. The Empty Sequence: \texttt{Proxy} (The
``0''
Atom)}{1. The Empty Sequence: Proxy (The ``0'' Atom)}}\label{the-empty-sequence-proxy-the-0-atom}

If there is absolutely no data of type \texttt{a} inside the structure,
what happens when we attempt to traverse it to sequence its effects?
There are no effects to sequence! The mathematical forced hand simply
lifts the empty box intact directly into the \texttt{Applicative}
context using \texttt{pure}.

\begin{Shaded}
\begin{Highlighting}[]
\KeywordTok{instance} \DataTypeTok{Traversable} \DataTypeTok{Proxy} \KeywordTok{where}
    \CommentTok{{-}{-} traverse :: Applicative f =\textgreater{} (a {-}\textgreater{} f b) {-}\textgreater{} Proxy a {-}\textgreater{} f (Proxy b)}
    \FunctionTok{traverse}\NormalTok{ \_ }\DataTypeTok{Proxy} \OtherTok{=} \FunctionTok{pure} \DataTypeTok{Proxy}
\end{Highlighting}
\end{Shaded}

This is mathematically absolute: navigating a zero-length sequence
requires zero sequenced operations, resulting purely in the vacuous
success of the overarching effect!

\paragraph{\texorpdfstring{2. The Single Effect: \texttt{Identity} (The
``1''
Atom)}{2. The Single Effect: Identity (The ``1'' Atom)}}\label{the-single-effect-identity-the-1-atom}

If there is exactly one generic element, we must unconditionally
evaluate our effect on it.

\begin{Shaded}
\begin{Highlighting}[]
\KeywordTok{instance} \DataTypeTok{Traversable} \DataTypeTok{Identity} \KeywordTok{where}
    \CommentTok{{-}{-} traverse :: Applicative f =\textgreater{} (a {-}\textgreater{} f b) {-}\textgreater{} Identity a {-}\textgreater{} f (Identity b)}
    \FunctionTok{traverse}\NormalTok{ f (}\DataTypeTok{Identity}\NormalTok{ x) }\OtherTok{=} \FunctionTok{fmap} \DataTypeTok{Identity}\NormalTok{ (f x) }
\end{Highlighting}
\end{Shaded}

Here, \texttt{f\ x} generates our \texttt{Applicative} effect (e.g., an
\texttt{IO} action or a \texttt{Maybe} computation). We mathematically
map (\texttt{fmap}) the \texttt{Identity} constructor \emph{inside} that
effect to strictly reconstruct our \(x^1\) bound!

\subsubsection{Section 8.3: The Algebra of
Traversables}\label{section-8.3-the-algebra-of-traversables}

Now, let's look at how the categorical binary operations effortlessly
scale into \texttt{Traversable}.

\paragraph{\texorpdfstring{Traversable Sums
(\texttt{f\ +\ g})}{Traversable Sums (f + g)}}\label{traversable-sums-f-g}

Just as \texttt{Either} allowed branching combinations for folds, it
allows us to delegate effectful traversals.

\begin{Shaded}
\begin{Highlighting}[]
\KeywordTok{instance}\NormalTok{ (}\DataTypeTok{Traversable}\NormalTok{ f, }\DataTypeTok{Traversable}\NormalTok{ g) }\OtherTok{=\textgreater{}} \DataTypeTok{Traversable}\NormalTok{ (}\DataTypeTok{Sum}\NormalTok{ f g) }\KeywordTok{where}
    \FunctionTok{traverse}\NormalTok{ fn (}\DataTypeTok{InL}\NormalTok{ fa) }\OtherTok{=} \FunctionTok{fmap} \DataTypeTok{InL}\NormalTok{ (}\FunctionTok{traverse}\NormalTok{ fn fa)}
    \FunctionTok{traverse}\NormalTok{ fn (}\DataTypeTok{InR}\NormalTok{ ga) }\OtherTok{=} \FunctionTok{fmap} \DataTypeTok{InR}\NormalTok{ (}\FunctionTok{traverse}\NormalTok{ fn ga)}
\end{Highlighting}
\end{Shaded}

If the execution pathway ventures into the \texttt{InL} branch, we
mathematically sequence the Left structure. Since
\texttt{traverse\ fn\ fa} perfectly yields an \texttt{f\ (fa\ b)}, we
\texttt{fmap} the \texttt{InL} boundary tag to cleanly rebuild the
correct Sum geometry. The exact same operation handles the Right branch
symmetrically!

\paragraph{\texorpdfstring{Traversable Products
(\texttt{f\ *\ g})}{Traversable Products (f * g)}}\label{traversable-products-f-g}

A Product possesses both structures simultaneously. To traverse a
\texttt{Product} from left to right, we unequivocally must: 1. Traverse
the first structure. 2. Traverse the second structure. 3. Bundle the two
executing effects together so their combined structures are perfectly
preserved inside the resulting \texttt{Applicative} effect.

\begin{Shaded}
\begin{Highlighting}[]
\KeywordTok{instance}\NormalTok{ (}\DataTypeTok{Traversable}\NormalTok{ f, }\DataTypeTok{Traversable}\NormalTok{ g) }\OtherTok{=\textgreater{}} \DataTypeTok{Traversable}\NormalTok{ (}\DataTypeTok{Product}\NormalTok{ f g) }\KeywordTok{where}
    \FunctionTok{traverse}\NormalTok{ fn (}\DataTypeTok{Pair}\NormalTok{ fa ga) }\OtherTok{=} 
\NormalTok{        liftA2 }\DataTypeTok{Pair}\NormalTok{ (}\FunctionTok{traverse}\NormalTok{ fn fa) (}\FunctionTok{traverse}\NormalTok{ fn ga)}
        \CommentTok{{-}{-} Equivalent to: Pair \textless{}$\textgreater{} traverse fn fa \textless{}*\textgreater{} traverse fn ga}
\end{Highlighting}
\end{Shaded}

Here, we see the profound elegance of the Applicative
\texttt{\textless{}*\textgreater{}} operator. It natively handles the
exact simultaneous product combination we require, seamlessly executing
the effects of \texttt{fa} before \texttt{ga} while merging their
resultant shapes inside the \texttt{Pair} constructor!

\paragraph{\texorpdfstring{Traversable Composition
(\texttt{f\ ∘\ g})}{Traversable Composition (f ∘ g)}}\label{traversable-composition-f-g}

Just as we nested \texttt{Foldable} loops, we can strictly nest
\texttt{Traversable} effects.

\begin{Shaded}
\begin{Highlighting}[]
\KeywordTok{instance}\NormalTok{ (}\DataTypeTok{Traversable}\NormalTok{ f, }\DataTypeTok{Traversable}\NormalTok{ g) }\OtherTok{=\textgreater{}} \DataTypeTok{Traversable}\NormalTok{ (}\DataTypeTok{Compose}\NormalTok{ f g) }\KeywordTok{where}
    \FunctionTok{traverse}\NormalTok{ fn (}\DataTypeTok{Compose}\NormalTok{ fga) }\OtherTok{=} 
        \FunctionTok{fmap} \DataTypeTok{Compose}\NormalTok{ (}\FunctionTok{traverse}\NormalTok{ (}\FunctionTok{traverse}\NormalTok{ fn) fga)}
\end{Highlighting}
\end{Shaded}

To traverse the outer nested layer, what function do we apply? Our inner
traversing function \texttt{traverse\ fn}! The inner sequence resolves
its Applicative effects into the outer sequence, compounding them
perfectly.

\subsubsection{The Grand Architectural
Synthesis}\label{the-grand-architectural-synthesis}

Look deeply at what we just proved algebraically. By using strictly: 1.
The \textbf{Zero Atom} (\texttt{Proxy} / \(0\)) 2. The \textbf{One Atom}
(\texttt{Identity} / \(1\)) 3. \textbf{Sums} (\texttt{Either} / \(+\))
4. \textbf{Products} (\texttt{(,)} / \(\times\))

We systematically programmed flawless implementations for
\textbf{Functor, Foldable, and Traversable}.

This concludes a magnificent piece of abstract mathematical geometry:
\textbf{Because these Typeclasses are perfectly closed over the
polynomial operators, any Algebraic Data Type (ADT) formulated using
Sums, Products, Zeros, and Ones is rigidly mathematically guaranteed to
be a valid Functor, Foldable, AND Traversable!}

This algebraic theorem is exactly what powers the
\texttt{DeriveFunctor}, \texttt{DeriveFoldable}, and
\texttt{DeriveTraversable} compiler extensions. The Haskell compiler
does not guess; it systematically parses your data type as a structural
mathematical polynomial and algebraically applies these exact
foundational atoms and combinators to write the canonical,
mathematically flawless instances for you.

\begin{center}\rule{0.5\linewidth}{0.5pt}\end{center}

\subsection{Chapter 4: Applicative (Context
Aggregation)}\label{chapter-4-applicative-context-aggregation}

Now we step up in power. An \texttt{Applicative} is a Functor equipped
with two new powers: \texttt{pure} (to lift values) and
\texttt{\textless{}*\textgreater{}} (to lift application).

\subsubsection{Section 2.1: The Applicative
Atoms}\label{section-2.1-the-applicative-atoms}

Let's see how our atomic structures ``upgrade'' to this new level.

\paragraph{\texorpdfstring{1. \texttt{Proxy}}{1. Proxy}}\label{proxy}

\begin{Shaded}
\begin{Highlighting}[]
\KeywordTok{instance} \DataTypeTok{Applicative} \DataTypeTok{Proxy} \KeywordTok{where}
    \FunctionTok{pure}\NormalTok{ \_ }\OtherTok{=} \DataTypeTok{Proxy}
    \DataTypeTok{Proxy} \OperatorTok{\textless{}*\textgreater{}} \DataTypeTok{Proxy} \OtherTok{=} \DataTypeTok{Proxy}
\end{Highlighting}
\end{Shaded}

\textbf{The ``Why''}: Our hands are tied. \texttt{pure} gives us an
\texttt{a}, which we must discard (as \texttt{Proxy} holds no data).
\texttt{\textless{}*\textgreater{}} combines two empty boxes into one.

\paragraph{\texorpdfstring{2. \texttt{Const\ r} (The Monoid
Requirement)}{2. Const r (The Monoid Requirement)}}\label{const-r-the-monoid-requirement}

This is the most critical upgrade in the minimal universe.

\begin{Shaded}
\begin{Highlighting}[]
\KeywordTok{instance} \DataTypeTok{Monoid}\NormalTok{ r }\OtherTok{=\textgreater{}} \DataTypeTok{Applicative}\NormalTok{ (}\DataTypeTok{Const}\NormalTok{ r) }\KeywordTok{where}
    \FunctionTok{pure}\NormalTok{ \_ }\OtherTok{=} \DataTypeTok{Const} \FunctionTok{mempty}
    \DataTypeTok{Const}\NormalTok{ r1 }\OperatorTok{\textless{}*\textgreater{}} \DataTypeTok{Const}\NormalTok{ r2 }\OtherTok{=} \DataTypeTok{Const}\NormalTok{ (r1 }\OtherTok{\textasciigrave{}mappend\textasciigrave{}}\NormalTok{ r2)}
\end{Highlighting}
\end{Shaded}

\textbf{The ``Why''}: * \texttt{pure} requires us to produce an
\texttt{r} out of nothing. We must use the \textbf{Identity element}
(\texttt{mempty}). * \texttt{\textless{}*\textgreater{}} gives us two
\texttt{r} values and needs one result. We must use the \textbf{Binary
operation} (\texttt{mappend}). This precisely defines why \texttt{Const}
requires its context to be a \texttt{Monoid} to achieve Applicative
status.

\paragraph{\texorpdfstring{3.
\texttt{Identity}}{3. Identity}}\label{identity}

\begin{Shaded}
\begin{Highlighting}[]
\KeywordTok{instance} \DataTypeTok{Applicative} \DataTypeTok{Identity} \KeywordTok{where}
    \FunctionTok{pure}\NormalTok{ x }\OtherTok{=} \DataTypeTok{Identity}\NormalTok{ x}
    \DataTypeTok{Identity}\NormalTok{ f }\OperatorTok{\textless{}*\textgreater{}} \DataTypeTok{Identity}\NormalTok{ x }\OtherTok{=} \DataTypeTok{Identity}\NormalTok{ (f x)}
\end{Highlighting}
\end{Shaded}

\textbf{The ``Why''}: Trivial application. We unwrap, apply, and rewrap.

\subsubsection{\texorpdfstring{Section 2.2: The Applicative Analog to
foldMap
(\texttt{traverse})}{Section 2.2: The Applicative Analog to foldMap (traverse)}}\label{section-2.2-the-applicative-analog-to-foldmap-traverse}

When working with \texttt{Foldable}, we saw how \texttt{foldMap} allows
us to elegantly collapse a structure by mapping each element to a
\texttt{Monoid} and combining them. With \texttt{Applicative}, we gain a
structurally analogous, but strictly more powerful operation from the
\texttt{Traversable} class: \texttt{traverse}.

\begin{Shaded}
\begin{Highlighting}[]
\CommentTok{{-}{-} The type signatures conceptually mirror each other beautifully}
\FunctionTok{foldMap}\OtherTok{  ::}\NormalTok{ (}\DataTypeTok{Foldable}\NormalTok{ t,    }\DataTypeTok{Monoid}\NormalTok{ m)      }\OtherTok{=\textgreater{}}\NormalTok{ (a }\OtherTok{{-}\textgreater{}}\NormalTok{ m)   }\OtherTok{{-}\textgreater{}}\NormalTok{ t a }\OtherTok{{-}\textgreater{}}\NormalTok{ m}
\FunctionTok{traverse}\OtherTok{ ::}\NormalTok{ (}\DataTypeTok{Traversable}\NormalTok{ t, }\DataTypeTok{Applicative}\NormalTok{ f) }\OtherTok{=\textgreater{}}\NormalTok{ (a }\OtherTok{{-}\textgreater{}}\NormalTok{ f b) }\OtherTok{{-}\textgreater{}}\NormalTok{ t a }\OtherTok{{-}\textgreater{}}\NormalTok{ f (t b)}
\end{Highlighting}
\end{Shaded}

Just as \texttt{foldMap} walks a structure and aggregates values using a
\texttt{Monoid}'s \texttt{mappend} and \texttt{mempty},
\texttt{traverse} walks a structure and aggregates \emph{computational
effects} using an \texttt{Applicative}'s
\texttt{\textless{}*\textgreater{}} and \texttt{pure}.

They are fundamentally so intimately related that we can completely and
perfectly recreate \texttt{foldMap} using exactly our minimal
\texttt{Const} functor! Recall from Section 2.1 that \texttt{Const\ m}
acts as an \texttt{Applicative} precisely when its context \texttt{m} is
a \texttt{Monoid}. Its \texttt{\textless{}*\textgreater{}} behaves
exactly like \texttt{mappend}, and its \texttt{pure} behaves exactly
like \texttt{mempty}.

By passing \texttt{Const} to \texttt{traverse}, we essentially trick the
\texttt{Traversable} into performing a \texttt{Foldable} operation. We
tell it to discard the structural reconstruction (since \texttt{Const}
holds no computational result) and solely accumulate the internal
\texttt{Monoid} state:

\begin{Shaded}
\begin{Highlighting}[]
\KeywordTok{import} \DataTypeTok{Data.Functor.Const}

\CommentTok{{-}{-} Emulating foldMap perfectly using traverse and Const}
\OtherTok{foldMapTraverse ::}\NormalTok{ (}\DataTypeTok{Traversable}\NormalTok{ t, }\DataTypeTok{Monoid}\NormalTok{ m) }\OtherTok{=\textgreater{}}\NormalTok{ (a }\OtherTok{{-}\textgreater{}}\NormalTok{ m) }\OtherTok{{-}\textgreater{}}\NormalTok{ t a }\OtherTok{{-}\textgreater{}}\NormalTok{ m}
\NormalTok{foldMapTraverse f xs }\OtherTok{=}\NormalTok{ getConst }\OperatorTok{$} \FunctionTok{traverse}\NormalTok{ (}\DataTypeTok{Const} \OperatorTok{.}\NormalTok{ f) xs}
\end{Highlighting}
\end{Shaded}

This mathematical elegance proves that folding is essentially a special
case of traversal, where the ``effect'' being sequenced is simply the
accumulation of a Monoid. It perfectly bridges the worlds of
\texttt{Monoid} and \texttt{Applicative} using our minimal atom,
\texttt{Const}.

\subsubsection{Section 2.3: Automated Law
Testing}\label{section-2.3-automated-law-testing}

Just as with Functors, we can verify our Applicative instances using
\texttt{tasty-checkers}. This is where the library truly shines, as the
number of Applicative laws (Identity, Homomorphism, Interchange, and
Composition) is significantly higher:

\begin{Shaded}
\begin{Highlighting}[]
  \CommentTok{{-}{-} Automatically tests all Applicative laws}
\NormalTok{  testBatch (applicative (}\FunctionTok{undefined}\OtherTok{ ::} \DataTypeTok{Maybe}\NormalTok{ (}\DataTypeTok{Int}\NormalTok{, }\DataTypeTok{String}\NormalTok{, }\DataTypeTok{Int}\NormalTok{)))}
\end{Highlighting}
\end{Shaded}

\begin{center}\rule{0.5\linewidth}{0.5pt}\end{center}

\subsection{Chapter 5: Monad (Effectful
Sequencing)}\label{chapter-5-monad-effectful-sequencing}

The \texttt{Monad} adds the power of \textbf{Context-Dependent
Sequencing} via \texttt{bind} (\texttt{\textgreater{}\textgreater{}=})
or \texttt{join}.

\subsubsection{Section 3.1: The Final
Upgrades}\label{section-3.1-the-final-upgrades}

\paragraph{\texorpdfstring{1. \texttt{Proxy}}{1. Proxy}}\label{proxy-1}

\begin{Shaded}
\begin{Highlighting}[]
\KeywordTok{instance} \DataTypeTok{Monad} \DataTypeTok{Proxy} \KeywordTok{where}
    \DataTypeTok{Proxy} \OperatorTok{\textgreater{}\textgreater{}=}\NormalTok{ \_ }\OtherTok{=} \DataTypeTok{Proxy}
\end{Highlighting}
\end{Shaded}

Flattening an empty box inside an empty box still yields an empty box.

\paragraph{\texorpdfstring{2.
\texttt{Identity}}{2. Identity}}\label{identity-1}

\begin{Shaded}
\begin{Highlighting}[]
\KeywordTok{instance} \DataTypeTok{Monad} \DataTypeTok{Identity} \KeywordTok{where}
    \DataTypeTok{Identity}\NormalTok{ x }\OperatorTok{\textgreater{}\textgreater{}=}\NormalTok{ f }\OtherTok{=}\NormalTok{ f x}
\end{Highlighting}
\end{Shaded}

Pure function application.

\paragraph{\texorpdfstring{3. \texttt{Const\ r} (The Monad
Barrier)}{3. Const r (The Monad Barrier)}}\label{const-r-the-monad-barrier}

\textbf{Crucially, \texttt{Const\ r} cannot be a Monad.}

\begin{Shaded}
\begin{Highlighting}[]
\OtherTok{(\textgreater{}\textgreater{}=) ::} \DataTypeTok{Const}\NormalTok{ r a }\OtherTok{{-}\textgreater{}}\NormalTok{ (a }\OtherTok{{-}\textgreater{}} \DataTypeTok{Const}\NormalTok{ r b) }\OtherTok{{-}\textgreater{}} \DataTypeTok{Const}\NormalTok{ r b}
\end{Highlighting}
\end{Shaded}

Because \texttt{Const} contains no \texttt{a}, we can never execute the
function \texttt{(a\ -\textgreater{}\ Const\ r\ b)}. We completely lose
whatever \texttt{r} value the function \emph{would} have produced,
violating the \textbf{Left Identity law}
(\texttt{pure\ a\ \textgreater{}\textgreater{}=\ f\ ==\ f\ a}). The
evolution stops here.

\subsubsection{Section 3.2: Automated Law
Testing}\label{section-3.2-automated-law-testing}

Finally, we can verify our Monad instances (Left Identity, Right
Identity, and Associativity) with a single check:

\begin{Shaded}
\begin{Highlighting}[]
  \CommentTok{{-}{-} Automatically tests all Monad laws}
\NormalTok{  testBatch (monad (}\FunctionTok{undefined}\OtherTok{ ::} \DataTypeTok{Maybe}\NormalTok{ (}\DataTypeTok{Int}\NormalTok{, }\DataTypeTok{String}\NormalTok{, }\DataTypeTok{Int}\NormalTok{)))}
\end{Highlighting}
\end{Shaded}

\begin{center}\rule{0.5\linewidth}{0.5pt}\end{center}

\subsection{Conclusion: The Tale of Three
Minimals}\label{conclusion-the-tale-of-three-minimals}

By starting from these absolute minimal examples, the ``magic''
evaporates, leaving the elegant logic of types and the algebraic
discovery of everything from \texttt{Maybe} to \texttt{List}.

\begin{center}\rule{0.5\linewidth}{0.5pt}\end{center}

\begin{center}\rule{0.5\linewidth}{0.5pt}\end{center}

\begin{center}\rule{0.5\linewidth}{0.5pt}\end{center}

\section{Part 5: The Functor Monoids}\label{part-5-the-functor-monoids}

This document captures a profound mathematical exploration into the
foundations of Algebraic Data Types (ADTs). Instead of looking at simple
types, we will investigate the \textbf{Category of Endofunctors}
(\texttt{Type\ -\textgreater{}\ Type}).

The goal is to discover the irreducible minimal generating set of the
Functor universe.

\subsection{Chapter 1: The Functor Monoid (The True
Engine)}\label{chapter-1-the-functor-monoid-the-true-engine}

Before we can build infinitely large data structures, we need to
understand how two functors mathematically combine.

If we have a universe of Functors, a mathematically perfect way to
define a ``structural glue'' is to define a \textbf{Monoid over
Functors}. A Monoid requires exactly two things: an Identity Element
(the 0-ary base case) and a Binary Combinator (the 2-ary joining
operator).

We can represent this concept as a type-level record or class:

\begin{Shaded}
\begin{Highlighting}[]
\CommentTok{{-}{-} | A Monoid over the Functor Category}
\KeywordTok{class} \DataTypeTok{FunctorMonoid}\NormalTok{ (}\OtherTok{m ::}\NormalTok{ [}\DataTypeTok{Type} \OtherTok{{-}\textgreater{}} \DataTypeTok{Type}\NormalTok{] }\OtherTok{{-}\textgreater{}} \OperatorTok{*} \OtherTok{{-}\textgreater{}} \OperatorTok{*}\NormalTok{) }\KeywordTok{where}
    \CommentTok{{-}{-} The 0{-}ary Base Case (Identity Element)}
    \CommentTok{{-}{-} Must be a mathematically valid Functor.}
    \KeywordTok{type} \DataTypeTok{Atom}\OtherTok{ m ::} \OperatorTok{*} \OtherTok{{-}\textgreater{}} \OperatorTok{*}
    
    \CommentTok{{-}{-} The 2{-}ary Combinator}
    \CommentTok{{-}{-} Given two valid Functors \textasciigrave{}f\textasciigrave{} and \textasciigrave{}g\textasciigrave{}, must produce a valid Functor.}
    \KeywordTok{type} \DataTypeTok{Bin}\OtherTok{ m ::}\NormalTok{ (}\OperatorTok{*} \OtherTok{{-}\textgreater{}} \OperatorTok{*}\NormalTok{) }\OtherTok{{-}\textgreater{}}\NormalTok{ (}\OperatorTok{*} \OtherTok{{-}\textgreater{}} \OperatorTok{*}\NormalTok{) }\OtherTok{{-}\textgreater{}}\NormalTok{ (}\OperatorTok{*} \OtherTok{{-}\textgreater{}} \OperatorTok{*}\NormalTok{)}
\end{Highlighting}
\end{Shaded}

Any record that validly implements this class completely defines a
structural family of data types. The associated mathematical laws
dictate that \texttt{Bin\ m\ (Atom\ m)\ f} must be structurally
isomorphic to \texttt{f} (the Identity Law).

\subsubsection{Section 1.1: The Minimal
Generators}\label{section-1.1-the-minimal-generators}

If we look at the universe of Functors, what are the fundamental
building blocks that can satisfy this Monoid, and how far do they get
us?

There are bounds to the universe we can manipulate.

\subsubsection{A. The Dummy Bounds (The Constant
Monoids)}\label{a.-the-dummy-bounds-the-constant-monoids}

The absolute simplest monoids are those that completely ignore the
\texttt{Bin} operation, throwing away all data combinations and always
returning the 0-ary \texttt{Atom}.

\begin{itemize}
\tightlist
\item
  \textbf{Constant Zero}: \texttt{Atom\ =\ Zero},
  \texttt{Bin\ f\ g\ =\ Zero}. Folding this defines the
  \textbf{\texttt{NaryZeroF}} glue (The mathematical Black Hole, which
  holds 0 constructors).
\item
  \textbf{Constant Proxy}: \texttt{Atom\ =\ Proxy},
  \texttt{Bin\ f\ g\ =\ Proxy}. Folding this defines the
  \textbf{\texttt{NaryProxyF}} glue (The mathematical Empty Box, holding
  1 empty constructor).
\end{itemize}

Both of these represent flat closures.

\subsubsection{\texorpdfstring{B. Semigroups vs.~Monoids (The
Destructive Fakes \texttt{First} and
\texttt{Second})}{B. Semigroups vs.~Monoids (The Destructive Fakes First and Second)}}\label{b.-semigroups-vs.-monoids-the-destructive-fakes-first-and-second}

What if we drop the \texttt{Atom} requirement entirely and only keep the
\texttt{Bin} operator? If we just enforce an associativity law on
\texttt{Bin}, we have defined a \textbf{Semigroup over Functors}.

With just a Semigroup, you can define combinators like \texttt{First}
(where \texttt{type\ Bin\ f\ g\ =\ f}) or \texttt{Second}. They satisfy
the
\texttt{(*\ -\textgreater{}\ *)\ -\textgreater{}\ (*\ -\textgreater{}\ *)\ -\textgreater{}\ (*\ -\textgreater{}\ *)}
kind requirement perfectly.

However, they are mathematically ``incomplete'' to build N-ary
structures. If you use a Semigroup \texttt{Bin} to fold an N-ary list,
you can only fold lists that have \emph{at least one functor in them}
(\(N \ge 1\)).

If you hand a Semigroup an empty list \texttt{{[}{]}} to fold, it
mathematically catastrophically crashes. Because you dropped the
\texttt{Atom} (the identity element), the space is undefined at \(N=0\).

A Monoid explicitly requires an \texttt{Atom} such that
\texttt{Bin\ e\ f\ =\ f} (Left Identity) and \texttt{Bin\ f\ e\ =\ f}
(Right Identity). It is impossible to define an \texttt{Atom} base case
for \texttt{First} because the \texttt{e} would have to magically morph
its type to equal whatever arbitrary \texttt{f} was passed in!

Therefore, \texttt{First} and \texttt{Second} are mathematically valid
\emph{Semigroups}, but they cannot ever be \emph{Monoids}. And because
they have no \(N=0\) case, they cannot generate complete ADT universes!

There are exactly two fundamental active properties we can manipulate:
\textbf{Choice} and \textbf{Conjunction}.

\subsubsection{C. The Minimum Generators for Choice (The Sum
Record)}\label{c.-the-minimum-generators-for-choice-the-sum-record}

If our \texttt{Bin} combinator represents a branching path (an
\texttt{OR} relationship), it is the \textbf{Sum} Functor (evaluating to
\texttt{Either}).

What is the \texttt{Atom} (identity element) for Sum? It must be the
Functor that adds zero choices. If \(X + e = X\), then \(e\)
mathematically must be the \textbf{\texttt{Zero} Functor} (an impossible
state with zero constructors, effectively \texttt{Void}).

\begin{itemize}
\tightlist
\item
  \textbf{Combinator (\texttt{Bin})}: Sum (\texttt{Either})
\item
  \textbf{Atom (\texttt{Atom})}: Zero (\texttt{Void})
\end{itemize}

\subsubsection{D. The Minimum Generators for Conjunction (The Product
Record)}\label{d.-the-minimum-generators-for-conjunction-the-product-record}

If our \texttt{Bin} combinator represents holding data simultaneously
(an \texttt{AND} relationship), it is the \textbf{Product} Functor
(evaluating to a tuple \texttt{(,)}).

What is the \texttt{Atom} (identity element) for Product? It must be the
Functor that adds zero information. If \(X \times e = X\), then \(e\)
mathematically must be the \textbf{\texttt{Proxy} Functor} (a single
stateless constructor \texttt{()}).

\begin{itemize}
\tightlist
\item
  \textbf{Combinator (\texttt{Bin})}: Product \texttt{(,)}
\item
  \textbf{Atom (\texttt{Atom})}: Proxy \texttt{()}
\end{itemize}

\subsubsection{\texorpdfstring{E. The Missing Infinites:
\texttt{Fix}}{E. The Missing Infinites: Fix}}\label{e.-the-missing-infinites-fix}

Sum and Product can generate any finite data shape. However, to generate
infinite or recursive structures (like \texttt{List} or \texttt{Tree}),
we strictly need a new minimal operator that bends a Functor back onto
itself.

\begin{itemize}
\tightlist
\item
  \textbf{The Recursive Minimal}: \texttt{Fix\ f\ =\ In\ (f\ (Fix\ f))}
  By feeding a combination of Sum and Product into \texttt{Fix}, we
  escape the finite universe.
\end{itemize}

\subsubsection{\texorpdfstring{F. The Alien Minimal:
\texttt{(-\textgreater{})}}{F. The Alien Minimal: (-\textgreater)}}\label{f.-the-alien-minimal--}

Computations (functions) cannot be generated by Sums, Products, or Fix.
Exponentiation requires a fundamentally distinct 2-ary minimal operator:
the Arrow \texttt{(-\textgreater{})}.

\subsubsection{\texorpdfstring{G. The Hidden Third Monoid:
\texttt{Compose}}{G. The Hidden Third Monoid: Compose}}\label{g.-the-hidden-third-monoid-compose}

It turns out there is a profound third mathematical Functor Monoid that
exists perfectly symmetrical to Sum and Product: \textbf{Functor
Composition}. Instead of branching (Sum) or pairing (Product), what if
we nest one Functor strictly inside another?

\begin{Shaded}
\begin{Highlighting}[]
\KeywordTok{class} \DataTypeTok{FunctorMonoid}\NormalTok{ (}\OtherTok{m ::}\NormalTok{ (}\OperatorTok{*} \OtherTok{{-}\textgreater{}} \OperatorTok{*}\NormalTok{) }\OtherTok{{-}\textgreater{}}\NormalTok{ (}\OperatorTok{*} \OtherTok{{-}\textgreater{}} \OperatorTok{*}\NormalTok{) }\OtherTok{{-}\textgreater{}} \OperatorTok{*} \OtherTok{{-}\textgreater{}} \OperatorTok{*}\NormalTok{) }\KeywordTok{where}
    \KeywordTok{type} \DataTypeTok{Unit}\OtherTok{ m ::} \OperatorTok{*} \OtherTok{{-}\textgreater{}} \OperatorTok{*}
    \CommentTok{{-}{-} The Combinator is inherently represented by the \textquotesingle{}m\textquotesingle{} parameter!}

\KeywordTok{instance} \DataTypeTok{FunctorMonoid} \DataTypeTok{Compose} \KeywordTok{where}
    \KeywordTok{type} \DataTypeTok{Unit} \DataTypeTok{Compose} \OtherTok{=} \DataTypeTok{Identity}
    \CommentTok{{-}{-} Law 3 (Identity): Compose Identity f  ≅ f}
    \CommentTok{{-}{-} Law 4 (Assoc):    Compose f (Compose g h) ≅ Compose (Compose f g) h}
\end{Highlighting}
\end{Shaded}

\begin{itemize}
\tightlist
\item
  \textbf{Combinator (\texttt{Bin})}: Compose
  (\texttt{Compose\ f\ g\ a\ =\ Compose\ (f\ (g\ a))})
\item
  \textbf{Atom (\texttt{Atom})}: Identity
  (\texttt{Identity\ a\ =\ Identity\ a})
\end{itemize}

Just as \(0\) is the atom for Addition, and \(1\) for Multiplication,
the \texttt{Identity} Functor is the perfect mathematical atom for
Nesting because nesting a functor inside \texttt{Identity} functionally
does absolutely nothing: it preserves the exact original structure!

\subsection{Chapter 2: From Monoids to N-ary
Glues}\label{chapter-2-from-monoids-to-n-ary-glues}

Why is defining the 0-ary Atom and the 2-ary Bin so profound?

Because if you have that \texttt{FunctorMonoid} record, you have
mathematically perfectly defined an \textbf{N-ary Glue}. You can write
exactly one universal folding function that magically turns any valid
Monoid into an infinitely scaling N-ary combinator!

\begin{Shaded}
\begin{Highlighting}[]
\CommentTok{{-}{-} The Universal N{-}ary Factory}
\KeywordTok{type} \KeywordTok{family} \DataTypeTok{FoldGlue}\NormalTok{ (}\OtherTok{m ::}\NormalTok{ [}\DataTypeTok{Type} \OtherTok{{-}\textgreater{}} \DataTypeTok{Type}\NormalTok{] }\OtherTok{{-}\textgreater{}} \OperatorTok{*} \OtherTok{{-}\textgreater{}} \OperatorTok{*}\NormalTok{) (}\OtherTok{fs ::}\NormalTok{ [}\DataTypeTok{Type} \OtherTok{{-}\textgreater{}} \DataTypeTok{Type}\NormalTok{])}\OtherTok{ ::} \OperatorTok{*} \OtherTok{{-}\textgreater{}} \OperatorTok{*} \KeywordTok{where}
    \DataTypeTok{FoldGlue}\NormalTok{ m \textquotesingle{}[]       }\OtherTok{=} \DataTypeTok{Atom}\NormalTok{ m}
    \DataTypeTok{FoldGlue}\NormalTok{ m (f \textquotesingle{}}\OperatorTok{:}\NormalTok{ fs) }\OtherTok{=} \DataTypeTok{Bin}\NormalTok{ m f (}\DataTypeTok{FoldGlue}\NormalTok{ m fs)}
\end{Highlighting}
\end{Shaded}

Thanks to parametricity, folding the \textbf{Sum Record} dynamically
generates the profound construct \texttt{UnionF}:

\begin{Shaded}
\begin{Highlighting}[]
\CommentTok{{-}{-} The Output of folding the Sum Record}
\KeywordTok{data} \DataTypeTok{UnionF}\NormalTok{ (}\OtherTok{fs ::}\NormalTok{ [}\DataTypeTok{Type} \OtherTok{{-}\textgreater{}} \DataTypeTok{Type}\NormalTok{]) a }\KeywordTok{where}
    \DataTypeTok{ThisF}\OtherTok{ ::}\NormalTok{ f a }\OtherTok{{-}\textgreater{}} \DataTypeTok{UnionF}\NormalTok{ (f \textquotesingle{}}\OperatorTok{:}\NormalTok{ fs) a}
    \DataTypeTok{ThatF}\OtherTok{ ::} \DataTypeTok{UnionF}\NormalTok{ fs a }\OtherTok{{-}\textgreater{}} \DataTypeTok{UnionF}\NormalTok{ (f \textquotesingle{}}\OperatorTok{:}\NormalTok{ fs) a}
\end{Highlighting}
\end{Shaded}

And folding the \textbf{Product Record} dynamically generates
\texttt{HListF}:

\begin{Shaded}
\begin{Highlighting}[]
\CommentTok{{-}{-} The Output of folding the Product Record}
\KeywordTok{data} \DataTypeTok{HListF}\NormalTok{ (}\OtherTok{fs ::}\NormalTok{ [}\DataTypeTok{Type} \OtherTok{{-}\textgreater{}} \DataTypeTok{Type}\NormalTok{]) a }\KeywordTok{where}
    \DataTypeTok{HNilF}\OtherTok{  ::} \DataTypeTok{HListF}\NormalTok{ \textquotesingle{}[] a                     }
    \DataTypeTok{HConsF}\OtherTok{ ::}\NormalTok{ f a }\OtherTok{{-}\textgreater{}} \DataTypeTok{HListF}\NormalTok{ fs a }\OtherTok{{-}\textgreater{}} \DataTypeTok{HListF}\NormalTok{ (f \textquotesingle{}}\OperatorTok{:}\NormalTok{ fs) a  }
\end{Highlighting}
\end{Shaded}

We do not need to explicitly declare N-ary glues. They are merely the
syntactic, inductive evaluation of a Functor Monoid folded over a
type-level list!

\subsubsection{Section 2.2: The Ultimate Generator: System
F}\label{section-2.2-the-ultimate-generator-system-f}

As a final profound twist: if you introduce the minimal function arrow
\texttt{(-\textgreater{})} and pair it with Polymorphism
(\texttt{forall}), it completely cannibalizes the rest of the universe.

In Type Theory (System F), using \textbf{Church Encodings}, you can
generate the ENTIRE universe of functors purely out of the Exponential
\texttt{(-\textgreater{})} glue---rendering the primitive Sum and
Product monoids entirely unnecessary!

\begin{itemize}
\tightlist
\item
  \textbf{Product Functor}:
  \texttt{type\ Prod\ f\ g\ a\ =\ forall\ r.\ (f\ a\ -\textgreater{}\ g\ a\ -\textgreater{}\ r)\ -\textgreater{}\ r}
\item
  \textbf{Sum Functor}:
  \texttt{type\ Sum\ f\ g\ a\ =\ forall\ r.\ (f\ a\ -\textgreater{}\ r)\ -\textgreater{}\ (g\ a\ -\textgreater{}\ r)\ -\textgreater{}\ r}
\item
  \textbf{Zero Functor}: \texttt{type\ Zero\ a\ =\ forall\ r.\ r}
\item
  \textbf{Proxy Functor}:
  \texttt{type\ Proxy\ a\ =\ forall\ r.\ r\ -\textgreater{}\ r}
\item
  \textbf{List Functor (Recursive)}:
  \texttt{type\ List\ f\ a\ =\ forall\ r.\ r\ -\textgreater{}\ (f\ a\ -\textgreater{}\ r\ -\textgreater{}\ r)\ -\textgreater{}\ r}
\end{itemize}

If you have the Polymorphic Arrow, its mathematical closure contains
absolutely all possible Functors.

\subsection{Chapter 3: The Value-Level
Symmetry}\label{chapter-3-the-value-level-symmetry}

We have defined \texttt{FunctorMonoid} as a monoid operating at the type
level to combine structural shapes. Does this concept exist at the value
level? Yes! It is the exact symmetry that defines two of Haskell's most
famous typeclasses:

\subsubsection{\texorpdfstring{Section 3.1: The Value-Level Product:
\texttt{Applicative}}{Section 3.1: The Value-Level Product: Applicative}}\label{section-3.1-the-value-level-product-applicative}

If the Product Record is the compile-time combination of functor shapes,
\texttt{Applicative} is the mathematical equivalent applied to runtime
values. * \textbf{0-ary Identity (\texttt{pure})}: Injects EXACTLY ONE
value out of nothing (\(1\)). * \textbf{2-ary Combinator
(\texttt{liftA2\ (,)})}: Takes two values from identical functor shapes
and combines them into a Product holding both.

\subsubsection{\texorpdfstring{Section 3.2: The Value-Level Sum:
\texttt{Alternative}}{Section 3.2: The Value-Level Sum: Alternative}}\label{section-3.2-the-value-level-sum-alternative}

If the Sum Record is the compile-time choice of functor shapes,
\texttt{Alternative} is the mathematical equivalent applied to runtime
values. * \textbf{0-ary Identity (\texttt{empty})}: Represents EXACTLY
ZERO choices (\(0\)). * \textbf{2-ary Combinator
(\texttt{\textless{}\textbar{}\textgreater{}})}: Takes two values from
identical functor shapes and provides a choice (or merging) into one.

The compile-time N-ary glues build new varying data structures
(\texttt{UnionF} and \texttt{HListF}); the runtime \texttt{Applicative}
and \texttt{Alternative} merge values within identical existing
structures. Yet, they are governed by the fundamentally identical
mathematical laws of Functor Monoids!

\subsection{Chapter 4: The Formal
Lexicon}\label{chapter-4-the-formal-lexicon}

There isn't exactly \emph{one} single buzzword that covers the entire
N-ary signature
\texttt{{[}Type\ -\textgreater{}\ Type{]}\ -\textgreater{}\ (*\ -\textgreater{}\ *)},
as the name changes depending on the domain:

\begin{enumerate}
\def\labelenumi{\arabic{enumi}.}
\tightlist
\item
  \textbf{Category Theory}: In the category of endofunctors,
  \texttt{UnionF} is formally the \textbf{N-ary Coproduct of
  Endofunctors}, and \texttt{HListF} is the \textbf{N-ary Product of
  Endofunctors}. Our dummy zero and one glues are the \emph{Initial and
  Terminal Objects} of the Functor Category. The 0-ary and 2-ary pairing
  we defined is formally a \textbf{Monoidal Category over Functors}.
\item
  \textbf{Advanced Haskell (\texttt{generics-sop})}: In libraries like
  Generics-SOP, these N-ary glues represent the true foundational
  bedrock (\texttt{NS} and \texttt{NP}). They are generally categorized
  as \textbf{Functor Combinators} or \textbf{Higher-Order Functors}.
\item
  \textbf{Type Theory (``Polynomials'')}: When building data structures
  entirely out of composed Sums and Products of Functors, mathematicians
  call the resulting space \textbf{Polynomial Endofunctors}. The
  structure perfectly mimics high-school algebra
  (\(F(X) = 1 + X \times F(X)\)).
\item
  \textbf{``Data Types à la Carte''}: In famous research regarding
  composing programming languages dynamically from smaller ADT islands,
  the 0-ary and 2-ary records are referred to as \textbf{Functor
  Coproducts} and Algebraic signatures.
\end{enumerate}

\section{Part 6: The Deep Math}\label{part-6-the-deep-math}

\subsection{Chapter 4: Deep Dive into
Bifunctors}\label{chapter-4-deep-dive-into-bifunctors}

\subsubsection{Section 4.1: The True Nature of
Bifunctors}\label{section-4.1-the-true-nature-of-bifunctors}

In Category Theory, a \textbf{Bifunctor} is not actually a special new
structure; it is quite literally just a standard Functor whose domain
happens to be a \textbf{Product Category}.

If you have two categories, \(\mathcal{C}\) and \(\mathcal{D}\), you can
create a Product Category \(\mathcal{C} \times \mathcal{D}\). The
objects in this category are pairs of objects \((c, d)\), and the
morphisms are pairs of morphisms \((f, g)\). A Bifunctor is simply a
normal Functor \(F\) that maps from that Product Category into a third
category \(\mathcal{E}\):

\[F: \mathcal{C} \times \mathcal{D} \to \mathcal{E}\]

In Haskell, everything happens in the single category \texttt{Hask}. So,
a Haskell \texttt{Bifunctor} is just a standard functor mapping from the
product category to the base category:

\[F: \mathbf{Hask} \times \mathbf{Hask} \to \mathbf{Hask}\]

Because it's just a normal Functor from a Product Category, the mapping
operation (\texttt{bimap}) takes a pair of morphisms (which in Haskell
means two functions: \texttt{(a\ -\textgreater{}\ c)} and
\texttt{(b\ -\textgreater{}\ d)}) and applies them to the pair of
objects inside the structure. \#\#\# Section 4.2: The Laws of Bifunctors

When you implement an \texttt{instance\ Bifunctor} in Haskell, you must
satisfy laws analogous to the standard Functor laws, just extended over
two dimensions. Since a Bifunctor is just a functor from a Product
Category, the morphisms we are mapping are pairs of functions.

The two laws are:

\begin{enumerate}
\def\labelenumi{\arabic{enumi}.}
\item
  \textbf{Identity Law}:

\begin{Shaded}
\begin{Highlighting}[]
\NormalTok{bimap }\FunctionTok{id} \FunctionTok{id} \OperatorTok{==} \FunctionTok{id}
\end{Highlighting}
\end{Shaded}

  \emph{Meaning}: If you apply the identity function to both the left
  and right sides simultaneously, the structure and its contents must
  remain completely unchanged.
\item
  \textbf{Composition Law}:

\begin{Shaded}
\begin{Highlighting}[]
\NormalTok{bimap (f }\OperatorTok{.}\NormalTok{ g) (h }\OperatorTok{.}\NormalTok{ i) }\OperatorTok{==}\NormalTok{ bimap f h }\OperatorTok{.}\NormalTok{ bimap g i}
\end{Highlighting}
\end{Shaded}

  \emph{Meaning}: Composing two functions and then mapping them over a
  Bifunctor is identical to mapping the first pair of functions, and
  then mapping the second pair of functions over the result.
\end{enumerate}

\textbf{The \texttt{first} and \texttt{second} Equivalences} The
\texttt{Data.Bifunctor} typeclass in Haskell also provides the helper
functions \texttt{first} and \texttt{second} to map over only one side
of the Bifunctor. The definition of a Bifunctor inextricably links
\texttt{bimap}, \texttt{first}, and \texttt{second} through these
properties: * \texttt{bimap\ f\ g\ ==\ first\ f\ .\ second\ g} *
\texttt{first\ f\ ==\ bimap\ f\ id} *
\texttt{second\ g\ ==\ bimap\ id\ g}

This reinforces the concept described in
\hyperref[section-13-the-algebra-of-functors-bifunctors]{Section 1.4}:
if you fix the identity function to one side of a Bifunctor, it
mathematically collapses into a standard Endofunctor.

\subsection{Chapter 5: Monoidal
Categories}\label{chapter-5-monoidal-categories}

\subsubsection{Section 5.1: The Pentagon and Triangle
Laws}\label{section-5.1-the-pentagon-and-triangle-laws}

A \textbf{Monoidal Category} is a higher-level structure that uses a
specific Bifunctor as a ``tensor product''.

It is defined by: 1. A base category (like \texttt{Hask}). 2. A specific
\textbf{Bifunctor} acting as the tensor product (like \texttt{(,)} or
\texttt{Either}). 3. A unit object (like \texttt{()} for products, or
\texttt{Void} for sums). 4. Associativity and Unit natural isomorphisms.
5. \textbf{The Coherence Conditions}: This is where the \textbf{pentagon
identity} (ensuring associativity associates consistently) and the
\textbf{triangle identity} (ensuring the unit behaves consistently) come
into play.

While \texttt{Either} and \texttt{(,)} are Bifunctors, they are
\emph{special} Bifunctors because they act as the tensor products that
turn \texttt{Hask} into a Monoidal Category. Other Bifunctors (like
\texttt{BiProxy} or \texttt{ConstContext\ r}) are perfectly valid
Bifunctors without forming a monoidal category with strict
pentagon/triangle laws.

\emph{(This section is a placeholder for a future deep-dive into the
formal definitions of tensor products, and how they interact
structurally via the pentagon and triangle identities).}

\begin{center}\rule{0.5\linewidth}{0.5pt}\end{center}

\subsubsection{Parametricity in Category
Theory}\label{parametricity-in-category-theory}

Yes, there is a very deep and well-established categorical notion of
Haskell's parametricity. In fact, category theory provides the exact
mathematical language needed to formalize what Philip Wadler famously
called ``Theorems for Free!''

Depending on how deep you want to go, parametricity can be understood
categorically in a few distinct layers, ranging from natural
transformations up to relational fibrations. Here is a breakdown of how
category theory models parametric polymorphism.

\paragraph{1. Natural Transformations (The Basic
View)}\label{natural-transformations-the-basic-view}

At the simplest level, if a type variable only appears in covariant
positions (like the output of a function or inside a standard data
structure), parametricity corresponds exactly to natural
transformations.

Suppose you have a polymorphic function in Haskell:

\begin{Shaded}
\begin{Highlighting}[]
\OtherTok{f ::} \KeywordTok{forall}\NormalTok{ a}\OperatorTok{.}\NormalTok{ [a] }\OtherTok{{-}\textgreater{}} \DataTypeTok{Maybe}\NormalTok{ a}
\end{Highlighting}
\end{Shaded}

Categorically, \texttt{{[}{]}} (List) and \texttt{Maybe} are functors
from the category of Haskell types (\(\mathbf{Hask}\)) to itself. The
polymorphic function \texttt{f} is a natural transformation
\(\eta : \text{List} \to \text{Maybe}\).

The defining property of a natural transformation is that for any
function \(h : A \to B\), the following square commutes: \$ \eta\_B
\circ \text{List}(h) = \text{Maybe}(h) \circ \eta\_A \$

In Haskell syntax, this is exactly the ``free theorem'' for \texttt{f}:

\begin{Shaded}
\begin{Highlighting}[]
\NormalTok{f }\OperatorTok{.} \FunctionTok{fmap}\NormalTok{ h }\OperatorTok{==} \FunctionTok{fmap}\NormalTok{ h }\OperatorTok{.}\NormalTok{ f}
\end{Highlighting}
\end{Shaded}

Because the function is parametrically polymorphic, it must be a natural
transformation, meaning it cannot inspect the elements it is moving
around.

\paragraph{2. Ends and Coends (Universal and Existential
Types)}\label{ends-and-coends-universal-and-existential-types}

When type variables appear in both contravariant (input) and covariant
(output) positions, we need a stronger categorical concept to represent
the \texttt{forall} keyword. This is where \textbf{Ends} come in.

Consider the identity function type:

\begin{Shaded}
\begin{Highlighting}[]
\FunctionTok{id}\OtherTok{ ::} \KeywordTok{forall}\NormalTok{ a}\OperatorTok{.}\NormalTok{ a }\OtherTok{{-}\textgreater{}}\NormalTok{ a}
\end{Highlighting}
\end{Shaded}

Here, \texttt{a} is an argument to the function constructor
\texttt{(-\textgreater{})}, which is a functor that is contravariant in
its first argument and covariant in its second. So,
\texttt{(-\textgreater{})} is a functor
\(H : \mathbf{Hask}^{op} \times \mathbf{Hask} \to \mathbf{Hask}\).

The \texttt{forall} quantifier is interpreted as a categorical end over
this mixed-variance functor: \[ \int_{A \in \mathbf{Hask}} H(A, A) \] An
end conceptually ``takes the intersection'' over all objects \(A\),
giving you the family of morphisms that act uniformly across all types.
(Conversely, existential types like \texttt{exists\ a.\ ...} are modeled
by coends, \(\int^{A} H(A, A)\)).

\paragraph{3. Dinatural Transformations (Mixed
Variance)}\label{dinatural-transformations-mixed-variance}

The ``elements'' (or points) of the end \(\int_A H(A, A)\) are called
\textbf{dinatural transformations}.

Standard natural transformations only work between functors of the same
variance. Because \texttt{id} maps from a contravariant position to a
covariant position, its uniformity is expressed as a dinatural
transformation. The dinaturality hexagon (the commuting diagram for
dinatural transformations) gives you the exact free theorem for
functions with mixed-variance type signatures.

\paragraph{4. Reflexive Graphs and Relational Fibrations (The Deep
View)}\label{reflexive-graphs-and-relational-fibrations-the-deep-view}

While naturality and ends describe how polymorphic functions behave
structurally, John C. Reynolds' original abstraction theorem
(\textbf{Relational Parametricity}) states that polymorphic functions
must preserve relations, not just functions.

To model this categorically, we have to move beyond just looking at the
category of types and functions. We use a structure often called a
\textbf{Reflexive Graph Category} or a \textbf{Relational Fibration}.

\begin{enumerate}
\def\labelenumi{\arabic{enumi}.}
\tightlist
\item
  We construct a base category \(\mathbb{B}\) where objects are types
  and morphisms are functions.
\item
  We construct a total category \(\mathbb{E}\) where objects are
  relations between types, and morphisms are pairs of functions that
  preserve those relations.
\item
  There is a functor \(p : \mathbb{E} \to \mathbb{B} \times \mathbb{B}\)
  that projects a relation down to the two types it relates.
\end{enumerate}

In this setting, a type operator (like List) isn't just a functor; it
must be a functor that lifts to relations (e.g., if you have a relation
\(R\) between \(A\) and \(B\), you automatically get a relation
\(\text{List}(R)\) between \(\text{List}(A)\) and \(\text{List}(B)\)). A
parametrically polymorphic function is then an object in this higher
category that intrinsically preserves all relations, fulfilling
Reynolds' exact definition.

\begin{center}\rule{0.5\linewidth}{0.5pt}\end{center}

\subsubsection{Summary and Type Bundle
Taxonomy}\label{summary-and-type-bundle-taxonomy}

At this level, Functors are entirely about \textbf{Shape and
Preservation}. Whether we are dealing with an empty box
(\texttt{Proxy}), a wrapper (\texttt{Identity}), or an infinite chain
(\texttt{List}), \texttt{fmap} ensures that the structure of the data
remains physically identical while the values inside are transformed.

Before moving to Applicatives, remember the three tools Haskell gives us
to bundle these shapes:

\begin{enumerate}
\def\labelenumi{\arabic{enumi}.}
\tightlist
\item
  \textbf{\texttt{type} (Alias)}: No new type created, zero overhead.
  Use for readability.
\item
  \textbf{\texttt{newtype} (Strict Wrapper)}: Distinct type, zero
  overhead. Use for type safety (e.g., \texttt{UserId}).
\item
  \textbf{\texttt{data} (Full ADT)}: Flexible, supports multiple
  constructors. Use for complex shapes.
\end{enumerate}

\begin{center}\rule{0.5\linewidth}{0.5pt}\end{center}

\subsection{Annex A: Proof of 2-Inhabitant
Associativity}\label{annex-a-proof-of-2-inhabitant-associativity}

\emph{Proof that all 2-inhabitant logical operations possessing a valid
two-sided identity element are automatically associative.}

Assume we have a 2-element type \texttt{\{E,\ X\}}. We declare that
\texttt{E} is the identity element. Because \texttt{E} is the identity,
the rules \texttt{E\ ⋄\ a\ =\ a} and \texttt{a\ ⋄\ E\ =\ a} instantly
lock in 3 of the 4 slots in our logical multiplication table (Cayley
Table):

\begin{longtable}[]{@{}ccc@{}}
\toprule\noalign{}
\texttt{x\ \textbackslash{}\ y} & E & X \\
\midrule\noalign{}
\endhead
\bottomrule\noalign{}
\endlastfoot
\textbf{E} & \texttt{E} & \texttt{X} \\
\textbf{X} & \texttt{X} & \textbf{???} \\
\end{longtable}

Because this table represents a closed binary operation on a
2-inhabitant type, there are only two possible values we can pick for
the single empty box (\texttt{X\ ⋄\ X}): 1. \textbf{Case 1}:
\texttt{X\ ⋄\ X\ =\ E} (This forms the XOR or Equivalence gate) 2.
\textbf{Case 2}: \texttt{X\ ⋄\ X\ =\ X} (This forms the OR or AND gate)

Now, we must test the Law of Associativity:
\texttt{(a\ ⋄\ b)\ ⋄\ c\ ==\ a\ ⋄\ (b\ ⋄\ c)}.

\begin{itemize}
\tightlist
\item
  \textbf{Trivial Cases}: If any of \texttt{a}, \texttt{b}, or
  \texttt{c} is the identity \texttt{E}, the associativity law trivially
  holds without evaluation. (For example, if \texttt{a\ =\ E}, then
  \texttt{(E\ ⋄\ b)\ ⋄\ c\ =\ b\ ⋄\ c}, and
  \texttt{E\ ⋄\ (b\ ⋄\ c)\ =\ b\ ⋄\ c}).
\item
  \textbf{The Single Non-Trivial Case}: The only case where none of the
  operands are the identity is when \texttt{a}, \texttt{b}, and
  \texttt{c} are all \texttt{X}. Thus, the entire proof hinges on
  whether \texttt{(X\ ⋄\ X)\ ⋄\ X\ ==\ X\ ⋄\ (X\ ⋄\ X)}.
\end{itemize}

Let us evaluate this final equation for both possible remaining tables:

\textbf{Case 1 (where \texttt{X\ ⋄\ X\ =\ E})}: * Left side:
\texttt{(X\ ⋄\ X)\ ⋄\ X\ =\ E\ ⋄\ X\ =\ X} * Right side:
\texttt{X\ ⋄\ (X\ ⋄\ X)\ =\ X\ ⋄\ E\ =\ X} * Since \texttt{X\ =\ X},
Associativity unconditionally holds.

\textbf{Case 2 (where \texttt{X\ ⋄\ X\ =\ X})}: * Left side:
\texttt{(X\ ⋄\ X)\ ⋄\ X\ =\ X\ ⋄\ X\ =\ X} * Right side:
\texttt{X\ ⋄\ (X\ ⋄\ X)\ =\ X\ ⋄\ X\ =\ X} * Since \texttt{X\ =\ X},
Associativity unconditionally holds.

\textbf{Q.E.D.} Once you successfully lock in an identity element on a
2-inhabitant type, there is simply no remaining mathematical room in the
\(2 \times 2\) matrix for associativity to break!

\begin{center}\rule{0.5\linewidth}{0.5pt}\end{center}

\subsection{Annex: Proofs and
Derivations}\label{annex-proofs-and-derivations}

\subsubsection{Proof of Monad
Equivalence}\label{proof-of-monad-equivalence}

If \texttt{bind}, \texttt{join}, and \texttt{kleisli} are equivalent, we
should be able to derive them algebraically from one another (assuming
we have \texttt{fmap} and \texttt{pure}).

\textbf{Deriving \texttt{join} from \texttt{bind}}:

\begin{Shaded}
\begin{Highlighting}[]
\NormalTok{join mm }\OtherTok{=}\NormalTok{ mm }\OperatorTok{\textgreater{}\textgreater{}=} \FunctionTok{id}  
\CommentTok{{-}{-} Where mm :: m (m a) and id :: m a {-}\textgreater{} m a}
\end{Highlighting}
\end{Shaded}

\textbf{Deriving \texttt{bind} from \texttt{join} and \texttt{fmap}}:

\begin{Shaded}
\begin{Highlighting}[]
\NormalTok{m }\OperatorTok{\textgreater{}\textgreater{}=}\NormalTok{ f }\OtherTok{=}\NormalTok{ join (}\FunctionTok{fmap}\NormalTok{ f m)}
\CommentTok{{-}{-} fmap f m produces m (m b)}
\CommentTok{{-}{-} join then flattens it to m b}
\end{Highlighting}
\end{Shaded}

\textbf{Deriving \texttt{kleisli} from \texttt{bind}}:

\begin{Shaded}
\begin{Highlighting}[]
\NormalTok{(f }\OperatorTok{\textgreater{}=\textgreater{}}\NormalTok{ g) x }\OtherTok{=}\NormalTok{ f x }\OperatorTok{\textgreater{}\textgreater{}=}\NormalTok{ g}
\end{Highlighting}
\end{Shaded}

\subsubsection{Proof of Unique Functor
Identity}\label{proof-of-unique-functor-identity}

Given \texttt{data\ Proxy\ a\ =\ Proxy}, how do we formally prove the
only valid function of type
\texttt{(a\ -\textgreater{}\ a)\ -\textgreater{}\ Proxy\ a\ -\textgreater{}\ Proxy\ a}
preserving structure is the identity?

\begin{enumerate}
\def\labelenumi{\arabic{enumi}.}
\tightlist
\item
  Let \texttt{g\ ::\ Proxy\ a\ -\textgreater{}\ Proxy\ a} be a total,
  terminating function.
\item
  The only inhabited value of \texttt{Proxy\ a} at the term level is
  \texttt{Proxy}.
\item
  Therefore, \texttt{g\ Proxy\ =\ Proxy}.
\item
  By definition, the \texttt{id} function is \texttt{id\ x\ =\ x}.
\item
  Thus, \texttt{id\ Proxy\ =\ Proxy}.
\item
  Since \texttt{g\ Proxy\ =\ id\ Proxy} for the sole value of the type,
  \texttt{g\ =\ id}. Because \texttt{g} is the only total mapping,
  \texttt{fmap\ id\ =\ id} trivially holds, and no other lawful
  interpretation exists.
\end{enumerate}

\subsubsection{Proof of Identity Implies
Composition}\label{proof-of-identity-implies-composition}

In Haskell, if \texttt{fmap\ id\ =\ id} (Identity Law) holds for a
parametrically polymorphic \texttt{fmap}, then
\texttt{fmap\ (f\ .\ g)\ =\ fmap\ f\ .\ fmap\ g} (Composition Law) is
automatically satisfied. This is a direct consequence of the
\textbf{Naturality} of \texttt{fmap}.

\begin{enumerate}
\def\labelenumi{\arabic{enumi}.}
\tightlist
\item
  \textbf{The Signature}:
  \texttt{fmap\ ::\ forall\ a\ b.\ (a\ -\textgreater{}\ b)\ -\textgreater{}\ F\ a\ -\textgreater{}\ F\ b}.
\item
  \textbf{Naturality Condition}: For any natural transformation
  \(\eta : F \to G\), and any function \(k : a \to b\), the condition
  \(\eta_b \circ F(k) = G(k) \circ \eta_a\) must hold.
\item
  \textbf{Treating \texttt{fmap} as a transformation}: We can view
  \(fmap(f)\) as a transformation from the functor \texttt{F} to itself.
\item
  \textbf{Free Theorem}: The ``Free Theorem'' for the type of
  \texttt{fmap} (as derived in Wadler's paper) states:
  \texttt{fmap\ f\ .\ fmap\ g\ =\ fmap\ (f\ .\ g)} This equality holds
  because the type of \texttt{fmap} is so restrictive that it cannot
  differentiate between ``applying a composition'' and ``composing two
  applications'' without knowing the internal structure of the
  types---which parametricity forbids.
\end{enumerate}

\subsubsection{Proof of Identity
Uniqueness}\label{proof-of-identity-uniqueness}

If a binary operation \(B\) has a left identity \(I_L\) and a right
identity \(I_R\), they must be structurally isomorphic
(\(I_L \cong I_R\)).

\begin{enumerate}
\def\labelenumi{\arabic{enumi}.}
\tightlist
\item
  By definition of a Left Identity, for \emph{any} type \(A\):
  \(B(I_L, A) \cong A\).
\item
  Let's choose \(A = I_R\). Therefore: \(B(I_L, I_R) \cong I_R\).
\item
  By definition of a Right Identity, for \emph{any} type \(A\):
  \(B(A, I_R) \cong A\).
\item
  Let's choose \(A = I_L\). Therefore: \(B(I_L, I_R) \cong I_L\).
\item
  Since both \(I_L\) and \(I_R\) are mathematically isomorphic to the
  exact same structural element \(B(I_L, I_R)\), then logically
  \(I_L \cong I_R\).
\end{enumerate}

Thus, if both exist, they are structurally identical.

\subsubsection{Strict Equality vs.~Structural
Isomorphism}\label{strict-equality-vs.-structural-isomorphism}

It is critical mathematically and computationally to distinguish between
typeclass ``laws'' and structural ``isomorphisms'' or Unitors.

\textbf{Typeclass Laws (Behavioral strict equality)}: These govern how
typeclass methods (like \texttt{fmap}, \texttt{bimap},
\texttt{\textgreater{}\textgreater{}=}) must behave computationally.
When we write a law like \texttt{fmap\ id\ ==\ id}, we mean
\textbf{strict equality}. The two sides must evaluate to the exact same
value of the exact same type. If you map the identity over
\texttt{{[}1,\ 2{]}}, you must get the exact \texttt{{[}1,\ 2{]}} back,
not a copy or an equivalent wrapper.

\textbf{Structural Isomorphisms (Type-level mapping)}: These govern the
``shape'' of the types themselves. When we say \(B(I, A) \cong A\)
(e.g., \texttt{(Either\ Void\ A)\ ≅\ A}), we are describing
\textbf{structural isomorphism}. The compiler knows that
\texttt{Either\ Void\ Bool} and \texttt{Bool} are two entirely different
types (\texttt{Left\ True} vs \texttt{True}). However, because
\texttt{Void} contains no information, we can write a perfect, lossless
two-way mapping between the two structures. These mappings are the exact
``Left/Right Unitors''. They are not equalities; they are natural
transformations between non-equal types.

\begin{quote}
For references, papers, and further reading on these advanced concepts,
refer to \href{09_bibliography.md}{Part 9: Bibliography}.
\end{quote}

\section{Part 9: Bibliography}\label{part-9-bibliography}

This section centralizes all the foundational papers, influential
articles, and recommended reading for the concepts discussed throughout
this series.

\subsubsection{Type Theory \& Algebraic Data
Types}\label{type-theory-algebraic-data-types}

\begin{enumerate}
\def\labelenumi{\arabic{enumi}.}
\tightlist
\item
  \textbf{Howard, W. A. (1980).} \emph{The formulae-as-types notion of
  construction}. In To H. B. Curry: Essays on Combinatory Logic, Lambda
  Calculus and Formalism (pp.~479-490). Academic Press. \emph{(The
  seminal work establishing the Curry-Howard correspondence connecting
  uninhabited types to logical falsity).}
\item
  \textbf{Maguire, S. (2018).}
  \emph{\href{https://thinkingwithtypes.com/}{Thinking with Types}}.
  \emph{(Features early chapters specifically focusing on the
  ``Cardinality'' and algebra of types, using structural counting to
  prove which functions can mathematically exist).}
\item
  \textbf{McBride, C. (2001).} \emph{The derivative of a regular type is
  its type of one-hole contexts}. \emph{(A mind-bending proof that
  taking the calculus derivative of a data type's polynomial
  structurally generates its exact Zipper).}
\item
  \textbf{Swierstra, W. (2008).} \emph{Data types à la carte}. \emph{(A
  seminal paper proving how to use the algebraic Sum operator over
  Functors to modularly compose distinct data types and interpreters).}
\item
  \textbf{Taylor, C. (2013).}
  \emph{\href{https://chris-taylor.github.io/blog/2013/02/10/the-algebra-of-algebraic-data-types/}{The
  Algebra of Algebraic Data Types}}. \emph{(A famous blog series
  expanding the structural analogy by directly using high-school algebra
  to calculate isomorphic data types).}
\end{enumerate}

\subsubsection{Category Theory \& Deep
Math}\label{category-theory-deep-math}

\begin{enumerate}
\def\labelenumi{\arabic{enumi}.}
\setcounter{enumi}{5}
\tightlist
\item
  \textbf{Bird, R. \& de Moor, O. (1997).} \emph{Algebra of
  Programming}. \emph{(A foundational text exploring how algebras and
  functor subcategories are derived systematically from building blocks
  like Bifunctors).}
\item
  \textbf{Milewski, B. (2014).}
  \emph{\href{https://bartoszmilewski.com/2014/10/28/category-theory-for-programmers-the-preface/}{Category
  Theory for Programmers}}. \emph{(A highly acclaimed resource that
  connects structural logic to Haskell, including the definitions of the
  Initial Object (0) and Terminal Object (1) purely structurally before
  ever introducing Functors and Monads).}
\item
  \textbf{Wadler, P. (1989).} \emph{Theorems for free!}. \emph{(The
  definitive paper explaining parametricity in Category Theory and why
  we get free laws for our functions and data types).}
\end{enumerate}

\subsubsection{Typeclasses, Functors, and
Monads}\label{typeclasses-functors-and-monads}

\begin{enumerate}
\def\labelenumi{\arabic{enumi}.}
\setcounter{enumi}{8}
\tightlist
\item
  \textbf{Bhargava, A.} \emph{Functors, Applicatives, And Monads In
  Pictures}. \emph{(A highly recommended visual guide for grasping the
  basic intuition behind these three crucial typeclasses).}
\item
  \textbf{Moggi, E. (1991).} \emph{Notions of computation and monads}.
  \emph{(The foundational paper introducing the concept of monads to
  programming languages to model side-effects).}
\item
  \textbf{Yorgey, B. (2009).}
  \emph{\href{https://wiki.haskell.org/Typeclassopedia}{The
  Typeclassopedia}}. The Monad Reader Issue 13. \emph{(The definitive
  guide to mapping out the core Haskell typeclasses, their
  relationships, and their mathematical laws).}
\end{enumerate}

\subsubsection{Practical Haskell \&
Reasoning}\label{practical-haskell-reasoning}

\begin{enumerate}
\def\labelenumi{\arabic{enumi}.}
\setcounter{enumi}{11}
\tightlist
\item
  \textbf{Danielsson, N. A., Hughes, J., Jansson, P., \& Gibbons, J.
  (2006).} \emph{Fast and Loose Reasoning is Morally Correct}. ACM
  SIGPLAN Notices, 41(1), 273-284. \emph{(A formal justification for
  reasoning about Haskell programs while ignoring
  \texttt{\_\textbar{}\_}, widely accepted as standard practice in the
  Haskell community).}
\item
  \textbf{Diehl, S.}
  \emph{\href{https://smunix.github.io/dev.stephendiehl.com/hask/tutorial.pdf}{What
  I Wish I Knew When Learning Haskell}}. \emph{(A comprehensive guide to
  practical Haskell, covering many advanced type-level mechanics
  including \texttt{Void} and phantom types).}
\item
  \textbf{King, A. (2019).}
  \emph{\href{https://lexi-lambda.github.io/blog/2019/11/05/parse-don-t-validate/}{Parse,
  don't validate}}. \emph{(A highly influential post demonstrating how
  to use the type system, including uninhabited types, to prove
  properties and prevent invalid states).}
\item
  \textbf{Leijen, D., \& Meijer, E. (1999).} \emph{Domain Specific
  Embedded Compilers}. ACM SIGPLAN Notices, 35(1), 109-122. \emph{(An
  early and influential paper showcasing the use of Phantom Types in
  Haskell).}
\end{enumerate}

\end{document}
